% The atom — Physics Upper Sixth — Cours signé OnBuch+
% Official MINESEC syllabus. Compiler avec : tectonic PHY11-the-atom.tex
\newcommand{\DOCMATIERE}{Physics}
\newcommand{\DOCNIVEAU}{Upper Sixth}
\newcommand{\DOCMODULE}{Upper Sixth — Physics}
\newcommand{\DOCLECON}{11}
\newcommand{\DOCTITRE}{The atom}
% OnBuch+ — préambule « cours signés » ENRICHI (compilé avec tectonic / XeLaTeX)
% Système visuel « Pop » d'OnBuch+ : fond crème (jamais blanc), bordures encre
% épaisses, ombres « dures » décalées, titres Archivo Black en capitales.
% Version MATHÉMATIQUES : graphes (pgfplots/tikz), boîtes pédagogiques
% (définition, propriété, méthode, attention, le savais-tu, exercice résolu,
% exercice, TP, bilan), page de garde et signature OnBuch+ ancrée en bas.
%
% Macros attendues dans main.tex AVANT \input{preamble} :
%   \DOCMATIERE  \DOCNIVEAU  \DOCMODULE  \DOCLECON (numéro)  \DOCTITRE
\documentclass[11pt]{article}

\usepackage{fontspec}
\usepackage[english]{babel}
\usepackage[a4paper,margin=2cm,top=3.4cm,bottom=2.8cm,headheight=1.4cm,footskip=1.3cm]{geometry}
\usepackage{amsmath,amssymb,amsfonts}
\usepackage{mathtools} % \xmapsto, \coloneqq... souvent utilisés par les modèles
\usepackage{cancel} % simplifications barrées (bilans de forces)
% \abs{z} (module, valeur absolue) et \norm{u} (norme), utilisés par les modèles
\providecommand{\abs}{}\let\abs\relax\DeclarePairedDelimiter{\abs}{\lvert}{\rvert}
\providecommand{\norm}{}\let\norm\relax\DeclarePairedDelimiter{\norm}{\lVert}{\rVert}
\usepackage[table]{xcolor} % [table] : fournit aussi \rowcolors (alternance de lignes)
\usepackage{pagecolor}
\usepackage{tikz}
\usetikzlibrary{arrows.meta,positioning,calc,shapes.geometric,shapes.misc,shapes.arrows,decorations.pathmorphing,decorations.pathreplacing,decorations.markings,patterns,shadows,mindmap,trees,fit,backgrounds,matrix,intersections,angles,quotes}
\usepackage{pgfplots}
\usepackage[version=4]{mhchem} % équations nucléaires \ce{^{235}_{92}U}
\usepackage[european,siunitx]{circuitikz} % schémas électriques
\pgfplotsset{compat=1.18}
\usepgfplotslibrary{fillbetween,groupplots} % aires entre courbes (calcul intégral)
\usepackage{enumitem}
\usepackage{fancyhdr}
% [most] charge toutes les bibliothèques tcolorbox (listings, minted, raster,
% documentation...) et gonfle inutilement la mémoire TeX sur un document long
% (~300 boîtes) : on ne charge que ce qui est réellement utilisé.
\usepackage[skins,breakable]{tcolorbox}
\usepackage{graphicx}
\usepackage{colortbl}
\usepackage{booktabs}
\usepackage{tabularx}
\usepackage{diagbox} % case d'en-tête diagonale des tableaux à double entrée
% Colonnes extensibles centrée / gauche / droite (C, L, R), souvent utilisées par les modèles
\newcolumntype{C}{>{\centering\arraybackslash}X}
\newcolumntype{L}{>{\raggedright\arraybackslash}X}
\newcolumntype{R}{>{\raggedleft\arraybackslash}X}
\usepackage{array}
\usepackage{multicol}
\usepackage{siunitx}
% parse-numbers=false : accepte aussi \SI{2,5 \times 10^{-3}}{...}
\sisetup{locale=UK,output-decimal-marker={.},group-minimum-digits=4,parse-numbers=false}
\DeclareSIUnit\ohm{\ensuremath{\symup{\Omega}}} % U+2126 absent de la police
\DeclareSIUnit\division{div} % réglages d'oscilloscope (ms/div, V/div)
\DeclareSIUnit\tr{tr} % tours (vitesse de rotation, ex. \SI{1500}{\tr\per\minute})
\usepackage{qrcode}
% \degree : commande fréquemment utilisée par les modèles pour un angle ou une
% température en dehors de \SI{}{} (non fournie par les paquets chargés).
\providecommand{\degree}{^{\circ}}
% \qed (fin de démonstration), utilisé par les modèles sans amsthm
\providecommand{\qed}{\hfill$\square$}
% \barre : double barre des valeurs interdites dans un tableau de variations (tkz-tab)
\providecommand{\barre}{\ensuremath{\|}}
% environnement proof (amsthm non chargé) utilisé par les modèles
\ifdefined\proof\else\newenvironment{proof}[1][Proof]{\par\noindent\textit{#1.}\ }{\qed\par}\fi
\usepackage{lastpage}
\usepackage{hyperref}
\hypersetup{hidelinks}

% ============================================================
% Palette « Pop » OnBuch+ — mêmes hex que la classe P (plus_ui.dart)
% ============================================================
\definecolor{poporange}{HTML}{F59321}
\definecolor{popdark}{HTML}{E07A0C}
\definecolor{popcream}{HTML}{FFF6E8}
\definecolor{popink}{HTML}{1C1714}
\definecolor{poppurple}{HTML}{7A5AE0}
\definecolor{popgreen}{HTML}{2E9E6A}
\definecolor{poppink}{HTML}{E0457B}
\definecolor{popblue}{HTML}{2F7DE1}
\definecolor{popgold}{HTML}{C98A12}
\definecolor{popmuted}{HTML}{8A7F76}
\definecolor{popline}{HTML}{EADFD2}
\definecolor{popgray}{HTML}{8A7F76}
\colorlet{popgrey}{popgray}
% Alias tolérants (le modèle nomme parfois les couleurs différemment)
\colorlet{popwhite}{white}
\colorlet{popblack}{popink}
\colorlet{popred}{poppink}
\colorlet{popyellow}{popgold}
% Teintes claires pour fonds de tableaux / zones
\colorlet{popblueL}{popblue!10!white}
\colorlet{poporangeL}{poporange!14!white}
\colorlet{popgreenL}{popgreen!12!white}
\colorlet{poppurpleL}{poppurple!12!white}
\colorlet{poppinkL}{poppink!12!white}
\colorlet{popgoldL}{popgold!14!white}

\pagecolor{popcream}

% ============================================================
% Typographie
% ============================================================
\defaultfontfeatures{Ligatures=TeX}
\setmainfont{PlusJakartaSans-Medium.ttf}[Path=../../../fonts/,BoldFont=PlusJakartaSans-Bold.ttf]
% Maths en sans-serif (Fira Math) avec chiffres et lettres droites de Plus
% Jakarta, pour que les formules se fondent dans le texte.
\usepackage[math-style=ISO,bold-style=ISO]{unicode-math}
\setmathfont{FiraMath-Regular.otf}[Path=../../../fonts/]
\setmathfont{PlusJakartaSans-Medium.ttf}[Path=../../../fonts/,range={up/num,up/{Latin,latin},\mathup}]
% (icomma retiré : décimales avec point en anglais)
% Caractères mathématiques Unicode absents des polices de texte (Plus Jakarta,
% Archivo Black) : rendus par leur équivalent mathématique, en texte comme en
% formule — sinon ils s'affichent en carré vide (ex. « Arithmétique dans ℤ »).
\usepackage{newunicodechar}
\newunicodechar{ℝ}{\ensuremath{\mathbb{R}}}
\newunicodechar{ℤ}{\ensuremath{\mathbb{Z}}}
\newunicodechar{ℂ}{\ensuremath{\mathbb{C}}}
\newunicodechar{ℕ}{\ensuremath{\mathbb{N}}}
\newunicodechar{ℚ}{\ensuremath{\mathbb{Q}}}
\newunicodechar{𝔻}{\ensuremath{\mathbb{D}}}
\newunicodechar{α}{\ensuremath{\alpha}}
\newunicodechar{β}{\ensuremath{\beta}}
\newunicodechar{γ}{\ensuremath{\gamma}}
\newunicodechar{δ}{\ensuremath{\delta}}
\newunicodechar{ε}{\ensuremath{\varepsilon}}
\newunicodechar{θ}{\ensuremath{\theta}}
\newunicodechar{λ}{\ensuremath{\lambda}}
\newunicodechar{μ}{\ensuremath{\mu}}
\newunicodechar{σ}{\ensuremath{\sigma}}
\newunicodechar{τ}{\ensuremath{\tau}}
\newunicodechar{φ}{\ensuremath{\varphi}}
\newunicodechar{ψ}{\ensuremath{\psi}}
\newunicodechar{ω}{\ensuremath{\omega}}
\newunicodechar{Σ}{\ensuremath{\Sigma}}
% majuscules produites par \MakeUppercase dans les titres (α-aminés -> Α)
\newunicodechar{Α}{\ensuremath{\alpha}}
\newunicodechar{Β}{\ensuremath{\beta}}
\newunicodechar{Γ}{\ensuremath{\Gamma}}
\newunicodechar{Δ}{\ensuremath{\Delta}}
\newunicodechar{Ε}{\ensuremath{\varepsilon}}
\newunicodechar{Θ}{\ensuremath{\Theta}}
\newunicodechar{Λ}{\ensuremath{\Lambda}}
\newunicodechar{Μ}{\ensuremath{\mu}}
\newunicodechar{Τ}{\ensuremath{\tau}}
\newunicodechar{Φ}{\ensuremath{\Phi}}
\newunicodechar{Ψ}{\ensuremath{\Psi}}
\newunicodechar{Ω}{\ensuremath{\Omega}}
\newunicodechar{↦}{\ensuremath{\mapsto}}
\newunicodechar{∈}{\ensuremath{\in}}
\newunicodechar{∉}{\ensuremath{\notin}}
\newunicodechar{∧}{\ensuremath{\wedge}}
\newunicodechar{⇔}{\ensuremath{\Leftrightarrow}}
\newunicodechar{⟺}{\ensuremath{\Longleftrightarrow}}
\newunicodechar{⇒}{\ensuremath{\Rightarrow}}
\newunicodechar{⟹}{\ensuremath{\Longrightarrow}}
\newunicodechar{∘}{\ensuremath{\circ}}
\newunicodechar{∪}{\ensuremath{\cup}}
\newunicodechar{∩}{\ensuremath{\cap}}
\newunicodechar{≡}{\ensuremath{\equiv}}
\newunicodechar{⊂}{\ensuremath{\subset}}
\newunicodechar{⊥}{\ensuremath{\perp}}
\newunicodechar{∃}{\ensuremath{\exists}}
\newunicodechar{∀}{\ensuremath{\forall}}
\newunicodechar{∛}{\ensuremath{\sqrt[3]{\,}}}
\newunicodechar{∜}{\ensuremath{\sqrt[4]{\,}}}
\newunicodechar{⃗}{\ensuremath{{}^{\to}}}
\newunicodechar{ⁿ}{\ifmmode^{n}\else\textsuperscript{\ensuremath{n}}\fi}
\newunicodechar{ˣ}{\ifmmode^{x}\else\textsuperscript{\ensuremath{x}}\fi}
\newunicodechar{ᵇ}{\ifmmode^{b}\else\textsuperscript{\ensuremath{b}}\fi}
\newunicodechar{ʳ}{\ifmmode^{r}\else\textsuperscript{\ensuremath{r}}\fi}
\newunicodechar{ᵘ}{\ifmmode^{u}\else\textsuperscript{\ensuremath{u}}\fi}
\newunicodechar{ʸ}{\ifmmode^{y}\else\textsuperscript{\ensuremath{y}}\fi}
\newunicodechar{ᵃ}{\ifmmode^{a}\else\textsuperscript{\ensuremath{a}}\fi}
\newunicodechar{ᵉ}{\ifmmode^{e}\else\textsuperscript{\ensuremath{e}}\fi}
\newunicodechar{ᵏ}{\ifmmode^{k}\else\textsuperscript{\ensuremath{k}}\fi}
\newunicodechar{ᵖ}{\ifmmode^{p}\else\textsuperscript{\ensuremath{p}}\fi}
\newunicodechar{ᴺ}{\ifmmode^{N}\else\textsuperscript{\ensuremath{N}}\fi}
\newunicodechar{ᶜ}{\ifmmode^{c}\else\textsuperscript{\ensuremath{c}}\fi}
\newunicodechar{ʰ}{\ifmmode^{h}\else\textsuperscript{\ensuremath{h}}\fi}
\newunicodechar{ᵠ}{\ifmmode^{q}\else\textsuperscript{\ensuremath{q}}\fi}
\newunicodechar{ₙ}{\ifmmode_{n}\else\textsubscript{\ensuremath{n}}\fi}
\newunicodechar{ₐ}{\ifmmode_{a}\else\textsubscript{\ensuremath{a}}\fi}
\newunicodechar{ᵢ}{\ifmmode_{i}\else\textsubscript{\ensuremath{i}}\fi}
\newunicodechar{ᵣ}{\ifmmode_{r}\else\textsubscript{\ensuremath{r}}\fi}
\newunicodechar{ₖ}{\ifmmode_{k}\else\textsubscript{\ensuremath{k}}\fi}
\newunicodechar{ᵦ}{\ifmmode_{\beta}\else\textsubscript{\ensuremath{\beta}}\fi}
\newunicodechar{ₓ}{\ifmmode_{x}\else\textsubscript{\ensuremath{x}}\fi}
\newfontfamily\popdisplay{ArchivoBlack-Regular.ttf}[Path=../../../fonts/]
\newfontfamily\poplogo{SpaceGrotesk-Bold.ttf}[Path=../../../fonts/]
\newfontfamily\popsemi{PlusJakartaSans-SemiBold.ttf}[Path=../../../fonts/]
\newfontfamily\popxbold{PlusJakartaSans-ExtraBold.ttf}[Path=../../../fonts/]
\newfontfamily\popxboldls{PlusJakartaSans-ExtraBold.ttf}[Path=../../../fonts/,LetterSpace=3.0]

\setlist[enumerate]{leftmargin=1.7em,itemsep=5pt,topsep=4pt}
\setlist[itemize]{leftmargin=1.7em,itemsep=4pt,topsep=4pt,label={\color{poporange}\textbullet}}
\setlength{\parindent}{0pt}
\setlength{\parskip}{5pt}
\renewcommand{\arraystretch}{1.25}

% pgfplots : style Pop par défaut
\pgfplotsset{
  every axis/.append style={axis line style={popink,line width=1pt},tick style={popink},
    grid style={popline,line width=0.5pt},label style={font=\small},tick label style={font=\footnotesize}},
  popcurve/.style={poporange,line width=1.8pt,no markers,smooth},
}
% Même style disponible dans un \draw TikZ simple (hors pgfplots)
\tikzset{popcurve/.style={poporange,line width=1.8pt}}
\tikzset{popgrid/.style={popline,very thin}}
\colorlet{popcurve}{poporange} % le nom du style est parfois utilisé comme couleur

% ============================================================
% Pill / squiggle
% ============================================================
\newcommand{\poppill}[3]{%
  \tikz[baseline=-0.55ex]\node[draw=popink,line width=1.2pt,rounded corners=2.6pt,%
    fill=#2,inner xsep=6.5pt,inner ysep=3pt]{%
    \popxboldls\fontsize{7}{7}\selectfont\color{#3}\MakeUppercase{#1}};%
}
\newcommand{\popsquiggle}[1]{%
  \tikz[baseline=-0.4ex]{
    \draw[#1,line width=2.6pt,line cap=round]
      plot [smooth] coordinates {(0,0.09) (0.34,0.01) (0.68,0.21) (1.02,0.01) (1.36,0.10)};
  }%
}
% Mot-clé surligné (marqueur orange sous le texte)
\newcommand{\cle}[1]{{\popxbold\color{popdark}#1}}

% ============================================================
% En-tête / pied
% ============================================================
\pagestyle{fancy}
\fancyhf{}
\renewcommand{\headrulewidth}{0pt}
\renewcommand{\footrulewidth}{0pt}
\lhead{\raisebox{-0.15ex}{\poplogo\fontsize{13}{13}\selectfont\color{popink}OnBuch\color{poporange}+}}
\rhead{\poppill{\DOCMATIERE\ \textbullet\ \DOCNIVEAU\ \textbullet\ Chapter \DOCLECON}{white}{popink}}
\lfoot{\popxbold\fontsize{7.6}{7.6}\selectfont\color{popmuted}\MakeUppercase{Signed course OnBuch+}\ \textbullet\ {\popsemi\DOCTITRE}}
\rfoot{\tikz[baseline=-0.6ex]\node[rounded corners=8.5pt,fill=popink,minimum height=17pt,minimum width=17pt,inner xsep=5pt,inner sep=0]{\popxbold\fontsize{8}{8}\selectfont\color{popcream}\thepage\,/\,\pageref*{LastPage}};}

% ============================================================
% Titres
% ============================================================
\newcommand{\chaptitre}[2]{%
  \begin{tcolorbox}[enhanced,colback=poporange,colframe=popink,boxrule=2.6pt,arc=11pt,
      drop shadow=popink,left=15pt,right=15pt,top=12pt,bottom=11pt,before skip=3pt,after skip=18pt]
    \poppill{#2}{popink}{popcream}\par\vspace{9pt}
    {\raggedright\hyphenpenalty=10000\exhyphenpenalty=10000\popdisplay\fontsize{22}{24}\selectfont\color{popcream}\MakeUppercase{#1}\par}
  \end{tcolorbox}
}
% Section numérotée : pastille orange avec le numéro + titre Archivo
\newcounter{popsec}
\newcommand{\coursec}[1]{%
  \stepcounter{popsec}\setcounter{popsub}{0}%
  \par\vspace{16pt}\noindent
  \begin{minipage}{\linewidth}
  \tikz[baseline=-0.7ex]\node[draw=popink,line width=1.6pt,fill=poporange,rounded corners=4pt,minimum size=22pt,inner sep=0]{\popdisplay\fontsize{12}{12}\selectfont\color{popcream}\thepopsec};%
  \hspace{8pt}{\popdisplay\fontsize{15}{17}\selectfont\color{popink}\MakeUppercase{#1}}\par
  \vspace{4pt}\noindent{\color{poporange}\rule{36pt}{3.6pt}}
  \end{minipage}\par\nopagebreak\vspace{8pt}
}
\newcounter{popsub}
\newcommand{\courssub}[1]{%
  \stepcounter{popsub}%
  \par\vspace{9pt}\noindent{\popxbold\fontsize{11.5}{13}\selectfont\color{popink}{\color{poporange}\thepopsec.\thepopsub}\hspace{6pt}#1}\par\nopagebreak\vspace{4pt}
}

% ============================================================
% Boîtes pédagogiques — même recette : fond blanc, bordure épaisse colorée,
% ombre dure encre, badge plein. Titre optionnel [#1].
% ============================================================
\tcbset{popbase/.style={breakable,enhanced,colback=white,boxrule=2.3pt,arc=9pt,
  shadow={3pt}{-3pt}{0pt}{popink},left=14pt,right=14pt,top=11pt,bottom=11pt,before skip=11pt,after skip=15pt,
  fontupper=\popsemi\fontsize{10.6}{14.5}\selectfont\color{popink},
  fontlower=\popsemi\fontsize{10.6}{14.5}\selectfont\color{popink}}}
% Coche dessinée (le glyphe ✓ n'existe pas dans les polices du document)
\newcommand{\popcheck}[1]{\tikz[baseline=-0.5ex]\draw[#1,line width=1.6pt,line cap=round,line join=round](-0.13,0.0)--(-0.04,-0.09)--(0.13,0.1);}
\providecommand{\coche}{\popcheck{popgreen}}
\providecommand{\resultat}[1]{\par\smallskip\noindent\fcolorbox{popgreen}{popgreenL}{\parbox{\dimexpr\linewidth-2\fboxsep-2\fboxrule\relax}{\textbf{Result.} #1}}\par\smallskip}
\newcommand{\popboxhead}[3]{\poppill{#1}{#2}{white}\ifx\relax#3\relax\else\hspace{6pt}{\popxbold\fontsize{10.6}{12}\selectfont\color{popink}#3}\fi\par\vspace{7pt}}

\newtcolorbox{aretenir}[1][]{popbase,colframe=popgreen,before upper={\popboxhead{Key points}{popgreen}{#1}}}
\newtcolorbox{exemplebox}[1][]{popbase,colframe=poporange,before upper={\popboxhead{Exemple}{poporange}{#1}}}
\newtcolorbox{definition}[1][]{popbase,colframe=popblue,before upper={\popboxhead{Definition}{popblue}{#1}}}
\newtcolorbox{propriete}[1][]{popbase,colframe=poppurple,before upper={\popboxhead{Property}{poppurple}{#1}}}
\newtcolorbox{methode}[1][]{popbase,colframe=popgold,before upper={\popboxhead{Method}{popgold}{#1}}}
\newtcolorbox{attention}[1][]{popbase,colframe=poppink,before upper={\popboxhead{Warning}{poppink}{#1}}}
\newtcolorbox{experience}[1][]{popbase,colframe=popblue,colback=popblueL,before upper={\popboxhead{Experiment}{popblue}{#1}}}
% « Le savais-tu ? » : carte violette pleine (contexte, culture, Cameroun)
\newtcolorbox{savaistu}[1][]{popbase,colback=poppurple,colframe=popink,
  fontupper=\popsemi\fontsize{10.4}{14.2}\selectfont\color{white},
  before upper={\poppill{Did you know?}{white}{poppurple}\ifx\relax#1\relax\else\hspace{6pt}{\popxbold\color{white}#1}\fi\par\vspace{7pt}}}
% Exercice résolu : énoncé en haut, \tcblower puis la solution sur fond crème
\newcounter{popexo}\newcounter{popexr}
\newtcolorbox{exoresolu}[1][]{popbase,colframe=poporange,colbacklower=popcream,
  segmentation style={popink,line width=1.2pt,dashed},
  before upper={\stepcounter{popexr}\popboxhead{Worked example \thepopexr}{poporange}{#1}},
  before lower={\poppill{Solution}{popink}{popcream}\par\vspace{6pt}}}
% Exercice (non résolu en place) — la correction est dans la partie Corrigés
\newtcolorbox{exercice}[1][]{popbase,colframe=popink,
  before upper={\stepcounter{popexo}\popboxhead{Exercise \thepopexo}{popink}{#1}}}
% Corrigé
\newtcolorbox{corrige}[1][]{popbase,colframe=popgreen,colback=popgreenL,
  before upper={\popboxhead{Solution}{popgreen}{#1}}}
% Objectifs / prérequis / bilan
\newtcolorbox{objectifs}{popbase,colframe=popink,colback=poporangeL,
  before upper={\popboxhead{Chapter objectives}{popink}{}}}
\newtcolorbox{prerequis}{popbase,colframe=popink,colback=popblueL,
  before upper={\popboxhead{Prerequisites}{popblue}{}}}
\newtcolorbox{bilan}{popbase,colframe=popink,colback=white,boxrule=2.8pt,
  before upper={\popboxhead{Fiche bilan}{poporange}{L'essentiel à connaître pour l'examen}}}
% Figure encadrée avec légende
\newtcolorbox{popfigure}[1][]{enhanced,breakable=false,colback=white,colframe=popline,boxrule=1.4pt,arc=8pt,
  left=10pt,right=10pt,top=10pt,bottom=8pt,before skip=10pt,after skip=12pt,halign=center,
  lower separated=false,halign lower=center,
  fontlower=\popsemi\fontsize{9}{11.5}\selectfont\color{popmuted},
  #1}
\newcommand{\legende}[1]{\par\vspace{6pt}{\popsemi\fontsize{9}{11.5}\selectfont\color{popmuted}#1}}

% ============================================================
% Signature OnBuch+
% ============================================================
\newcommand{\poplogoinline}[1]{{\poplogo\fontsize{#1}{#1}\selectfont\color{popink}OnBuch\color{poporange}+}}
\newcommand{\signatureonbuch}{%
  \begin{tcolorbox}[enhanced,colback=popink,colframe=popink,boxrule=2.2pt,arc=10pt,
      drop shadow=poporange,left=14pt,right=14pt,top=10pt,bottom=10pt,before skip=0pt,after skip=0pt]
    \tikz[baseline=-0.7ex]\node[circle,fill=popgreen,draw=popcream,line width=1.3pt,minimum size=22pt,inner sep=0]{\popcheck{white}};%
    \hspace{9pt}\begin{minipage}[c]{0.62\linewidth}
      {\popxbold\fontsize{12}{13}\selectfont\color{popcream}Signed\ \textbullet\ {\poplogo OnBuch\color{poporange}+}}\par\vspace{2pt}
      {\raggedright\popsemi\fontsize{8}{10.5}\selectfont\color{popline}\MakeUppercase{OnBuch+ teaching team}\\\MakeUppercase{Follows the official MINESEC syllabus}\par}
    \end{minipage}\hfill
    {\poplogo\fontsize{22}{22}\selectfont\color{popcream}OnBuch\color{poporange}+}
  \end{tcolorbox}%
}

% ============================================================
% Page de garde
% #1 titre · #2 matière · #3 niveau · #4 module · #5 n° leçon · #6 durée ·
% #7 description / objectifs (texte ou itemize)
% ============================================================
\newcommand{\pagegarde}[7]{%
  \thispagestyle{empty}
  \newgeometry{a4paper,margin=1.7cm,top=1.6cm,bottom=1.4cm}%
  \begin{tikzpicture}[remember picture,overlay]
    \fill[poporange!18] ([xshift=-3.2cm,yshift=-3.6cm]current page.north east) circle (4.6cm);
    \fill[poppurple!14] ([xshift=2.4cm,yshift=7.6cm]current page.south west) circle (3.8cm);
  \end{tikzpicture}%
  {\poplogo\fontsize{30}{30}\selectfont\color{popink}OnBuch\color{poporange}+}\hfill
  \raisebox{0.6ex}{\poppill{Signed course \textbullet\ Premium}{popink}{popcream}}\par
  \vspace{1pt}\popsquiggle{poppurple}\par
  \vspace{8pt}
  \begin{tcolorbox}[enhanced,colback=poporange,colframe=popink,boxrule=2.8pt,arc=13pt,
      drop shadow=popink,left=18pt,right=18pt,top=12pt,bottom=13pt,after skip=10pt]
    \poppill{#2 \textbullet\ #3}{popink}{popcream}\hspace{5pt}\poppill{#4}{white}{popink}\par\vspace{8pt}
    {\popdisplay\fontsize{14}{14}\selectfont\color{popink}CHAPTER #5}\par\vspace{4pt}
    {\raggedright\hyphenpenalty=10000\exhyphenpenalty=10000\popdisplay\fontsize{25}{27}\selectfont\color{popcream}\MakeUppercase{#1}\par}
    \vspace{8pt}
    {\popxbold\fontsize{9.5}{9.5}\selectfont\color{popink}\MakeUppercase{Suggested time: #6}}
  \end{tcolorbox}
  \begin{tcolorbox}[enhanced,colback=white,colframe=popink,boxrule=2.2pt,arc=10pt,
      drop shadow=popink,left=16pt,right=16pt,top=10pt,bottom=10pt,after skip=8pt]
    \poppill{In this chapter}{poppurple}{white}\par\vspace{5pt}
    \popsemi\fontsize{9.3}{12.6}\selectfont\color{popink}#7
  \end{tcolorbox}
  \noindent\begin{minipage}[t]{0.487\linewidth}
    \begin{tcolorbox}[enhanced,colback=popgreen,colframe=popink,boxrule=2.2pt,arc=10pt,
        drop shadow=popink,left=11pt,right=11pt,top=10pt,bottom=10pt,height=3.1cm,valign=center]
      \tcbox[colback=white,colframe=popink,boxrule=1.3pt,arc=5pt,boxsep=0pt,nobeforeafter,
        left=3pt,right=3pt,top=3pt,bottom=3pt,valign=center]{\qrcode[height=1.9cm]{https://play.google.com/store/apps/details?id=com.luvvix.onbuch&hl=en}}%
      \hspace{8pt}\begin{minipage}[c]{0.5\linewidth}
        {\popdisplay\fontsize{11}{12.5}\selectfont\color{white}\MakeUppercase{Download OnBuch}}\par\vspace{4pt}
        {\raggedright\popsemi\fontsize{8.4}{11}\selectfont\color{white}Free \textbullet\ Google Play\\AI tutor, past papers, results\par}
      \end{minipage}
    \end{tcolorbox}
  \end{minipage}\hfill
  \begin{minipage}[t]{0.487\linewidth}
    \begin{tcolorbox}[enhanced,colback=poppurple,colframe=popink,boxrule=2.2pt,arc=10pt,
        drop shadow=popink,left=11pt,right=11pt,top=10pt,bottom=10pt,height=3.1cm,valign=center]
      \tikz[baseline=-0.7ex]\node[circle,fill=white,draw=popink,line width=1.3pt,minimum size=1.6cm,inner sep=0]{\popdisplay\fontsize{24}{24}\selectfont\color{poppurple}+};%
      \hspace{8pt}\begin{minipage}[c]{0.6\linewidth}
        {\popdisplay\fontsize{11}{12.5}\selectfont\color{white}\MakeUppercase{Go OnBuch+}}\par\vspace{4pt}
        {\raggedright\popsemi\fontsize{8.4}{11}\selectfont\color{white}All signed courses, unlimited AI tutor, PDF sheets — OnBuch+ tab in the app\par}
      \end{minipage}
    \end{tcolorbox}
  \end{minipage}
  \vspace{8pt}
  \popcontenu
  \vfill
  \signatureonbuch
  \restoregeometry
  \newpage
}

% Rangée « Ce cours contient » (page de garde)
\newcommand{\poptile}[3]{% #1 couleur #2 gros texte #3 légende
  \begin{tcolorbox}[enhanced,colback=#1,colframe=popink,boxrule=2pt,arc=9pt,shadow={2.5pt}{-2.5pt}{0pt}{popink},
      left=6pt,right=6pt,top=9pt,bottom=9pt,halign=center,valign=center,height=1.75cm,width=\linewidth,nobeforeafter]
    {\popdisplay\fontsize{12.5}{13}\selectfont\color{white}\MakeUppercase{#2}}\par\vspace{3pt}
    {\centering\popsemi\fontsize{7.8}{9.5}\selectfont\color{white}#3\par}
  \end{tcolorbox}}
\newcommand{\popcontenu}{%
  \par\noindent{\popxboldls\fontsize{7.5}{7.5}\selectfont\color{popmuted}THIS SIGNED COURSE CONTAINS}\par\vspace{7pt}
  \noindent\begin{minipage}[t]{0.235\linewidth}\poptile{popblue}{Course}{complete, illustrated, step by step}\end{minipage}\hfill
  \begin{minipage}[t]{0.235\linewidth}\poptile{poporange}{Methods}{and worked examples}\end{minipage}\hfill
  \begin{minipage}[t]{0.235\linewidth}\poptile{poppink}{Exercises}{exam style + solutions}\end{minipage}\hfill
  \begin{minipage}[t]{0.235\linewidth}\poptile{popgreen}{Summary sheet}{the essentials to remember}\end{minipage}\par}

% Sommaire
\newtcolorbox{sommaire}{enhanced,colback=white,colframe=popink,boxrule=2.2pt,arc=10pt,
  drop shadow=popink,left=18pt,right=18pt,top=13pt,bottom=13pt,before skip=6pt,after skip=20pt,
  before upper={\poppill{Contents}{poppurple}{white}\par\vspace{9pt}\popsemi\fontsize{10.6}{14.5}\selectfont\color{popink}}}

% Signature de fin de cours — poussée tout en bas de la dernière page
\newcommand{\findecours}{%
  \par\vfill\nopagebreak
  \begin{center}\popsquiggle{poporange}\end{center}\vspace{4pt}
  \signatureonbuch
}
\AtBeginDocument{\def\llbracket{\mathopen{[\mkern-3.5mu[}}\def\rrbracket{\mathclose{]\mkern-3.5mu]}}} % intervalles d'entiers (glyphe ⟦ absent de la police)
% Glyphes absents de Fira Math : redéfinis à partir de glyphes disponibles
\AtBeginDocument{%
  \def\setminus{\mathbin{\backslash}}\let\smallsetminus\setminus
  \def\vdots{\mathord{\vbox{\baselineskip3.2pt\lineskiplimit0pt\kern4pt\hbox{.}\hbox{.}\hbox{.}}}}%
  \def\ddots{\mathinner{\mkern1mu\raise6.5pt\hbox{.}\mkern2mu\raise3.6pt\hbox{.}\mkern2mu\raise0.7pt\hbox{.}\mkern1mu}}%
  \def\triangle{\mathord{\scalebox{0.9}{$\symup{\Delta}$}}}%
  \def\nabla{\mathord{\raisebox{\depth}{\scalebox{1}[-1]{$\symup{\Delta}$}}}}%
  \def\checkmark{\popcheck{popdark}}%
  \def\star{\mathbin{\ast}}\def\bigstar{\mathord{\ast}}%
  \def\diamond{\mathbin{\scalebox{0.6}{$\Diamond$}}}%
  \def\bigcup{\mathop{\mathchoice{\raisebox{-0.3em}{\scalebox{1.7}{$\cup$}}}{\scalebox{1.25}{$\cup$}}{\cup}{\cup}}\displaylimits}%
  \def\bigcap{\mathop{\mathchoice{\raisebox{-0.3em}{\scalebox{1.7}{$\cap$}}}{\scalebox{1.25}{$\cap$}}{\cap}{\cap}}\displaylimits}%
  \def\lesseqqgtr{\mathrel{\lessgtr}}\def\bigoplus{\mathop{\oplus}\displaylimits}%
}
\providecommand{\ssi}{if and only if} % abréviation fréquente
\AtBeginDocument{\providecommand{\wideparen}{\overparen}} % arc de cercle

%% FIGFIX-BEGIN v5 — correctifs globaux des figures (docs/onbuch-generation/tools/figfix)
\usepackage{adjustbox}\usepackage{etoolbox}\tcbuselibrary{raster}
% toute figure TikZ/pgfplots plus large que la ligne est réduite à la largeur disponible
\BeforeBeginEnvironment{tikzpicture}{\begin{adjustbox}{max width=\linewidth}}
\AfterEndEnvironment{tikzpicture}{\end{adjustbox}}
% légendes pgfplots lisibles par-dessus les courbes
\pgfplotsset{every axis/.append style={legend style={fill=white,fill opacity=0.92,text opacity=1,draw=black!15,inner sep=3pt}}}
% étiquettes d'arêtes (syntaxe edge["..."]) avec fond blanc
\tikzset{every edge quotes/.append style={fill=white,inner sep=1.2pt}}
%% FIGFIX-END
\begin{document}

\pagegarde{The atom}{Physics}{Upper Sixth}{Upper Sixth — Physics}{11}{3 periods}{This chapter explores the structure of the atom, from Rutherford's nuclear model revealed by alpha-particle scattering to the quantised energy levels that explain atomic emission and absorption spectra. It introduces the electron volt, the photon relation E = hf = E2 - E1, and the Schrodinger model of the hydrogen atom, which are essential tools for the GCE Advanced Level examination.
\vspace{4pt}
\begin{multicols}{2}\small
\begin{itemize}
\item By the end of this chapter I can describe the historical atomic models (Dalton, Thomson, Rutherford, Bohr) and state their limitations
\item By the end of this chapter I can describe Rutherford's alpha scattering experiment, its observations and the conclusions drawn about the nucleus
\item By the end of this chapter I can estimate the size of the nucleus from scattering data and explain why most alpha particles pass through undeflected
\item By the end of this chapter I can distinguish continuous, line emission and line absorption spectra and explain their origin
\item By the end of this chapter I can explain energy levels, ground state, excited states and ionisation energy of an atom
\item By the end of this chapter I can define and use the electron volt and convert between eV and joules
\end{itemize}\end{multicols}}

\begin{sommaire}
\begin{multicols}{2}
\begin{enumerate}
\item Atomic models and Rutherford's alpha scattering experiment
\item Emission and absorption spectra
\item Energy levels, the electron volt and the equation E = hf = E2 - E1
\item The Schrodinger model of the hydrogen atom
\item Integration activity
\item Methods and worked examples
\item Exercises
\item Solutions to the exercises
\item Summary sheet
\end{enumerate}
\end{multicols}
\end{sommaire}

\begin{objectifs}
\begin{itemize}
\item By the end of this chapter I can describe the historical atomic models (Dalton, Thomson, Rutherford, Bohr) and state their limitations
\item By the end of this chapter I can describe Rutherford's alpha scattering experiment, its observations and the conclusions drawn about the nucleus
\item By the end of this chapter I can estimate the size of the nucleus from scattering data and explain why most alpha particles pass through undeflected
\item By the end of this chapter I can distinguish continuous, line emission and line absorption spectra and explain their origin
\item By the end of this chapter I can explain energy levels, ground state, excited states and ionisation energy of an atom
\item By the end of this chapter I can define and use the electron volt and convert between eV and joules
\item By the end of this chapter I can apply E = hf = E2 - E1 to calculate photon frequencies, wavelengths and energies in transitions
\item By the end of this chapter I can outline the Schrodinger (quantum mechanical) model of the hydrogen atom and the meaning of quantum numbers
\item By the end of this chapter I can interpret a hydrogen energy level diagram and identify the Lyman, Balmer and Paschen series (GCE exam extra)
\end{itemize}
\end{objectifs}

\begin{prerequis}
\begin{itemize}
\item Wave nature of light: frequency, wavelength and the relation c = f lambda (Lower Sixth)
\item Photoelectric effect and the photon concept E = hf (Lower Sixth)
\item Electric charge, Coulomb force and electric potential energy
\item Kinetic energy and conservation of energy
\item Orders of magnitude and use of scientific notation with SI units
\end{itemize}
\end{prerequis}

\newpage

\coursec{Atomic models and Rutherford's alpha scattering experiment}

When a welder in Douala strikes an arc, the brilliant blue-white light hurts the eyes, yet the sodium street lamps of Yaoundé glow a soft orange, and the neon signs of Bamenda shine red. Why does each gas emit its own characteristic colours, and why does white sunlight, spread out by a CD or a prism, show a continuous rainbow while a gas discharge tube shows only a few sharp coloured lines? Even more puzzling: why does the negatively charged electron not simply spiral into the positive nucleus, as classical physics predicts? The answers lie inside the atom itself, and they were uncovered by bombarding gold foil with alpha particles and by analysing the light atoms emit and absorb. This chapter takes you on that journey, from Rutherford's scattering experiment to the quantum picture of the hydrogen atom.

Before we can explain the light, we must know what the atom looks like. That story begins with a simple question: is matter infinitely divisible, or is there a smallest piece?

\courssub{A brief history of atomic models}

\textbf{Dalton's indivisible sphere (early 19th century).}\par
John Dalton, an English schoolteacher, proposed that each chemical element is made of identical, tiny, \cle{indivisible} solid spheres called atoms. Atoms of different elements differ in mass, and chemical reactions simply rearrange atoms without creating or destroying them. This model explained the laws of chemical combination beautifully, but it said nothing about the \emph{internal structure} of the atom — for Dalton, there was none.

\textbf{Thomson and the discovery of the electron (1897).}\par
Studying cathode rays in evacuated discharge tubes, J. J. Thomson showed that the rays consist of negatively charged particles — \cle{electrons} — far lighter than any atom, and that they are the same whatever gas or metal is used. The electron is therefore a \emph{constituent} of every atom: Dalton's indivisible sphere had been divided.

Since atoms are electrically neutral, the negative electrons must be balanced by positive charge somewhere. Thomson proposed the \textbf{plum-pudding model}: the atom is a sphere of uniformly distributed positive charge (the ``pudding''), with electrons embedded in it like plums. The model is neat, but it was never tested directly — until Ernest Rutherford and his students aimed alpha particles at gold foil.

\begin{attention}[Two models, one test]
Dalton's atom is a hard, featureless ball; Thomson's atom is a soft, diffuse ball of positive charge with electrons inside. Both predict that a fast, heavy, positively charged projectile fired at a thin foil should pass through with only \emph{tiny} deflections, because no concentrated charge exists to push it hard. Rutherford's experiment was designed to check exactly this prediction.
\end{attention}

\courssub{Rutherford's alpha scattering experiment (Geiger and Marsden, 1909)}

\begin{experience}[The gold foil experiment]
In Rutherford's laboratory in Manchester, Hans Geiger and Ernest Marsden set up the following apparatus (see figure below):
\begin{itemize}
\item a radioactive source (radon) emitting \cle{alpha particles} — helium nuclei \ensuremath{\mathrm{^{4}_{2}He}}, charge $+2e$, mass about \SI{6.6e-27}{kg} — at speeds of about $2\times10^{7}\ \mathrm{m\,s^{-1}}$;
\item a thick \textbf{lead block} with a narrow hole, acting as a collimator to produce a fine beam;
\item an extremely \textbf{thin gold foil} (about \SI{1e-6}{m} thick, only a few thousand atoms) as target; gold is used because it can be beaten exceptionally thin and has a large atomic number $Z = 79$;
\item a \textbf{scintillation detector} (a zinc sulphide screen viewed through a microscope): each alpha particle striking it produces a tiny flash of light. The detector can be moved around the foil to count particles at any angle $\theta$;
\item the whole arrangement sits in an \textbf{evacuated chamber}, so that the alpha particles are not scattered by air molecules.
\end{itemize}
\end{experience}

\begin{popfigure}
\begin{tikzpicture}[scale=1.0]
% vacuum chamber
\draw[thick, popline] (-4.2,-2.6) rectangle (4.6,2.6);
\node[popmuted, anchor=north west, font=\small] at (-4.1,2.5) {evacuated chamber};
% lead collimator
\fill[gray!60] (-3.4,-0.9) rectangle (-2.9,-0.12);
\fill[gray!60] (-3.4,0.12) rectangle (-2.9,0.9);
\node[font=\small, align=center] at (-3.15,-1.3) {lead\\collimator};
% source
\fill[poporange] (-4.0,-0.25) rectangle (-3.4,0.25);
\node[font=\small, align=center] at (-4.35,0.9) {alpha\\source};
% beam
\draw[-latex, thick, popblue] (-3.35,0) -- (-0.35,0);
\node[popblue, font=\small, above] at (-1.8,0.05) {$\alpha$ beam};
% gold foil
\draw[line width=2.2pt, popgold] (0,-1.4) -- (0,1.4);
\node[popgold, font=\small, align=center] at (0.75,1.7) {gold foil};
% circular scale
\draw[popline, dashed] (0,0) circle (2.1);
% detector at several positions
\foreach \a/\lab in {20/{$\theta$ small}, 90/{$90^{\circ}$}, 150/{$\theta>90^{\circ}$}} {
  \begin{scope}[rotate=\a]
    \draw[fill=popgreenL, draw=popgreen] (2.1,-0.22) rectangle (2.75,0.22);
  \end{scope}
}
\node[popgreen, font=\small, align=center] at (3.6,1.9) {movable\\detector};
% scattered rays
\draw[-latex, poppurple] (0.1,0) -- (2.0,0.75);
\draw[-latex, poppurple] (0.1,0) -- (0.0,2.0);
\draw[-latex, poppurple] (0.1,0) -- (-1.7,1.35);
% angle arc
\draw[popink] (0.9,0) arc (0:21:0.9);
\node[popink, font=\small] at (1.25,0.18) {$\theta$};
\end{tikzpicture}
\legende{Rutherford's scattering apparatus: a collimated beam of alpha particles strikes a thin gold foil inside an evacuated chamber; the scintillation detector counts particles at each scattering angle $\theta$.}
\end{popfigure}

\courssub{Observations}

Counting the flashes at many angles, Geiger and Marsden found three striking facts:

\begin{enumerate}
\item The \textbf{vast majority} of alpha particles pass straight through the foil with little or no deflection ($\theta$ less than about $1^{\circ}$).
\item A \textbf{small fraction} are deflected through large angles.
\item About \textbf{1 in 8000} alpha particles is scattered through more than $90^{\circ}$ — it effectively \emph{bounces back} towards the source.
\end{enumerate}

Rutherford later said this was as surprising as firing a shell at a piece of tissue paper and having it come back and hit you. No diffuse ``pudding'' of positive charge can reverse the motion of a fast, heavy alpha particle; only a very concentrated charge can.

\courssub{Conclusions: the nuclear atom}

\begin{definition}[The atomic nucleus]
The \cle{nucleus} is the tiny, dense, positively charged centre of the atom. It contains nearly all the atom's mass and all its positive charge $+Ze$ (where $Z$ is the atomic number), concentrated in a region of radius about $10^{-15}$ m to $10^{-14}$ m, while the whole atom has a radius of about $10^{-10}$ m. The light, negatively charged electrons occupy the enormous, almost empty space around it. The nucleus is composed of \cle{protons} (charge $+e$) and \cle{neutrons} (charge $0$), collectively called \cle{nucleons}.
\end{definition}

The three observations translate directly into three conclusions:

\begin{itemize}
\item Most particles pass through undeflected $\Rightarrow$ the atom is \textbf{mostly empty space}.
\item Some particles are strongly deflected $\Rightarrow$ the positive charge is concentrated in a \textbf{very small region}, so that a close encounter produces a huge Coulomb force.
\item A few particles bounce back $\Rightarrow$ nearly all the \textbf{mass} is in that same tiny region, so the alpha particle (much lighter than a gold nucleus) rebounds.
\end{itemize}

\begin{attention}[Atom versus nucleus: a factor of $10^{4}$]
A classic examination trap is to confuse the radius of the \textbf{atom} ($\approx 10^{-10}$ m) with that of the \textbf{nucleus} ($\approx 10^{-14}$ m). If the nucleus were the size of a marble placed at the centre of the Ahmadou Ahidjo stadium, the electrons would be moving around the furthest stands — everything in between is empty space.
\end{attention}

\courssub{Qualitative explanation of the scattering}

The alpha particle carries charge $+2e$ and the nucleus carries charge $+Ze$. Between them acts the \textbf{Coulomb repulsion}
\[
F=\frac{1}{4\pi\varepsilon_{0}}\,\frac{(2e)(Ze)}{r^{2}},
\]
which grows rapidly as the separation $r$ decreases. An alpha particle aimed far from a nucleus feels only a weak, brief push and continues almost straight. One aimed close to a nucleus feels a strong force for a longer time and is deflected through a large angle. One aimed almost head-on is slowed, stopped momentarily, and repelled straight back — this is the rare $\theta > 90^{\circ}$ event.

\begin{attention}[Repulsion, not attraction]
A very common mistake is to say that alpha particles are \emph{attracted} by the nucleus. Both are \textbf{positively charged}, so the force is \textbf{repulsive}. The curved trajectories bend \emph{away} from the nucleus, never towards it.
\end{attention}

\begin{popfigure}
\begin{tikzpicture}[scale=1.05]
% nucleus
\fill[poporange] (0,0) circle (0.14);
\node[font=\small, below=2pt] at (0,-0.14) {nucleus $+Ze$};
% undeflected paths
\draw[-latex, thick, popblue] (-4,1.6) -- (3.2,1.6);
\draw[-latex, thick, popblue] (-4,-1.7) -- (3.2,-1.7);
% moderately deflected
\draw[-latex, thick, poppurple] (-4,0.85) .. controls (-0.6,0.8) and (0.3,0.55) .. (3.2,1.9);
\draw[-latex, thick, poppurple] (-4,-0.8) .. controls (-0.6,-0.75) and (0.3,-0.5) .. (3.2,-1.8);
% strongly deflected
\draw[-latex, thick, popgreen] (-4,0.35) .. controls (-0.9,0.3) and (-0.1,0.2) .. (1.6,2.3);
% back-scattered
\draw[-latex, thick, poppink] (-4,0.08) .. controls (-1.2,0.06) and (-0.55,0.05) .. (-0.45,0.05)
   .. controls (-0.7,0.35) and (-1.6,0.9) .. (-3.4,1.75);
\node[popblue, font=\small, align=center] at (2.2,2.15) {far aim: no\\deflection};
\node[poppurple, font=\small, align=center] at (3.35,-1.15) {near: large\\deflection};
\node[poppink, font=\small, align=center] at (-3.0,2.25) {head-on: bounces back\\($\theta>90^{\circ}$)};
\end{tikzpicture}
\legende{Alpha particle trajectories near a nucleus. The closer the line of approach to the nucleus, the stronger the Coulomb repulsion and the larger the deflection; a nearly head-on particle is scattered backwards.}
\end{popfigure}

\begin{propriete}[Factors affecting large-angle scattering]
The probability of large-angle scattering:
\begin{itemize}
\item \textbf{increases} with the thickness of the foil (more nuclei on the path);
\item \textbf{increases} with the atomic number $Z$ of the target (stronger repulsion);
\item \textbf{decreases} as the kinetic energy of the alpha particles increases (faster particles spend less time near the nucleus and are harder to deflect).
\end{itemize}
This qualitative dependence is examinable; the full derivation of Rutherford's scattering formula is not required.
\end{propriete}

\courssub{Estimating the size of the nucleus}

Consider an alpha particle fired directly at a nucleus of charge $+Ze$. Far away it has kinetic energy $E_k$. As it approaches, the Coulomb repulsion does negative work on it: it slows down, converting kinetic energy into electric potential energy. At the \cle{distance of closest approach} $r_{0}$, the particle is momentarily at rest, so all its initial kinetic energy has become potential energy.

\begin{methode}[Estimating the nuclear radius from closest approach]
\begin{enumerate}
\item Write the initial kinetic energy of the alpha particle: $E_k=\tfrac{1}{2}m_{\alpha}v^{2}$.
\item Write the electric potential energy of the pair (alpha $+2e$, nucleus $+Ze$) at separation $r$:
\[
E_p=\frac{1}{4\pi\varepsilon_{0}}\,\frac{(2e)(Ze)}{r}.
\]
\item At the distance of closest approach, apply conservation of energy: $E_k = E_p(r_0)$.
\item Solve for $r_{0}$:
\[
r_{0}=\frac{1}{4\pi\varepsilon_{0}}\,\frac{2Ze^{2}}{E_k}.
\]
\item Interpret: $r_{0}$ is an \textbf{upper limit} for the nuclear radius, since the particle stops before (or just at) the nuclear surface; if it penetrated the nucleus, the strong nuclear force would alter the interaction.
\end{enumerate}
\end{methode}

\begin{popfigure}
\begin{tikzpicture}[scale=1.0]
\draw[-latex, thick] (0,0) -- (7.2,0) node[below left, font=\small] {separation $r$};
\draw[-latex, thick] (0,0) -- (0,3.4) node[above, font=\small] {energy};
% potential energy curve
\draw[very thick, poppurple, domain=0.85:6.8, samples=60] plot (\x,{2.6/\x});
\node[poppurple, font=\small, align=center] at (5.3,1.15) {$E_p=\dfrac{1}{4\pi\varepsilon_0}\dfrac{2Ze^2}{r}$};
% total energy line
\draw[thick, popblue, dashed] (0,1.55) -- (6.8,1.55);
\node[popblue, font=\small, anchor=west] at (4.1,1.75) {$E_k$ (initial)};
% closest approach
\draw[dashed, poppink] (1.68,0) -- (1.68,1.55);
\fill[poppink] (1.68,1.55) circle (0.07);
\node[poppink, font=\small, below] at (1.68,0) {$r_0$};
\node[font=\small, align=center, popink] at (3.4,0.55) {kinetic energy is converted\\into potential energy};
\draw[-latex, popink] (2.6,0.75) -- (1.9,1.35);
\end{tikzpicture}
\legende{Energy diagram for a head-on collision. As the alpha particle approaches, its kinetic energy falls while the electric potential energy rises; at $r=r_0$ the particle stops momentarily and $E_k=E_p$.}
\end{popfigure}

\begin{exoresolu}[Upper limit for the radius of a gold nucleus]
An alpha particle of kinetic energy $E_k = 5.0\ \mathrm{MeV}$ is fired head-on at a gold nucleus ($Z = 79$). Estimate the distance of closest approach and hence an upper limit for the radius of the gold nucleus.\\
Data: $e = 1.60\times10^{-19}\ \mathrm{C}$, $\dfrac{1}{4\pi\varepsilon_{0}} = 8.99\times10^{9}\ \mathrm{N\,m^{2}\,C^{-2}}$, $1\ \mathrm{MeV} = 1.60\times10^{-13}\ \mathrm{J}$.
\tcblower
\textbf{Step 1 — Convert the energy to joules.}
\[
E_k = 5.0\times1.60\times10^{-13} = 8.0\times10^{-13}\ \mathrm{J}.
\]
\textbf{Step 2 — Apply energy conservation at closest approach.}
\[
E_k=\frac{1}{4\pi\varepsilon_{0}}\frac{(2e)(Ze)}{r_0}
\quad\Longrightarrow\quad
r_0=\frac{1}{4\pi\varepsilon_{0}}\frac{2Ze^{2}}{E_k}.
\]
\textbf{Step 3 — Substitute.}
\[
r_0=\frac{8.99\times10^{9}\times 2\times79\times(1.60\times10^{-19})^{2}}{8.0\times10^{-13}}
=\frac{8.99\times10^{9}\times158\times2.56\times10^{-38}}{8.0\times10^{-13}}.
\]
\[
r_0=\frac{3.64\times10^{-26}}{8.0\times10^{-13}}\approx 4.5\times10^{-14}\ \mathrm{m}.
\]
\textbf{Conclusion.} The gold nucleus has a radius \emph{smaller than about} $4.5\times10^{-14}$ m — roughly $10^{4}$ times smaller than the atom ($\approx10^{-10}$ m). This confirms that the atom is overwhelmingly empty space.
\end{exoresolu}

\courssub{Rutherford's planetary model — and its failure}

To complete the picture, Rutherford proposed the \textbf{nuclear (planetary) model}: the electrons orbit the nucleus as planets orbit the Sun, the Coulomb attraction providing the centripetal force:
\[
\frac{1}{4\pi\varepsilon_{0}}\frac{e^{2}}{r^{2}}=\frac{m_e v^{2}}{r}.
\]
The model explains the scattering data perfectly, but it contains a fatal flaw. According to classical electromagnetism, \emph{any accelerated charge radiates energy}. An orbiting electron is permanently accelerated (centripetal acceleration), so it should continuously emit electromagnetic radiation, lose energy, and \textbf{spiral into the nucleus} in a tiny fraction of a second.

\begin{attention}[Two failures of the planetary model]
\begin{itemize}
\item \textbf{Stability:} the radiating electron should collapse into the nucleus — yet atoms are stable.
\item \textbf{Spectra:} a spiralling electron would emit a \emph{continuous} range of frequencies — yet gases emit only a few \emph{sharp, discrete} lines (the orange of sodium, the red of neon).
\end{itemize}
Classical physics cannot save the planetary atom. A new idea is needed: energy in the atom is \emph{quantised}. That is the subject of the next sections.
\end{attention}

\begin{savaistu}[From Manchester to your examination]
Rutherford, a New Zealander working in Manchester, is often called the father of nuclear physics. His scattering technique — probing something too small to see by firing particles at it and studying how they bounce — is still the basic method of particle physics today, in machines such as the Large Hadron Collider, where protons replace alpha particles and gold nuclei are replaced by other protons.
\end{savaistu}

\begin{aretenir}[Key points of this section]
\begin{itemize}
\item Dalton: indivisible spheres; Thomson: electrons embedded in a positive ``pudding''.
\item Rutherford's experiment: alpha particles on thin gold foil in vacuum, detected by scintillations at all angles.
\item Observations: most particles pass straight through; a few are deflected strongly; about 1 in 8000 bounces back.
\item Conclusions: the atom is mostly empty space; a tiny, dense, positive \cle{nucleus} ($\approx10^{-14}$ m) holds nearly all the mass; the atom is $\approx10^{-10}$ m across.
\item Scattering is due to \textbf{Coulomb repulsion} between the positive alpha particle and the positive nucleus; closer approach means larger deflection.
\item Nuclear size estimate: equate initial kinetic energy to potential energy at closest approach, $E_k=\dfrac{1}{4\pi\varepsilon_{0}}\dfrac{2Ze^{2}}{r_0}$.
\item The planetary model fails: an orbiting electron should radiate and spiral into the nucleus; it cannot explain atomic stability or line spectra.
\end{itemize}
\end{aretenir}

\begin{exercice}[Checking your understanding]
\begin{enumerate}
\item State the three main observations of the Geiger--Marsden experiment and the conclusion drawn from each.
\item Explain why the scattering chamber must be evacuated.
\item An alpha particle of kinetic energy $7.7\ \mathrm{MeV}$ approaches a silver nucleus ($Z = 47$) head-on. Calculate the distance of closest approach. ($1\ \mathrm{MeV}=1.60\times10^{-13}\ \mathrm{J}$.)
\item Why would an aluminium foil ($Z = 13$) of the same thickness produce fewer large-angle scatterings than the gold foil?
\end{enumerate}
\end{exercice}

\begin{corrige}[Exercise 1]
\begin{enumerate}
\item Most alpha particles pass through undeflected $\Rightarrow$ the atom is mostly empty space. A few are deflected through large angles $\Rightarrow$ the positive charge is concentrated in a very small region. About 1 in 8000 rebounds $\Rightarrow$ nearly all the mass is concentrated in that tiny nucleus.
\item Air molecules would scatter and absorb the alpha particles, masking the scattering by the foil itself.
\item $r_0=\dfrac{8.99\times10^{9}\times2\times47\times(1.60\times10^{-19})^{2}}{7.7\times1.60\times10^{-13}}=\dfrac{2.16\times10^{-26}}{1.23\times10^{-12}}\approx1.8\times10^{-14}\ \mathrm{m}$.
\item The Coulomb repulsion is proportional to $Z$; with a smaller $Z$, the force on the alpha particle is weaker, so large deflections are less probable.
\end{enumerate}
\end{corrige}

\coursec{Emission and absorption spectra}

Rutherford's scattering experiment told us \emph{where} the mass and positive charge of the atom sit: in a tiny, dense nucleus. But it said nothing about how the electrons are arranged around that nucleus, nor about the energy they carry. The next great clue did not come from bombarding atoms with particles, but from something far gentler: simply \textbf{looking carefully at the light that matter emits and absorbs}. When we spread light out into its component colours — when we produce a \cle{spectrum} — matter reveals a hidden barcode that no chemical test can imitate. This section is about reading that barcode.

\courssub{Producing a spectrum: dispersion and the spectrometer}

\textbf{Observation.} Hold a glass prism in a sunbeam near a window in Yaounde on a bright afternoon and a rainbow appears on the wall. The white sunlight has been separated into its colours. This separation is called \cle{dispersion}: the refractive index of glass depends slightly on wavelength, so each colour is deviated through a slightly different angle (violet most, red least).

A \cle{diffraction grating} — a slide ruled with thousands of equally spaced parallel lines per millimetre — does the same job by interference rather than refraction, and usually spreads the colours out much more, allowing precise measurement of wavelengths using the grating equation $d\sin\theta = n\lambda$.

\begin{definition}[Spectrum and wavelength]
A \cle{spectrum} is the display obtained when light is separated into its component wavelengths. Each colour corresponds to a definite wavelength $\lambda$ (in metres, often quoted in nanometres, \SI{1}{nm}=\SI{e-9}{m}) and frequency $f$, linked by
\[ c = f\lambda, \qquad c = \SI{3.00e8}{m.s^{-1}}. \]
\end{definition}

To study spectra quantitatively we use a \cle{spectrometer}, which has three essential parts:

\begin{itemize}
\item a \textbf{collimator}: a narrow slit placed at the focus of a converging lens, producing a parallel beam from the source;
\item a \textbf{dispersing element}: a prism or a diffraction grating, which sends different wavelengths in different directions;
\item a \textbf{telescope} (or screen/camera): a second lens system that focuses each direction to a sharp image of the slit — one image per wavelength present.
\end{itemize}

\begin{popfigure}
\begin{tikzpicture}[scale=1.0, >=stealth]
% source and slit
\fill[popgold] (0,0) circle (0.12);
\node[below, popdark] at (0,-0.2) {\small source};
\draw[thick, popink] (1.1,-0.5) -- (1.1,-0.08);
\draw[thick, popink] (1.1,0.08) -- (1.1,0.5);
\node[below, popdark] at (1.1,-0.55) {\small slit};
% collimator lens
\draw[thick, popblue] (2.2,-0.7) -- (2.2,0.7);
\draw[popblue] (2.2,0.7) -- (2.1,0.55) (2.2,0.7) -- (2.3,0.55);
\draw[popblue] (2.2,-0.7) -- (2.1,-0.55) (2.2,-0.7) -- (2.3,-0.55);
\node[above, popblue] at (2.2,0.75) {\small collimator};
% prism
\draw[thick, poppurple, fill=poppurple!15] (4.6,0.9) -- (3.9,-0.7) -- (5.3,-0.7) -- cycle;
\node[popdark] at (4.6,-1.0) {\small prism / grating};
% telescope lens
\draw[thick, popgreen] (7.0,-0.7) -- (7.0,0.7);
\draw[popgreen] (7.0,0.7) -- (6.9,0.55) (7.0,0.7) -- (7.1,0.55);
\draw[popgreen] (7.0,-0.7) -- (6.9,-0.55) (7.0,-0.7) -- (7.1,-0.55);
\node[above, popgreen] at (7.0,0.75) {\small telescope};
% screen
\draw[very thick, popink] (8.6,-0.9) -- (8.6,0.9);
\node[right, popdark] at (8.65,0) {\small screen};
% rays: white then split
\draw[thick, popgold] (0.12,0) -- (1.1,0);
\draw[thick, popgold] (1.1,0.02) -- (2.2,0.25);
\draw[thick, popgold] (1.1,-0.02) -- (2.2,-0.25);
\draw[thick, popgold] (2.2,0.25) -- (4.35,0.25);
\draw[thick, popgold] (2.2,-0.25) -- (4.35,-0.25);
% dispersed rays
\draw[thick, red] (4.9,-0.1) -- (7.0,0.15) -- (8.6,0.55);
\draw[thick, green!60!black] (4.9,-0.1) -- (7.0,0.0) -- (8.6,0.1);
\draw[thick, blue!70!black] (4.9,-0.1) -- (7.0,-0.15) -- (8.6,-0.4);
% spectrum marks on screen
\fill[red] (8.6,0.55) circle (0.06);
\fill[green!60!black] (8.6,0.1) circle (0.06);
\fill[blue!70!black] (8.6,-0.4) circle (0.06);
\end{tikzpicture}
\legende{A simple prism spectrometer: the collimator makes a parallel beam, the prism disperses it, and the telescope focuses each wavelength to a separate line image on the screen.}
\end{popfigure}

\courssub{The three types of spectra}

When we point a spectrometer at different light sources, exactly \textbf{three} kinds of spectra appear. Learning to recognise them instantly is essential for the GCE.

\textbf{1. The continuous spectrum.} Heat a solid until it glows — the tungsten filament of an incandescent lamp, the iron of a blacksmith's forge in Douala — and its spectrum is an \emph{unbroken band of colour} from red to violet. Every wavelength in a broad range is present.

\begin{definition}[Continuous spectrum]
A \cle{continuous spectrum} contains all wavelengths over a broad range, with no gaps. It is emitted by \textbf{hot solids, hot liquids, and dense (high-pressure) gases}. The Sun's photosphere and a filament lamp are standard examples.
\end{definition}

\textbf{2. The line emission spectrum.} Now replace the source by a \textbf{gas at low pressure} in a discharge tube (a glass tube with two electrodes, like a neon sign) or by a flame coloured by a metallic salt. The rainbow collapses: instead of a continuous band you see a few \textbf{bright coloured lines on a dark background}, each line being an image of the slit at one particular wavelength.

\begin{definition}[Line spectrum]
A \cle{line spectrum} is a spectrum consisting of discrete bright (\emph{emission}) or dark (\emph{absorption}) lines at definite wavelengths, characteristic of an element.
\end{definition}

Hydrogen, for example, shows in the visible region a red line ($\lambda = \SI{656}{nm}$), a blue-green line ($\SI{486}{nm}$), and violet lines ($\SI{434}{nm}$, $\SI{410}{nm}$) — and nothing in between. Sodium vapour is dominated by a brilliant yellow doublet near $\SI{589}{nm}$, which is why sodium street lamps bathe everything in yellow.

\textbf{3. The line absorption spectrum.} Finally, let white light (which would give a continuous spectrum) pass through a \textbf{cooler gas} before entering the spectrometer. The continuous band reappears, but it is now \emph{crossed by dark lines} — narrow wavelengths missing from the rainbow.

\begin{propriete}[Kirchhoff's laws of spectroscopy]
\begin{enumerate}
\item A hot solid, liquid, or dense gas produces a \textbf{continuous spectrum}.
\item A hot, low-pressure gas produces a \textbf{line emission spectrum} (bright lines on dark background).
\item A continuous spectrum source viewed through a cooler, low-pressure gas produces a \textbf{line absorption spectrum} (dark lines on continuous background). The dark lines occur at \textbf{exactly the same wavelengths} as the bright lines in the emission spectrum of the same gas.
\end{enumerate}
\end{propriete}

\begin{popfigure}
\begin{tikzpicture}[scale=1.0]
% Continuous spectrum
\shade[left color=red, right color=blue!80!black] (0,3.4) rectangle (8,4.2);
\draw[popink] (0,3.4) rectangle (8,4.2);
\node[left, popdark] at (0,3.8) {\small continuous};
% Emission spectrum (hydrogen-like)
\fill[black] (0,1.9) rectangle (8,2.7);
\draw[line width=1.6pt, red] (6.2,1.9) -- (6.2,2.7);
\draw[line width=1.6pt, green!50!blue] (4.1,1.9) -- (4.1,2.7);
\draw[line width=1.6pt, blue!70!black] (3.3,1.9) -- (3.3,2.7);
\draw[line width=1.6pt, violet] (2.9,1.9) -- (2.9,2.7);
\node[left, popdark] at (0,2.3) {\small emission};
% Absorption spectrum
\shade[left color=red, right color=blue!80!black] (0,0.4) rectangle (8,1.2);
\draw[line width=2.2pt, black] (6.2,0.4) -- (6.2,1.2);
\draw[line width=2.2pt, black] (4.1,0.4) -- (4.1,1.2);
\draw[line width=2.2pt, black] (3.3,0.4) -- (3.3,1.2);
\draw[line width=2.2pt, black] (2.9,0.4) -- (2.9,1.2);
\draw[popink] (0,0.4) rectangle (8,1.2);
\node[left, popdark] at (0,0.8) {\small absorption};
% wavelength axis
\draw[->, popink] (0,-0.05) -- (8.4,-0.05);
\node[below, popdark] at (2.9,-0.1) {\small $\SI{410}{nm}$};
\node[below, popdark] at (6.2,-0.1) {\small $\SI{656}{nm}$};
\node[below, popdark] at (4.1,-0.35) {\small $\lambda$ increasing};
\end{tikzpicture}
\legende{The three types of spectra. Top: continuous spectrum of a hot solid. Middle: emission line spectrum of hydrogen (visible Balmer lines). Bottom: absorption spectrum of the same gas — the dark lines sit at exactly the same wavelengths as the bright emission lines.}
\end{popfigure}

\courssub{Why absorption lines appear: a closer look}

The absorption spectrum deserves careful thought, because it hides a subtlety that examiners love to test.

\begin{popfigure}
\begin{tikzpicture}[scale=1.0, >=stealth]
% white light source
\fill[popgold] (0,0) circle (0.14);
\node[below, popdark] at (0,-0.25) {\small white light};
% gas cell
\draw[thick, popblue, fill=popblueL] (2.2,-0.7) rectangle (4.4,0.7);
\node[popdark, align=center] at (3.3,0) {\small cool sodium\\ \small vapour};
% screen
\shade[left color=red, right color=blue!80!black] (6.2,-0.7) rectangle (9.2,0.7);
\draw[line width=2.4pt, black] (7.6,-0.7) -- (7.6,0.7);
\draw[popink] (6.2,-0.7) rectangle (9.2,0.7);
\node[above, popdark] at (7.7,0.8) {\small spectrum with dark line};
% forward beam
\draw[very thick, popgold, ->] (0.2,0) -- (2.2,0);
\draw[very thick, orange!80!red, ->] (4.4,0) -- (6.2,0);
% re-emission arrows in all directions
\draw[->, poporange] (3.3,0.7) -- (3.0,1.4);
\draw[->, poporange] (3.3,0.7) -- (3.9,1.3);
\draw[->, poporange] (2.2,0.2) -- (1.4,0.6);
\draw[->, poporange] (4.4,0.3) -- (5.1,1.0);
\draw[->, poporange] (3.3,-0.7) -- (3.6,-1.4);
\node[poporange, align=center, text width=2.5cm] at (5.6,1.5) {\small absorbed light re-emitted in all directions};
\end{tikzpicture}
\legende{White light crossing cool sodium vapour. Atoms absorb strongly at their characteristic wavelengths and re-emit the energy in random directions, so the forward beam is weakened at those wavelengths: a dark line appears.}
\end{popfigure}

When white light traverses the cool gas, atoms absorb photons of their characteristic wavelengths very strongly. But the atoms do not keep this energy: a tiny fraction of a second later they re-emit it — \textbf{in random directions}. Almost none of the re-emitted light travels forward into the spectrometer, so the forward beam is depleted at exactly those wavelengths: a dark line. The gas has not "ignored" those wavelengths — quite the opposite, it has interacted with them \emph{more} strongly than with any other.

\begin{attention}[Common mistake: misreading the dark lines]
The dark absorption lines do \textbf{not} mean that the gas is transparent to those wavelengths or fails to interact with them. The gas absorbs those wavelengths \emph{very strongly} and then re-emits the energy in all directions, so little of it continues along the original beam. Dark line $=$ strong absorption, not absence of interaction.
\end{attention}

\begin{attention}[Common mistake: confusing continuous and line spectra]
A \textbf{continuous spectrum} contains \emph{every} colour in a range (hot solid, filament lamp). A \textbf{line spectrum} contains only a \emph{few discrete} wavelengths (low-pressure gas). Never describe a discharge-tube spectrum as "all the colours of the rainbow", and never say a filament lamp emits "lines".
\end{attention}

\courssub{Spectra as fingerprints of the elements}

Here is the fact that turned spectroscopy into the most powerful identification tool in science: \textbf{no two elements have the same line spectrum}. Hydrogen's four visible lines, sodium's yellow doublet, mercury's green and violet lines — each element signs its light with a unique pattern.

\begin{methode}[Identifying an element from its spectrum]
\begin{enumerate}
\item Produce the spectrum of the unknown sample (discharge tube, flame, or starlight through a spectrometer).
\item Measure the wavelengths of the bright lines (emission) or dark lines (absorption).
\item Compare these wavelengths with \textbf{reference spectra} tabulated for known elements.
\item A match of several line wavelengths identifies the element unambiguously — even if it is present only as a trace, or located in a star millions of kilometres away.
\end{enumerate}
\end{methode}

\begin{savaistu}[Helium: the element discovered in the Sun before the Earth]
In 1868, astronomers studying the Sun's spectrum during an eclipse found a yellow line that matched \emph{no known element}. They concluded the Sun contained a new substance and named it \textbf{helium}, from the Greek \emph{helios}, "Sun". Helium was only found on Earth 27 years later. The dark lines in the Sun's spectrum — the \cle{Fraunhofer lines} — are absorption lines produced when light from the hot photosphere crosses the cooler gases of the solar atmosphere; they let us determine the chemical composition of a star we can never visit.
\end{savaistu}

\courssub{The electron volt (eV)}

In atomic physics, the joule is inconveniently large. We define a much smaller unit of energy: the \cle{electron volt}.

\begin{definition}[Electron volt]
One \cle{electron volt (\SI{1}{eV})} is the kinetic energy gained by an electron when it is accelerated through a potential difference of \SI{1}{V}.
\[ \SI{1}{eV} = e \times \SI{1}{V} = \SI{1.60e-19}{C} \times \SI{1}{J.C^{-1}} = \SI{1.60e-19}{J}. \]
\end{definition}

\begin{aretenir}[Using the electron volt]
\begin{itemize}
\item Atomic energy levels are typically a few \SI{}{eV} (e.g. hydrogen ground state \SI{-13.6}{eV}).
\item Photon energies in the visible range are \SI{1.8}{eV} (red) to \SI{3.1}{eV} (violet).
\item Always convert to joules when using $E=hf$ with $h=\SI{6.63e-34}{J.s}$, or use $h=\SI{4.14e-15}{eV.s}$.
\end{itemize}
\end{aretenir}

\courssub{The qualitative link to the atom: energy is quantised}

Now comes the deep question: \emph{why} do atoms emit and absorb only certain frequencies? A hot solid emits everything; a rarefied gas of atoms emits a handful of precise notes, like an instrument that can play only a few keys.

The inescapable conclusion — which the next section develops quantitatively — is this:

\begin{aretenir}[The message of line spectra]
Atoms emit or absorb light only at \textbf{specific, well-defined frequencies}. Since light of frequency $f$ carries energy $E = hf$, this means an atom can only gain or lose \textbf{definite amounts of energy}. The internal energy of an atom is \cle{quantised}: it exists only in certain allowed values called \textbf{energy levels}. Each spectral line corresponds to a transition between two such levels.
\end{aretenir}

\begin{popfigure}
\begin{tikzpicture}[scale=1.2, >=stealth]
% Energy levels
\draw[thick, popdark] (0,0) -- (2,0) node[right, popdark] {$E_4 = \SI{-0.85}{eV}$};
\draw[thick, popdark] (0,1.5) -- (2,1.5) node[right, popdark] {$E_3 = \SI{-1.51}{eV}$};
\draw[thick, popdark] (0,3.0) -- (2,3.0) node[right, popdark] {$E_2 = \SI{-3.40}{eV}$};
\draw[thick, popdark] (0,5.0) -- (2,5.0) node[right, popdark] {$E_1 = \SI{-13.6}{eV}$};
% Transitions: Absorption (up arrows)
\draw[->, thick, popblue] (2.5,0) -- (2.5,1.5) node[midway, right, popblue] {Absorption $\Delta E = E_3-E_4$};
\draw[->, thick, popblue] (2.5,1.5) -- (2.5,3.0) node[midway, right, popblue] {Absorption $\Delta E = E_2-E_3$};
\draw[->, thick, popblue] (2.5,3.0) -- (2.5,5.0) node[midway, right, popblue] {Absorption $\Delta E = E_1-E_2$};
% Transitions: Emission (down arrows)
\draw[->, thick, poporange] (3.5,5.0) -- (3.5,3.0) node[midway, right, poporange] {Emission $hf = E_1-E_2$};
\draw[->, thick, poporange] (3.5,3.0) -- (3.5,1.5) node[midway, right, poporange] {Emission $hf = E_2-E_3$};
\draw[->, thick, poporange] (3.5,1.5) -- (3.5,0) node[midway, right, poporange] {Emission $hf = E_3-E_4$};
% Photon wavy lines
\draw[popblue, decorate, decoration={snake, amplitude=1pt, segment length=4pt}] (2.2,0.75) -- (2.8,0.75);
\draw[poporange, decorate, decoration={snake, amplitude=1pt, segment length=4pt}] (3.2,4.0) -- (3.8,4.0);
\node[popblue] at (2.5, -0.3) {Photon absorbed};
\node[poporange] at (3.5, 5.3) {Photon emitted};
\end{tikzpicture}
\legende{Energy level diagram for hydrogen (not to scale). Upward arrows (blue): absorption of a photon of exact energy $\Delta E = hf$. Downward arrows (orange): emission of a photon of exact energy $\Delta E = hf$. The same energy gaps appear in both processes.}
\end{popfigure}

This single idea — quantised energy levels — will explain the hydrogen spectrum numerically and lead us, by the end of the chapter, to Schrödinger's model of the atom.

\courssub{Applications of spectra}

\begin{itemize}
\item \textbf{Flame tests (practical work):} a platinum wire dipped in a metallic salt and held in a Bunsen flame colours it characteristically — sodium: intense yellow; potassium: lilac; calcium: brick red; copper(II): blue-green. This is the line emission spectrum seen by eye.
\item \textbf{Street lamps:} sodium vapour lamps give an efficient yellow light dominated by the $\SI{589}{nm}$ doublet; mercury vapour lamps give a bluish-white light rich in green and violet lines (with ultraviolet that requires a phosphor coating).
\item \textbf{Neon advertising signs:} a discharge through neon at low pressure glows red-orange; other gases or coatings give other colours. The colour is the element's emission spectrum.
\item \textbf{Astronomy:} the Fraunhofer absorption lines in starlight reveal the elements present in stellar atmospheres, their temperature, and even (through small wavelength shifts — Doppler effect) their motion toward or away from us.
\end{itemize}

\begin{exoresolu}[Identifying a gas and finding a frequency]
White light is shone through the cool vapour of an unknown element and then into a spectrometer. The continuous spectrum obtained is crossed by dark lines at $\SI{589.0}{nm}$, $\SI{589.6}{nm}$ and several other wavelengths identical to those of sodium's emission spectrum. (a) Identify the element. (b) Calculate the frequency of the $\SI{589.0}{nm}$ line. (c) Explain why the lines are dark rather than bright.
\tcblower
\textbf{(a)} The dark absorption lines coincide exactly with sodium's emission lines. Since each element has a unique line spectrum, the vapour is \textbf{sodium}.

\textbf{(b)} Using $c = f\lambda$ with $\lambda = \SI{589.0}{nm} = \SI{5.890e-7}{m}$:
\[ f = \frac{c}{\lambda} = \frac{\SI{3.00e8}{m.s^{-1}}}{\SI{5.890e-7}{m}} = \SI{5.09e14}{Hz}. \]

\textbf{(c)} Sodium atoms absorb photons of these wavelengths strongly, then re-emit the energy in \emph{random directions}. Very little of the re-emitted light travels forward into the spectrometer, so the beam is weakened at those wavelengths and the lines appear dark against the continuous background.
\end{exoresolu}

\begin{exercice}[Spectra around us]
\begin{enumerate}
\item Classify each source by the type of spectrum it produces: (i) a candle flame made yellow by sodium impurities; (ii) the glowing filament of a torch bulb; (iii) a low-pressure mercury discharge tube; (iv) starlight crossing the star's cooler atmosphere.
\item A line in the hydrogen spectrum has wavelength $\SI{486}{nm}$. Calculate its frequency and the energy of one photon in joules and in electron volts (use $h = \SI{6.63e-34}{J.s}$, $e = \SI{1.60e-19}{C}$).
\item Explain why the absorption lines of a gas are found at exactly the same wavelengths as its emission lines.
\end{enumerate}
\end{exercice}

\begin{corrige}[Exercise 1]
\textbf{1.} (i) Line emission spectrum (the sodium atoms in the hot, low-density flame emit their yellow doublet). (ii) Continuous spectrum (hot solid). (iii) Line emission spectrum (low-pressure gas in a discharge). (iv) Line absorption spectrum (continuous spectrum crossed by dark Fraunhofer lines).

\textbf{2.} Frequency: $f = \dfrac{c}{\lambda} = \dfrac{\SI{3.00e8}{m.s^{-1}}}{\SI{4.86e-7}{m}} = \SI{6.17e14}{Hz}$.\\
Photon energy in joules: $E = hf = \SI{6.63e-34}{J.s} \times \SI{6.17e14}{Hz} = \SI{4.09e-19}{J}$.\\
Photon energy in eV: $E = \dfrac{\SI{4.09e-19}{J}}{\SI{1.60e-19}{J.eV^{-1}}} = \SI{2.56}{eV}$.

\textbf{3.} An atom can only absorb the precise photon energies that carry it between its allowed energy levels, and these are exactly the energies it releases when it de-excites. Since $E = hf = hc/\lambda$, equal energy jumps mean equal wavelengths: the wavelengths a gas emits are precisely those it absorbs.
\end{corrige}

\begin{aretenir}[Key points of this section]
\begin{itemize}
\item A spectrometer (collimator + prism or grating + telescope) separates light into its wavelengths.
\item \textbf{Continuous spectrum}: all wavelengths; from hot solids, liquids, dense gases.
\item \textbf{Line emission spectrum}: discrete bright lines on a dark background; from low-pressure gases in discharge tubes or flames.
\item \textbf{Line absorption spectrum}: continuous spectrum crossed by dark lines; produced when white light crosses a cooler gas; dark lines match the emission wavelengths exactly (Kirchhoff's 3rd law).
\item Each element has a unique line spectrum — a fingerprint used in flame tests, street lamps, neon signs and the analysis of starlight.
\item Line spectra prove that atoms exchange energy only in definite amounts: atomic energy is \textbf{quantised}.
\item The \textbf{electron volt (\SI{1}{eV} = \SI{1.60e-19}{J})} is the practical energy unit for atomic transitions.
\end{itemize}
\end{aretenir}

\coursec{Energy levels, the electron volt and the equation $E = hf = E_2 - E_1$}

In the previous section we saw that a hot hydrogen gas emits light only at certain well-defined wavelengths, producing a line spectrum, and that a cool gas absorbs exactly those same wavelengths. This observation is incomprehensible if the electron in an atom could have \emph{any} energy. It becomes perfectly natural if the electron can only occupy certain \cle{allowed energies}. In this section we make that idea precise: we introduce \cle{energy levels}, a practical unit of energy at the atomic scale (the \cle{electron volt}), and the fundamental equation $E = hf = E_2 - E_1$ that links every spectral line to a jump between two levels.

\courssub{Quantisation of atomic energy}

Imagine climbing a staircase in the dark: your foot can rest on a step, but never between two steps. The electron in an atom behaves in the same way: it can only possess certain definite values of energy, called \cle{energy levels}.

\begin{definition}[Energy levels, ground state, excited states]
The energy of an electron bound in an atom is \cle{quantised}: it can only take certain discrete values $E_1, E_2, E_3, \dots$ called \textbf{energy levels}. The lowest level $E_1$ is the \textbf{ground state}; all higher levels are \textbf{excited states}. The atom is normally found in its ground state.
\end{definition}

For the hydrogen atom, the allowed energies are given by the remarkably simple formula (which we shall justify in the next section):
\[
E_n = -\frac{13.6}{n^2}\ \text{eV}, \qquad n = 1, 2, 3, \dots
\]
where the integer $n$ is called the \cle{principal quantum number}. Note the crucial features:

\begin{itemize}
\item The energies are \textbf{negative}. This is not a defect: it expresses that the electron is \emph{bound} to the nucleus.
\item The levels get closer and closer together as $n$ increases ($E_1 = -13.6$ eV, $E_2 = -3.40$ eV, $E_3 = -1.51$ eV, $E_4 = -0.85$ eV, \dots), and they accumulate towards $E_\infty = 0$.
\item $E = 0$ corresponds to an electron just removed from the atom: the atom is \textbf{ionised}.
\end{itemize}

\begin{propriete}[Meaning of the negative sign]
The zero of energy is chosen for an electron at rest infinitely far from the nucleus (a \emph{free} electron). A negative energy $E_n < 0$ therefore means the electron is \textbf{bound}: one must \emph{supply} energy $|E_n|$ to free it. The value $E = 0$ is the threshold of \cle{ionisation}.
\end{propriete}

\begin{definition}[Ionisation energy]
The \textbf{ionisation energy} of an atom is the minimum energy required to remove the electron completely from the ground state. For hydrogen:
\[
E_{\text{ion}} = E_\infty - E_1 = 0 - (-13.6\ \text{eV}) = 13.6\ \text{eV}.
\]
\end{definition}

\begin{definition}[Excitation and de-excitation]
\textbf{Excitation}: the atom absorbs energy (from an incident photon or from a collision with another particle) and its electron jumps from a lower level to a higher one. \textbf{De-excitation}: the electron falls back to a lower level and the atom releases the energy difference as a photon. Excited states are unstable: de-excitation occurs spontaneously after a very short time (typically $10^{-8}$ s).
\end{definition}

\begin{popfigure}
\begin{center}
\begin{tikzpicture}[scale=1.0]
% Energy axis
\draw[->, thick, popink] (0,-5.6) -- (0,0.6) node[above, popink]{\small $E$ (eV)};
% Levels (y scaled: y = E/2.72 roughly, using values)
\draw[thick, popblue] (0.6,-5.44) -- (6.4,-5.44) node[right, popink]{\small $E_1=-13.6$};
\node[left, popink] at (0.6,-5.44) {\small $n=1$};
\draw[thick, popblue] (0.6,-3.40) -- (6.4,-3.40) node[right, popink]{\small $E_2=-3.40$};
\node[left, popink] at (0.6,-3.40) {\small $n=2$};
\draw[thick, popblue] (0.6,-2.51) -- (6.4,-2.51) node[right, popink]{\small $E_3=-1.51$};
\node[left, popink] at (0.6,-2.51) {\small $n=3$};
\draw[thick, popblue] (0.6,-2.19) -- (6.4,-2.19) node[right, popink]{\small $E_4=-0.85$};
\node[left, popink] at (0.6,-2.19) {\small $n=4$};
\draw[thick, popblue] (0.6,-2.02) -- (6.4,-2.02) node[right, popink]{\small $E_5=-0.54$};
\node[left, popink] at (0.6,-2.02) {\small $n=5$};
\draw[thick, popdark, dashed] (0.6,-1.6) -- (6.4,-1.6) node[right, popink]{\small $E_\infty=0$};
\node[left, popink] at (0.6,-1.6) {\small $n=\infty$};
% Absorption arrow 1->3
\draw[->, very thick, popgreen] (1.6,-5.44) -- (1.6,-2.51) node[midway, left, popgreen, align=center]{\small absorption};
% Emission arrow 3->2
\draw[->, very thick, poporange] (3.4,-2.51) -- (3.4,-3.40) node[midway, right, poporange, align=center]{\small emission};
% Ionisation arrow
\draw[->, very thick, poppurple] (5.2,-5.44) -- (5.2,-1.6) node[midway, right, poppurple, align=center]{\small ionisation};
\end{tikzpicture}
\end{center}
\legende{Energy level diagram of the hydrogen atom. Levels crowd together as $n$ increases and converge to $E_\infty = 0$ (ionisation). Upward arrows: absorption or ionisation; downward arrows: emission.}
\end{popfigure}

\courssub{The electron volt}

Energies at the atomic scale are tiny: the ionisation energy of hydrogen is about $2.18 \times 10^{-18}$ J. Working in joules is clumsy, so physicists use a unit adapted to the atomic scale.

\begin{definition}[The electron volt]
One \textbf{electron volt} (symbol eV) is the kinetic energy gained by a particle of charge $e$ (an electron, for example) accelerated from rest through a potential difference of $1$ V:
\[
1\ \text{eV} = e \times 1\ \text{V} = 1.60 \times 10^{-19}\ \text{J}.
\]
Conversely, $1\ \text{J} = \dfrac{1}{1.60\times 10^{-19}}\ \text{eV} = 6.25\times 10^{18}\ \text{eV}$.
\end{definition}

The definition follows directly from the work of the electric force: a charge $q$ accelerated through a p.d. $U$ gains kinetic energy $E_c = qU$. Setting $q = e = 1.60\times 10^{-19}$ C and $U = 1$ V gives $E_c = 1.60\times 10^{-19}$ J.

\begin{exemplebox}[Converting between eV and joules]
The first excited level of hydrogen is $E_2 = -3.40$ eV. In joules:
\[
E_2 = -3.40 \times 1.60\times 10^{-19} = -5.44\times 10^{-19}\ \text{J}.
\]
Conversely, a photon of energy $3.31\times 10^{-19}$ J has, in eV:
\[
E = \frac{3.31\times 10^{-19}}{1.60\times 10^{-19}} = 2.07\ \text{eV}.
\]
\end{exemplebox}

\courssub{The photon relation $E = hf = E_2 - E_1$}

We now connect the quantised levels to the line spectra of the previous section. When an electron falls from a level $E_2$ to a lower level $E_1$, conservation of energy requires the atom to lose exactly $\Delta E = E_2 - E_1$. This energy is carried away by a single photon.

\begin{propriete}[Bohr's frequency condition (examinable)]
A photon emitted or absorbed during an atomic transition between two levels $E_2$ (higher) and $E_1$ (lower) has energy
\[
\[\boxed{\;E = hf = E_2 - E_1\;}\]
\]
where $h = 6.63\times 10^{-34}$ J s is Planck's constant and $f$ the photon's frequency. Since $f = c/\lambda$, the wavelength of the emitted or absorbed light is
\[
\lambda = \frac{hc}{E_2 - E_1}.
\]
Absorption is possible \textbf{only} if the photon energy $hf$ is \emph{exactly} equal to a difference between two levels.
\end{propriete}

This single equation explains everything observed in the previous section: since only certain level differences exist, only certain frequencies (hence certain lines) appear in emission and absorption spectra.

\begin{methode}[Solving a transition problem]
\begin{enumerate}
\item Identify the two levels and compute the energy difference $\Delta E = E_2 - E_1$ in eV (take the absolute value; emission if the electron falls, absorption if it rises).
\item Convert $\Delta E$ to joules: multiply by $1.60\times 10^{-19}$ J/eV.
\item Compute the frequency $f = \dfrac{\Delta E}{h}$ with $h = 6.63\times 10^{-34}$ J s.
\item Compute the wavelength $\lambda = \dfrac{c}{f} = \dfrac{hc}{\Delta E}$ with $c = 3.00\times 10^{8}$ m s$^{-1}$.
\item Check the order of magnitude: visible light corresponds to $\lambda$ between about 400 nm and 750 nm.
\end{enumerate}
\end{methode}

\begin{exoresolu}[The red line of hydrogen]
The transition of a hydrogen atom from level $n=3$ ($E_3 = -1.51$ eV) to level $n=2$ ($E_2 = -3.40$ eV) produces a well-known red line. Calculate (a) the energy of the emitted photon in eV and in joules; (b) its frequency; (c) its wavelength.
\tcblower
\textbf{(a) Energy of the photon.}
\[
\Delta E = E_3 - E_2 = -1.51 - (-3.40) = 1.89\ \text{eV}.
\]
In joules:
\[
\Delta E = 1.89 \times 1.60\times 10^{-19} = 3.02\times 10^{-19}\ \text{J}.
\]
\textbf{(b) Frequency.}
\[
f = \frac{\Delta E}{h} = \frac{3.02\times 10^{-19}}{6.63\times 10^{-34}} = 4.56\times 10^{14}\ \text{Hz}.
\]
\textbf{(c) Wavelength.}
\[
\lambda = \frac{c}{f} = \frac{3.00\times 10^{8}}{4.56\times 10^{14}} = 6.58\times 10^{-7}\ \text{m} = 658\ \text{nm}.
\]
This is indeed in the red region of the visible spectrum — the famous Balmer-$\alpha$ line, the strongest line in the visible spectrum of hydrogen.
\end{exoresolu}

\begin{attention}[Convert eV to joules before using $E = hf$]
Planck's constant $h = 6.63\times 10^{-34}$ J s is expressed in \textbf{joules}. If you insert an energy in eV directly into $f = E/h$, your frequency will be wrong by a factor $1.6\times 10^{19}$. Always convert: $\Delta E\,(\text{J}) = \Delta E\,(\text{eV}) \times 1.60\times 10^{-19}$.
\end{attention}

\begin{attention}[Not every photon can excite an atom]
A photon is absorbed \textbf{only} if its energy matches \emph{exactly} a level difference. For hydrogen in the ground state, the smallest absorbable energy is $E_2 - E_1 = 10.2$ eV. A photon of 9 eV or 11 eV passes through without any effect; a photon of exactly 10.2 eV is absorbed and lifts the electron to $n=2$. By contrast, a \emph{collision} with another particle can excite the atom even if the available energy exceeds the level difference, because the surplus can be carried away as kinetic energy.
\end{attention}

\begin{exoresolu}[Identifying a transition from a measured wavelength]
A line in the spectrum of atomic hydrogen is measured at $\lambda = 434$ nm. Using the levels $E_2 = -3.40$ eV, $E_3 = -1.51$ eV, $E_4 = -0.85$ eV, $E_5 = -0.54$ eV, identify the transition responsible.
\tcblower
Energy of the photon:
\[
\Delta E = \frac{hc}{\lambda} = \frac{6.63\times 10^{-34}\times 3.00\times 10^{8}}{434\times 10^{-9}} = 4.58\times 10^{-19}\ \text{J}.
\]
Converting to eV:
\[
\Delta E = \frac{4.58\times 10^{-19}}{1.60\times 10^{-19}} = 2.86\ \text{eV}.
\]
Testing the level differences from $E_2$: $E_3 - E_2 = 1.89$ eV; $E_4 - E_2 = 2.55$ eV; $E_5 - E_2 = 2.86$ eV. The line corresponds to the transition $n=5 \rightarrow n=2$.
\end{exoresolu}

\courssub{The spectral series of hydrogen}

All transitions of hydrogen can be grouped into \cle{series}, according to the \emph{final} level of the electron:

\begin{itemize}
\item \textbf{Lyman series}: transitions down to $n = 1$. The energies are large (from $10.2$ eV up to $13.6$ eV), so the lines lie in the \textbf{ultraviolet}.
\item \textbf{Balmer series}: transitions down to $n = 2$. The energies (from $1.89$ eV to $3.40$ eV) correspond to \textbf{visible} light — these are the lines seen with a school spectroscope.
\item \textbf{Paschen series}: transitions down to $n = 3$. These lines lie in the \textbf{infrared}.
\end{itemize}

\begin{popfigure}
\begin{center}
\begin{tikzpicture}[scale=1.0]
\draw[->, thick, popink] (0,-5.6) -- (0,0.4) node[above, popink]{\small $E$};
\draw[thick, popblue] (0.6,-5.44) -- (8.6,-5.44);
\node[left, popink] at (0.6,-5.44) {\small $n=1$};
\draw[thick, popblue] (0.6,-3.40) -- (8.6,-3.40);
\node[left, popink] at (0.6,-3.40) {\small $n=2$};
\draw[thick, popblue] (0.6,-2.51) -- (8.6,-2.51);
\node[left, popink] at (0.6,-2.51) {\small $n=3$};
\draw[thick, popblue] (0.6,-2.19) -- (8.6,-2.19);
\node[left, popink] at (0.6,-2.19) {\small $n=4$};
\draw[thick, popblue] (0.6,-2.02) -- (8.6,-2.02);
\node[left, popink] at (0.6,-2.02) {\small $n=5$};
\draw[thick, popdark, dashed] (0.6,-1.6) -- (8.6,-1.6);
\node[left, popink] at (0.6,-1.6) {\small $n=\infty$};
% Lyman (to n=1)
\draw[->, thick, poppurple] (1.8,-3.40) -- (1.8,-5.44);
\draw[->, thick, poppurple] (2.4,-2.51) -- (2.4,-5.44);
\draw[->, thick, poppurple] (3.0,-2.19) -- (3.0,-5.44);
\node[poppurple, align=center] at (2.4,-5.9) {\small Lyman (UV)};
% Balmer (to n=2)
\draw[->, thick, poporange] (4.6,-2.51) -- (4.6,-3.40);
\draw[->, thick, poporange] (5.2,-2.19) -- (5.2,-3.40);
\draw[->, thick, poporange] (5.8,-2.02) -- (5.8,-3.40);
\node[poporange, align=center] at (5.2,-5.9) {\small Balmer (visible)};
% Paschen (to n=3)
\draw[->, thick, popgreen] (7.2,-2.19) -- (7.2,-2.51);
\draw[->, thick, popgreen] (7.8,-2.02) -- (7.8,-2.51);
\node[popgreen, align=center] at (7.5,-5.9) {\small Paschen (IR)};
\end{tikzpicture}
\end{center}
\legende{The Lyman, Balmer and Paschen series of hydrogen. Each series gathers all transitions ending on the same lower level.}
\end{popfigure}

\begin{popfigure}
\begin{center}
\begin{tikzpicture}
\begin{axis}[
width=11cm, height=5.2cm,
xlabel={$\lambda$ (nm)}, ylabel={relative intensity},
xmin=380, xmax=700, ymin=0, ymax=1.2,
xtick={400,500,600,700}, ytick=\empty,
axis lines=left, thick]
\addplot[ycomb, poporange, very thick] coordinates {(656,1.0)};
\addplot[ycomb, popblue, very thick] coordinates {(486,0.55)};
\addplot[ycomb, poppurple, very thick] coordinates {(434,0.30)};
\addplot[ycomb, poppink, very thick] coordinates {(410,0.18)};
\node[poporange, align=center] at (axis cs:656,1.08) {\small $3\to2$};
\node[popblue, align=center] at (axis cs:486,0.63) {\small $4\to2$};
\node[poppurple, align=center] at (axis cs:434,0.38) {\small $5\to2$};
\node[poppink, align=center] at (axis cs:410,0.26) {\small $6\to2$};
\end{axis}
\end{tikzpicture}
\end{center}
\legende{Sketch of the visible Balmer lines. Each line corresponds to one transition ending on $n=2$: 656 nm (red), 486 nm (blue-green), 434 nm and 410 nm (violet).}
\end{popfigure}

\begin{exercice}[Transitions in hydrogen]
Take $E_n = -13.6/n^2$ eV, $h = 6.63\times 10^{-34}$ J s, $c = 3.00\times 10^{8}$ m s$^{-1}$, $e = 1.60\times 10^{-19}$ C.
\begin{enumerate}
\item Calculate the wavelength of the photon emitted in the Lyman-$\alpha$ transition ($n=2 \to n=1$). In which region of the spectrum does it lie?
\item A hydrogen atom in the ground state absorbs a photon of wavelength $97.3$ nm. To which level is the electron raised?
\item What is the shortest wavelength that hydrogen can emit in the Balmer series?
\end{enumerate}
\end{exercice}

\begin{corrige}[Exercise 1]
\textbf{1.} $\Delta E = E_2 - E_1 = -3.40-(-13.6) = 10.2$ eV $= 10.2\times 1.60\times 10^{-19} = 1.63\times 10^{-18}$ J. Then
\[
\lambda = \frac{hc}{\Delta E} = \frac{6.63\times 10^{-34}\times 3.00\times 10^{8}}{1.63\times 10^{-18}} = 1.22\times 10^{-7}\ \text{m} = 122\ \text{nm}.
\]
This is ultraviolet ($\lambda < 400$ nm), as expected for the Lyman series.

\textbf{2.} Photon energy:
\[
\Delta E = \frac{hc}{\lambda} = \frac{6.63\times 10^{-34}\times 3.00\times 10^{8}}{97.3\times 10^{-9}} = 2.04\times 10^{-18}\ \text{J} = 12.8\ \text{eV}.
\]
The electron reaches $E_1 + 12.8 = -13.6 + 12.8 = -0.8$ eV $\approx E_4 = -0.85$ eV: the electron is raised to the level $n=4$.

\textbf{3.} The shortest Balmer wavelength corresponds to the largest energy jump ending on $n=2$, i.e. from $n=\infty$: $\Delta E = 0-(-3.40) = 3.40$ eV $= 5.44\times 10^{-19}$ J, giving
\[
\lambda_{\min} = \frac{hc}{\Delta E} = \frac{6.63\times 10^{-34}\times 3.00\times 10^{8}}{5.44\times 10^{-19}} = 3.66\times 10^{-7}\ \text{m} = 366\ \text{nm},
\]
at the violet/ultraviolet edge of the visible.
\end{corrige}

\begin{aretenir}[Key points]
\begin{itemize}
\item Atomic energies are quantised: $E_n = -13.6/n^2$ eV for hydrogen; negative energies mean a bound electron, $E=0$ means ionisation (13.6 eV from the ground state).
\item $1\ \text{eV} = 1.60\times 10^{-19}$ J: energy gained by a charge $e$ through 1 V.
\item Every transition obeys Bohr's frequency condition $hf = E_2 - E_1$, with $\lambda = hc/\Delta E$. Convert eV to joules before using $h$ in SI units.
\item Absorption requires an \emph{exact} match between photon energy and a level difference.
\item Hydrogen series: Lyman ($\to n=1$, UV), Balmer ($\to n=2$, visible), Paschen ($\to n=3$, infrared).
\end{itemize}
\end{aretenir}

\begin{savaistu}[Why the sky of a neon sign is not white]
The orange-red glow of advertising tubes and the coloured lamps used in some streets of Douala and Yaoundé come from exactly this physics: an electric discharge excites the gas atoms, which then de-excite by emitting photons at their own characteristic wavelengths. Each gas has its own colour signature — its line spectrum is a true ``barcode'' of the element, which is how astronomers identify the chemical composition of stars millions of kilometres away.
\end{savaistu}

\coursec{The Schrödinger model of the hydrogen atom}

In the previous section we saw how Bohr's model explains the hydrogen spectrum: the electron occupies quantised energy levels $E_n = -\dfrac{13.6}{n^2}\ \mathrm{eV}$, and a transition between two levels emits or absorbs a photon of energy $hf = E_2 - E_1$. This is a remarkable achievement. But a careful scientist must always ask: \emph{is this the whole story?} In this final section we shall see that Bohr's picture of the electron as a tiny planet on a fixed circular orbit is only an approximation, and that the deeper, correct description of the hydrogen atom is given by \cle{wave mechanics}, developed by Erwin Schrödinger in 1926.

\courssub{The limits of the Bohr model}

Bohr's model was a brilliant bridge between classical physics and the quantum world, but it suffers from serious weaknesses:

\begin{itemize}
\item It works \textbf{only for hydrogen} (and hydrogen-like ions with a single electron, such as $\ce{He+}$). It fails completely for atoms with two or more electrons, such as helium.
\item It cannot explain the \textbf{fine details} of spectra: when examined with high-resolution instruments, many spectral lines are in fact doublets or multiplets, which Bohr's single quantum number $n$ cannot account for.
\item It says nothing about the \textbf{intensities} of spectral lines -- why some lines are bright and others faint.
\item Most fundamentally, it treats the electron as a \textbf{classical particle on a fixed orbit}, while simply \emph{imposing} quantisation by hand, with no deeper justification.
\end{itemize}

\begin{attention}[Do not oversell Bohr's model]
In the GCE examination, if you are asked why the Bohr model was replaced, do not say it is ``wrong''. It correctly predicts the hydrogen energy levels. Say rather that it is \textbf{incomplete}: it cannot treat multi-electron atoms, fine structure or line intensities, and it gives no reason \emph{why} angular momentum should be quantised.
\end{attention}

\courssub{The wave nature of the electron}

The key that unlocks the atom came from Louis de Broglie in 1924. He proposed that every particle of momentum $p$ has an associated \cle{wavelength}:
\[
\lambda = \frac{h}{p} = \frac{h}{mv}
\]
where $h = 6.63\times 10^{-34}\ \mathrm{J\,s}$ is Planck's constant. This \textbf{wave nature of matter} was later confirmed experimentally: electrons fired at a crystal produce \emph{diffraction patterns}, exactly as waves do.

Now imagine the electron in a hydrogen atom not as a particle on an orbit, but as a \textbf{wave confined around the nucleus}. A wave confined in a region of space behaves like the standing wave on a guitar string: only certain wavelengths ``fit''. For a wave wrapped around a circle of radius $r$, the circumference must contain a whole number of wavelengths, otherwise the wave would cancel itself out after many turns:
\[
2\pi r = n\lambda, \qquad n = 1, 2, 3, \dots
\]

\begin{exemplebox}[De Broglie's wave explains Bohr's quantisation]
Substituting $\lambda = h/(mv)$ into $2\pi r = n\lambda$:
\[
2\pi r = n\frac{h}{mv} \quad\Longrightarrow\quad mvr = n\frac{h}{2\pi}
\]
This is \emph{exactly} Bohr's quantisation of angular momentum! What Bohr imposed by hand, de Broglie \emph{derived} from the condition that the electron wave must close on itself smoothly. Only certain orbits -- certain energies -- are allowed, because only certain standing waves fit. Quantisation is no longer a mysterious rule: it is a natural consequence of the wave nature of the electron.
\end{exemplebox}

\begin{popfigure}
\begin{center}
\begin{tikzpicture}[scale=1.0]
% Standing wave that fits (n=4)
\begin{scope}
\draw[popmuted, dashed] (0,0) circle (1.4);
\draw[popblue, thick, domain=0:360, samples=200] plot ({(1.4+0.22*sin(4*\x))*cos(\x)}, {(1.4+0.22*sin(4*\x))*sin(\x)});
\fill[poporange] (0,0) circle (0.09);
\node[below] at (0,-1.9) {\small Wave fits: $2\pi r = 4\lambda$ (allowed)};
\end{scope}
% Standing wave that does not fit
\begin{scope}[xshift=6.2cm]
\draw[popmuted, dashed] (0,0) circle (1.4);
\draw[popink, thick, domain=0:360, samples=200] plot ({(1.4+0.22*sin(3.5*\x))*cos(\x)}, {(1.4+0.22*sin(3.5*\x))*sin(\x)});
\fill[poporange] (0,0) circle (0.09);
\node[below] at (0,-1.9) {\small Wave does not close (forbidden)};
\end{scope}
\end{tikzpicture}
\end{center}
\legende{The electron as a standing wave around the nucleus. Left: the wave closes smoothly on itself after a whole number of wavelengths -- this state is allowed. Right: the wave does not match up -- this state cannot exist.}
\end{popfigure}

\courssub{The Schrödinger equation}

In 1926, Erwin Schrödinger wrote down the equation that governs matter waves. Just as Newton's second law $\vec{F} = m\vec{a}$ governs the motion of classical particles, the \cle{Schrödinger equation} governs the behaviour of quantum particles:
\[
-\frac{\hbar^2}{2m}\nabla^2\psi + V\psi = E\psi
\]
Do not be frightened by the symbols -- \textbf{you are not expected to use or solve this equation at GCE level}. What matters is what it \emph{means}:

\begin{itemize}
\item It is a \textbf{wave equation}: it describes how the electron's matter wave behaves in the electric field of the nucleus (the term $V$ is the electric potential energy of the electron in the field of the proton).
\item Its solutions are functions $\psi$ called \cle{wavefunctions}. Each acceptable solution describes one \textbf{allowed state} of the electron.
\item The requirement that the solutions be physically acceptable (smooth, single-valued, finite) automatically \textbf{quantises the energy}: only certain energies $E$ give acceptable wavefunctions.
\end{itemize}

\begin{attention}[A frequent exam trap]
Candidates sometimes try to write down or manipulate the Schrödinger equation mathematically. At GCE Advanced Level this is \textbf{not required}. You must know: (1) it is a wave equation for the electron, (2) its solutions are wavefunctions $\psi$, (3) each $\psi$ describes an allowed state, and (4) $|\psi|^2$ gives a probability density. That is all.
\end{attention}

\courssub{Interpreting the wavefunction: probability, not orbits}

What does the wavefunction $\psi$ actually represent? Here quantum mechanics makes its most revolutionary statement. The wavefunction itself is not directly observable, but its \textbf{modulus squared} has a precise physical meaning:

\begin{propriete}[Born's interpretation of the wavefunction]
The quantity $|\psi|^2$ at a point is proportional to the \textbf{probability per unit volume} (probability density) of finding the electron at that point. Where $|\psi|^2$ is large, the electron is likely to be found; where $|\psi|^2 = 0$, the electron is never found.
\end{propriete}

This destroys the picture of the electron as a small ball tracing a definite path. The electron is spread out as a \cle{probability cloud} around the nucleus. The region of space where the probability of finding the electron is high is called an \cle{orbital}.

\begin{definition}[Orbital]
An \textbf{orbital} is a region of space around the nucleus where the probability of finding the electron is high. Each orbital is described by a wavefunction $\psi$, solution of the Schrödinger equation, and is characterised by a set of quantum numbers.
\end{definition}

\begin{attention}[The orbital is not an orbit!]
The most common mistake of all: picturing the electron as a small ball moving \emph{along} the orbital, like a car on a road. This is wrong. The orbital is a \textbf{probability distribution} -- a cloud showing where the electron is \emph{likely to be found}. The electron has no definite trajectory. Note also the spelling: an \emph{orbit} (Bohr) is a path; an \emph{orbital} (Schrödinger) is a probability cloud.
\end{attention}

\begin{popfigure}
\begin{center}
\begin{tikzpicture}[scale=1.0]
% Bohr orbit
\begin{scope}
\fill[poporange] (0,0) circle (0.1);
\node[below=2pt] at (0,-0.1) {\small nucleus};
\draw[popblue, thick] (0,0) circle (1.5);
\fill[popblue] (1.5,0) circle (0.09);
\node[popblue, right] at (1.55,0) {\small $e^-$};
\draw[->, popblue, thick] (115:1.5) arc (115:75:1.5);
\node[below] at (0,-1.9) {\small \textbf{Bohr}: electron on a fixed orbit};
\end{scope}
% Schrodinger cloud
\begin{scope}[xshift=7cm]
\shade[inner color=popblue, outer color=white, opacity=0.75] (0,0) circle (1.6);
\fill[poporange] (0,0) circle (0.1);
\node[below=2pt] at (0,-0.1) {\small nucleus};
\node[below] at (0,-1.9) {\small \textbf{Schrödinger}: probability cloud};
\end{scope}
\end{tikzpicture}
\end{center}
\legende{Two visions of the hydrogen ground state. Left: Bohr's electron follows a definite circular orbit of radius $a_0$. Right: Schrödinger's electron is a probability cloud, densest near the centre, with no definite path.}
\end{popfigure}

\courssub{Quantum numbers of the hydrogen atom}

When the Schrödinger equation is solved for the hydrogen atom, each allowed wavefunction is labelled by \textbf{three quantum numbers}. You must know clearly what each one describes.

\begin{center}
\begin{tabularx}{\linewidth}{|l|l|X|}
\hline
\rowcolor{popblueL} \textbf{Quantum number} & \textbf{Symbol and values} & \textbf{What it describes} \\
\hline
Principal & $n = 1, 2, 3, \dots$ & The \textbf{energy} of the state and the average \textbf{size} of the orbital. Larger $n$ means higher energy and a larger cloud. \\
\hline
Orbital (azimuthal) & $l = 0, 1, \dots, n-1$ & The \textbf{shape} of the orbital. $l = 0$ gives a spherical \emph{s} orbital; $l = 1$ gives a dumb-bell \emph{p} orbital. \\
\hline
Magnetic & $m_l = -l, \dots, 0, \dots, +l$ & The \textbf{orientation} of the orbital in space. \\
\hline
\end{tabularx}
\end{center}

Note how the quantum numbers are linked: for a given $n$, only $n$ values of $l$ are possible, and for each $l$ there are $2l+1$ values of $m_l$. (A fourth number, the \emph{spin}, describes the electron's own rotation; it is beyond the GCE syllabus.)

For the \textbf{ground state} of hydrogen, the lowest possible energy state:
\[
n = 1, \qquad l = 0, \qquad m_l = 0
\]
This is the \emph{1s orbital}: a spherical probability cloud centred on the nucleus.

\begin{methode}[Answering exam questions on quantum numbers]
\begin{enumerate}
\item Identify the state asked about (ground state: $n=1$; first excited state: $n=2$; etc.).
\item State each quantum number and \textbf{what it describes}: $n$ gives the energy and size; $l$ gives the shape of the orbital; $m_l$ gives its orientation in space.
\item Give the allowed ranges: $l$ runs from $0$ to $n-1$; $m_l$ runs from $-l$ to $+l$.
\item Name the orbital: $l=0$ is an \emph{s} orbital (spherical); $l=1$ is a \emph{p} orbital (dumb-bell).
\end{enumerate}
\end{methode}

\courssub{Orbitals: shapes and sizes}

The shape of the probability cloud depends on the orbital quantum number $l$:

\begin{itemize}
\item \textbf{s orbitals} ($l = 0$) are \textbf{spherical}: the probability density depends only on the distance from the nucleus, not on the direction.
\item \textbf{p orbitals} ($l = 1$) are \textbf{dumb-bell shaped}, with two lobes on opposite sides of the nucleus and a \emph{node} (a plane where the probability is zero) through the nucleus.
\end{itemize}

\begin{popfigure}
\begin{center}
\begin{tikzpicture}[scale=1.0]
% 1s orbital
\begin{scope}
\shade[inner color=popblue, outer color=white, opacity=0.8] (0,0) circle (1.3);
\fill[poporange] (0,0) circle (0.09);
\draw[popmuted, dashed] (0,0) circle (1.3);
\node[below] at (0,-1.7) {\small 1s orbital: spherical cloud};
\end{scope}
% p orbital
\begin{scope}[xshift=7cm]
\shade[left color=popgreen, right color=white, opacity=0.85] (-1.5,0) ellipse (1.0 and 0.55);
\shade[left color=white, right color=popgreen, opacity=0.85] (1.5,0) ellipse (1.0 and 0.55);
\fill[poporange] (0,0) circle (0.09);
\draw[popmuted, dashed] (0,-0.9) -- (0,0.9);
\node[below] at (0,-1.7) {\small p orbital: dumb-bell shape};
\end{scope}
\end{tikzpicture}
\end{center}
\legende{Shapes of orbitals. The 1s orbital is a spherical probability cloud; a p orbital has two lobes separated by a nodal plane (dashed) where the probability of finding the electron is zero.}
\end{popfigure}

A natural question is: \emph{at what distance from the nucleus is the electron most likely to be found?} For the ground state, we must consider the \textbf{radial probability} $P(r) = 4\pi r^2 |\psi|^2$, which combines the probability density with the volume of a spherical shell of radius $r$. This function is zero both at the nucleus and very far away, and reaches a \textbf{maximum} at exactly the \cle{Bohr radius}:
\[
a_0 = 5.3\times 10^{-11}\ \mathrm{m}
\]

\begin{popfigure}
\begin{center}
\begin{tikzpicture}
\begin{axis}[
    width=11cm, height=6.5cm,
    xlabel={$r\;(\times 10^{-10}\ \mathrm{m})$},
    ylabel={radial probability $P(r)$},
    xmin=0, xmax=5, ymin=0, ymax=1.15,
    axis lines=left,
    xtick={0,0.53,1,2,3,4,5},
    xticklabels={$0$,$a_0$,$1$,$2$,$3$,$4$,$5$},
    ytick=\empty,
    samples=200, domain=0:5,
]
\addplot[popblue, very thick] {x^2*exp(-2*x/0.53)/0.038};
\addplot[popmuted, dashed] coordinates {(0.53,0) (0.53,1.0)};
\node[popblue, anchor=south west] at (axis cs:0.62,1.0) {\small maximum at $r = a_0$};
\end{axis}
\end{tikzpicture}
\end{center}
\legende{Radial probability distribution for the hydrogen ground state (vertical scale arbitrary). The most probable distance of the electron from the nucleus is exactly the Bohr radius $a_0 = 5.3\times 10^{-11}\ \mathrm{m}$ -- a beautiful link between the two models.}
\end{popfigure}

This is a lovely result: the Schrödinger model does not place the electron \emph{on} the Bohr orbit, but the \textbf{most probable distance} coincides with Bohr's radius. The old model captured part of the truth, but the new model tells the fuller story.

\courssub{Successes of the Schrödinger model}

\begin{propriete}[Energy levels from the Schrödinger model]
Solving the Schrödinger equation for hydrogen gives exactly the same energy levels as Bohr's model:
\[
E_n = -\frac{13.6}{n^2}\ \mathrm{eV}, \qquad n = 1, 2, 3, \dots
\]
but the model explains far more: the fine structure of spectra, the intensities of spectral lines (through the probabilities of transitions), and it extends successfully to \textbf{all atoms and molecules}, not just hydrogen.
\end{propriete}

The Schrödinger model is therefore the foundation of the whole of modern atomic physics and chemistry: the arrangement of electrons in orbitals explains the periodic table, chemical bonding, and the behaviour of materials.

\begin{exoresolu}[Ground state of hydrogen in the Schrödinger model]
\textbf{Statement.} (a) Give the quantum numbers of the electron in the ground state of hydrogen and state what each describes. (b) Describe the shape of the corresponding orbital. (c) At what distance from the nucleus is the electron most likely to be found? (d) Calculate the energy of this state in eV and in joules.
\tcblower
\textbf{Solution.}

(a) Ground state: $n = 1$, $l = 0$, $m_l = 0$. Here $n$ fixes the energy and the size of the orbital; $l = 0$ means the orbital is spherical (an \emph{s} orbital); $m_l = 0$ is the only possible orientation.

(b) The orbital is the \emph{1s orbital}: a spherical probability cloud centred on the nucleus. The electron is not on a path -- $|\psi|^2$ tells us only the probability of finding it at each point.

(c) The radial probability is maximum at the Bohr radius:
\[
r = a_0 = 5.3\times 10^{-11}\ \mathrm{m}
\]

(d) Using $E_n = -\dfrac{13.6}{n^2}\ \mathrm{eV}$ with $n = 1$:
\[
E_1 = -13.6\ \mathrm{eV}
\]
Converting to joules with $1\ \mathrm{eV} = 1.60\times 10^{-19}\ \mathrm{J}$:
\[
E_1 = -13.6 \times 1.60\times 10^{-19} = -2.18\times 10^{-18}\ \mathrm{J}
\]
The negative sign means the electron is \textbf{bound} to the nucleus: $13.6\ \mathrm{eV}$ must be supplied to free it.
\end{exoresolu}

\courssub{Heisenberg's uncertainty idea}

Why must we content ourselves with probabilities? Werner Heisenberg gave the deep answer in 1927 with his \cle{uncertainty principle}: \textbf{the position and the momentum of a particle cannot both be known exactly at the same time}. The more precisely we know where the electron is, the less precisely we can know how it is moving, and vice versa.

This is not a failure of our instruments -- it is a fundamental property of nature. Because we can never know both the position and the velocity of the electron exactly, the very idea of a definite \emph{orbit} loses its meaning. The best any theory can do is to give \textbf{probabilities} -- and that is precisely what $|\psi|^2$ provides. The probability cloud is not a sign of ignorance that better instruments could remove; it is the complete description nature allows.

\begin{savaistu}[Quantum mechanics all around us]
The Schrödinger equation is not an abstract curiosity. It explains why atoms emit and absorb the exact colours we use to identify elements in stars; it underlies the laser, the transistor in every phone, and the LED lamps lighting homes across Cameroon; and it explains the chemical bonds that hold together every molecule in your body. Whenever you switch on a phone, you are using quantum mechanics.
\end{savaistu}

\begin{aretenir}[Key points of this section]
\begin{itemize}
\item The Bohr model is limited: it works only for hydrogen, ignores fine structure and line intensities, and treats the electron as a particle on a fixed orbit.
\item De Broglie: the electron is a wave of wavelength $\lambda = h/(mv)$; only standing waves that fit around the nucleus survive, which explains quantisation.
\item The Schrödinger equation is a wave equation whose solutions, the wavefunctions $\psi$, describe the allowed states of the electron (interpretation only at GCE level).
\item $|\psi|^2$ gives the probability density of finding the electron; the electron forms a probability cloud called an \textbf{orbital}, not a definite orbit.
\item Three quantum numbers: $n$ (energy and size), $l$ (shape: $l=0$ spherical s, $l=1$ dumb-bell p), $m_l$ (orientation). Ground state: $n=1$, $l=0$, $m_l=0$.
\item The most probable electron distance in the ground state is the Bohr radius $a_0 = 5.3\times 10^{-11}\ \mathrm{m}$.
\item The model reproduces $E_n = -13.6/n^2\ \mathrm{eV}$ and extends to all atoms and molecules.
\item Heisenberg: position and momentum cannot both be known exactly, so only probabilities can be given.
\end{itemize}
\end{aretenir}

\begin{exercice}[Checking your understanding]
\begin{enumerate}
\item State two limitations of the Bohr model that the Schrödinger model overcomes.
\item An electron in hydrogen has quantum numbers $n = 2$, $l = 1$, $m_l = -1$. What does each quantum number tell you, and what is the shape of the orbital?
\item Explain why the concept of a definite electron orbit had to be abandoned, referring to Heisenberg's uncertainty principle.
\item Calculate the energy of the $n = 3$ level of hydrogen in eV, and the wavelength of the photon emitted when the electron falls from $n = 3$ to the ground state.
\end{enumerate}
\end{exercice}

\begin{corrige}[Exercise 1]
\textbf{1.} Any two of: the Bohr model works only for one-electron atoms while the Schrödinger model extends to all atoms and molecules; Bohr cannot explain the fine details (splitting) of spectral lines; Bohr cannot explain why some lines are more intense than others; Bohr treats the electron as a classical particle on a fixed orbit, whereas the electron is really described by a probability cloud.

\textbf{2.} $n = 2$: the electron is in the second energy level (first excited level), with a larger orbital than the ground state. $l = 1$: the orbital is a \emph{p} orbital, dumb-bell shaped. $m_l = -1$: this specifies the orientation of the dumb-bell in space.

\textbf{3.} Heisenberg's uncertainty principle states that the position and momentum of the electron cannot both be known exactly at the same time. A definite orbit would require knowing both at every instant, which nature forbids. Hence the electron can only be described by the probability distribution $|\psi|^2$ -- the orbital -- and not by a trajectory.

\textbf{4.} Energy of the $n = 3$ level:
\[
E_3 = -\frac{13.6}{3^2} = -1.51\ \mathrm{eV}
\]
The photon emitted in the transition $n = 3 \to n = 1$ has energy:
\[
hf = E_3 - E_1 = (-1.51) - (-13.6) = 12.1\ \mathrm{eV} = 12.1 \times 1.60\times 10^{-19} = 1.94\times 10^{-18}\ \mathrm{J}
\]
Then the wavelength:
\[
\lambda = \frac{hc}{hf} = \frac{6.63\times 10^{-34} \times 3.00\times 10^{8}}{1.94\times 10^{-18}} = 1.03\times 10^{-7}\ \mathrm{m} = 103\ \mathrm{nm}
\]
This is an ultraviolet photon (a member of the Lyman series), exactly as observed in the hydrogen spectrum -- a quantitative triumph of the model.
\end{corrige}

\newpage

\coursec{Integration activity}

\courssub{Aims of the Activity}
\begin{aretenir}[Aims of the Activity]
At the end of this activity, you should be able to:
\begin{itemize}
\item use the relation $E = hf = \dfrac{hc}{\lambda}$ to convert a measured wavelength into a photon energy, expressed both in joules and in electron-volts;
\item use the quantised energy levels of hydrogen, $E_n = -\dfrac{13.6}{n^2}\ \mathrm{eV}$, to compute the energy of a transition $\Delta E = E_{\text{final}} - E_{\text{initial}}$ between two levels;
\item identify the spectral lines of the Balmer series and associate each measured wavelength with a well-defined electronic transition;
\item represent transitions by arrows on an energy level diagram;
\item argue, from experimental evidence, that the light emitted by a gas discharge tube carries the ``signature'' of the gas it contains, and explain how astronomers exploit this idea to determine the chemical composition of stars.
\end{itemize}
\end{aretenir}

\courssub{Situation}
\begin{experience}[Identifying an Unknown Gas by its Spectrum]
During the practical week at a Government High School in Yaound\'e, the physics club receives a visit from a technician of the Cameroon Radio Television science programme. He brings a \cle{gas discharge tube} connected to a high-voltage supply: a sealed glass tube containing a pure gas at low pressure, between two electrodes. When the high voltage is switched on, the gas glows with a beautiful pink-violet light. The label on the tube has been erased, and the technician challenges the students: \emph{``Identify the gas in this tube using only the light it emits.''}

The students observe the light through a \cle{diffraction grating} (600 lines per millimetre). Instead of a continuous rainbow, they see a few sharp, bright coloured lines at precise angles: an \cle{emission line spectrum}. Using the grating relation, they measure the wavelengths of the four visible lines:

\begin{center}
\begin{tabularx}{\linewidth}{|l|X|X|X|X|}
\hline
\rowcolor{popblueL} Line & 1 (red) & 2 (blue-green) & 3 (violet) & 4 (deep violet) \\
\hline
Wavelength $\lambda$ (nm) & 656 & 486 & 434 & 410 \\
\hline
\end{tabularx}
\end{center}

\textbf{Data provided to the groups:}
\begin{itemize}
\item Planck constant: $h = 6.63\times 10^{-34}\ \mathrm{J\,s}$; speed of light in vacuum: $c = 3.00\times 10^{8}\ \mathrm{m\,s^{-1}}$;
\item elementary charge: $e = 1.60\times 10^{-19}\ \mathrm{C}$; $1\ \mathrm{eV} = 1.60\times 10^{-19}\ \mathrm{J}$;
\item energy levels of the hydrogen atom: $E_n = -\dfrac{13.6}{n^2}\ \mathrm{eV}$, with $n = 1, 2, 3, \dots$
\item a sodium vapour lamp is known to emit a dominant \cle{yellow doublet}: two very close lines near $589\ \mathrm{nm}$ and $589.6\ \mathrm{nm}$.
\end{itemize}

The club suspects the gas might be hydrogen, since the four wavelengths look like the famous \cle{Balmer series}. Your mission, working in groups, is to prove (or disprove) this hypothesis quantitatively, and to present your conclusion to the technician.
\end{experience}

\begin{definition}[Gas Discharge Tube]
A sealed glass tube containing a gas at low pressure (a few torr) with electrodes at each end. When a high voltage (several kV) is applied, the gas ionises and conducts electricity, exciting atoms which then de-excite by emitting photons characteristic of the gas.
\end{definition}

\begin{definition}[Diffraction Grating]
An optical component with a periodic structure (e.g., 600 lines/mm) that splits light into its component wavelengths by diffraction. The grating equation is $d\sin\theta = m\lambda$ where $d$ is the slit spacing.
\end{definition}

\begin{definition}[Emission Line Spectrum]
A spectrum consisting of bright lines at specific wavelengths on a dark background, produced when excited atoms in a low-pressure gas return to lower energy states. Each element has a unique line spectrum.
\end{definition}

\begin{definition}[Balmer Series]
The series of spectral lines emitted by hydrogen when the electron falls to the $n=2$ level from higher levels ($n=3,4,5,\dots$). The first four lines (H$\alpha$, H$\beta$, H$\gamma$, H$\delta$) lie in the visible region.
\end{definition}

\courssub{Tasks}
\begin{exercice}[Spectral Analysis of the Unknown Gas]
\begin{enumerate}
\item \textbf{From wavelength to photon energy.} For each of the four measured lines, calculate the photon energy using $E = hf = \dfrac{hc}{\lambda}$. Give each result in joules, then convert it into electron-volts. Present your four results in a table. (Remember to convert nanometres into metres before substituting.)

\item \textbf{Energy levels of hydrogen.} Using $E_n = -\dfrac{13.6}{n^2}\ \mathrm{eV}$, calculate the energies $E_1, E_2, E_3, E_4, E_5$ and $E_6$ of the first six levels of the hydrogen atom. Give your answers in eV, to two decimal places.

\item \textbf{Identifying the transitions.} The Balmer series corresponds to transitions in which the electron falls \emph{down to the level $n = 2$} from a higher level $n = 3, 4, 5, 6, \dots$ Using $\Delta E = E_{n} - E_{2}$, compute the energies of the transitions $3 \rightarrow 2$, $4 \rightarrow 2$, $5 \rightarrow 2$ and $6 \rightarrow 2$. Compare them with the photon energies found in question 1 and match each measured line to a transition.

\item \textbf{Energy level diagram.} Draw (to a reasonable vertical scale) the energy level diagram of hydrogen for $n = 1$ to $n = 6$, including the ionisation limit $E = 0$. On this diagram, draw a downward arrow for each of the four transitions identified, and label each arrow with the wavelength of the emitted photon.

\item \textbf{Conclusion on the identity of the gas.} Write a short paragraph (5--8 lines) stating your conclusion. Is the gas hydrogen? Justify your answer by referring to your numerical results and to the idea of quantised energy levels.

\item \textbf{Why not sodium?} Explain, in a few sentences, what the students would have observed if the tube had contained sodium vapour instead, and why the observation of a dominant yellow doublet near $589\ \mathrm{nm}$ would rule out hydrogen.

\item \textbf{Extension: light from the stars.} Astronomers spread the light of a star with a grating (or a prism) and observe dark absorption lines superimposed on a continuous spectrum. Using the ideas developed in this activity, explain how they can deduce which chemical elements are present in the star's atmosphere, even though the star may be millions of kilometres... in fact light-years away.
\end{enumerate}
\end{exercice}

\begin{methode}[Converting Wavelength to Photon Energy]
\begin{enumerate}
\item Convert wavelength $\lambda$ from nanometres to metres: $\lambda(\mathrm{m}) = \lambda(\mathrm{nm}) \times 10^{-9}$.
\item Compute energy in joules: $E(\mathrm{J}) = \dfrac{hc}{\lambda}$ with $h = 6.63\times 10^{-34}\ \mathrm{J\,s}$, $c = 3.00\times 10^8\ \mathrm{m\,s^{-1}}$.
\item Convert to electron-volts: $E(\mathrm{eV}) = \dfrac{E(\mathrm{J})}{e}$ with $e = 1.60\times 10^{-19}\ \mathrm{C}$ (since $1\ \mathrm{eV} = 1.60\times 10^{-19}\ \mathrm{J}$).
\item Keep three significant figures throughout.
\end{enumerate}
\end{methode}

\begin{methode}[Identifying Transitions from Energy Levels]
\begin{enumerate}
\item List the energy levels $E_n$ (in eV) for the relevant principal quantum numbers $n$.
\item For a transition from initial level $n_i$ to final level $n_f$ ($n_i > n_f$), the photon energy is $\Delta E = E_{n_f} - E_{n_i}$ (a positive quantity).
\item Compare the calculated $\Delta E$ values with the experimental photon energies.
\item Match each experimental line to the transition giving the closest energy match.
\end{enumerate}
\end{methode}

\begin{attention}[Common Mistakes]
\begin{itemize}
\item Forgetting to convert nm to m before using $hc/\lambda$ (leads to energies $10^9$ times too large).
\item Using $E = hf$ with $f = c/\lambda$ but mixing units (e.g., $\lambda$ in nm, $c$ in m/s).
\item Confusing the sign of energy levels: $E_n$ are negative (bound states). The photon energy is $\Delta E = E_{\text{final}} - E_{\text{initial}} = E_{\text{lower}} - E_{\text{upper}} > 0$.
\item Rounding intermediate results too early; keep extra digits until the final answer.
\end{itemize}
\end{attention}

\courssub{Model Answer}
\begin{exoresolu}[Model Answer]
\textbf{1. Photon energies.}
For line 1, $\lambda_1 = 656\ \mathrm{nm} = 6.56\times 10^{-7}\ \mathrm{m}$:
\begin{align*}
E_1 &= \frac{hc}{\lambda_1} = \frac{(6.63\times 10^{-34}\ \mathrm{J\,s})(3.00\times 10^{8}\ \mathrm{m\,s^{-1}})}{6.56\times 10^{-7}\ \mathrm{m}} \\
&= 3.03\times 10^{-19}\ \mathrm{J} = \frac{3.03\times 10^{-19}\ \mathrm{J}}{1.60\times 10^{-19}\ \mathrm{J\,eV^{-1}}} = 1.89\ \mathrm{eV}.
\end{align*}
Repeating the same calculation for the other lines:

\begin{center}
\begin{tabularx}{\linewidth}{|l|X|X|X|X|}
\hline
\rowcolor{popblueL} Line & 1 & 2 & 3 & 4 \\
\hline
$\lambda$ (nm) & 656 & 486 & 434 & 410 \\
\hline
$E$ ($10^{-19}$ J) & 3.03 & 4.09 & 4.58 & 4.85 \\
\hline
$E$ (eV) & 1.89 & 2.56 & 2.86 & 3.03 \\
\hline
\end{tabularx}
\end{center}

\textbf{2. Energy levels.}
Applying $E_n = -\dfrac{13.6}{n^2}\ \mathrm{eV}$:
\begin{align*}
E_1 &= -13.60\ \mathrm{eV}, & E_2 &= -3.40\ \mathrm{eV}, & E_3 &= -1.51\ \mathrm{eV},\\
E_4 &= -0.85\ \mathrm{eV}, & E_5 &= -0.54\ \mathrm{eV}, & E_6 &= -0.38\ \mathrm{eV}.
\end{align*}

\textbf{3. Transition energies.}
For a transition from level $n$ down to level 2, the emitted photon carries away $\Delta E = E_n - E_2$ (since $E_n > E_2$):
\begin{align*}
3\rightarrow 2 &: \quad -1.51 - (-3.40) = 1.89\ \mathrm{eV} \quad \Rightarrow \text{line 1 (656 nm)},\\
4\rightarrow 2 &: \quad -0.85 - (-3.40) = 2.55\ \mathrm{eV} \quad \Rightarrow \text{line 2 (486 nm)},\\
5\rightarrow 2 &: \quad -0.54 - (-3.40) = 2.86\ \mathrm{eV} \quad \Rightarrow \text{line 3 (434 nm)},\\
6\rightarrow 2 &: \quad -0.38 - (-3.40) = 3.02\ \mathrm{eV} \quad \Rightarrow \text{line 4 (410 nm)}.
\end{align*}
Each measured photon energy coincides, within rounding errors, with one Balmer transition. The agreement is excellent.

\textbf{4. Energy level diagram.}

\begin{popfigure}
\begin{tikzpicture}[scale=0.8]
% Energy axis
\draw[->, thick] (0,0) -- (0,15) node[above] {$E$ (eV)};
\draw[thick] (-0.2,0) -- (0.2,0) node[right] {$0$ (ionisation)};
% Levels: y = -E_n (so E=0 at top, negative down)
% n=1: -13.6 -> y=13.6
% n=2: -3.40 -> y=3.40
% n=3: -1.51 -> y=1.51
% n=4: -0.85 -> y=0.85
% n=5: -0.54 -> y=0.54
% n=6: -0.38 -> y=0.38
% Scale factor for visibility: multiply y by 1.0 (1 eV = 1 cm approx)
\def\scale{1.0}
% n=1
\draw[thick, popink] (-2,13.6*\scale) -- (2,13.6*\scale) node[right] {$n=1,\ -13.60\ \mathrm{eV}$};
% n=2
\draw[thick, popink] (-2,3.40*\scale) -- (2,3.40*\scale) node[right] {$n=2,\ -3.40\ \mathrm{eV}$};
% n=3
\draw[thick, popink] (-2,1.51*\scale) -- (2,1.51*\scale) node[right] {$n=3,\ -1.51\ \mathrm{eV}$};
% n=4
\draw[thick, popink] (-2,0.85*\scale) -- (2,0.85*\scale) node[right] {$n=4,\ -0.85\ \mathrm{eV}$};
% n=5
\draw[thick, popink] (-2,0.54*\scale) -- (2,0.54*\scale) node[right] {$n=5,\ -0.54\ \mathrm{eV}$};
% n=6
\draw[thick, popink] (-2,0.38*\scale) -- (2,0.38*\scale) node[right] {$n=6,\ -0.38\ \mathrm{eV}$};
% Ionisation limit
\draw[thick, dashed, popmuted] (-2,0) -- (2,0) node[right] {$E=0$};
% Transitions (arrows downward)
\draw[->, very thick, poporange] (2.5,1.51*\scale) -- (2.5,3.40*\scale) node[midway, right] {656 nm};
\draw[->, very thick, popgreen] (3.5,0.85*\scale) -- (3.5,3.40*\scale) node[midway, right] {486 nm};
\draw[->, very thick, popblue] (4.5,0.54*\scale) -- (4.5,3.40*\scale) node[midway, right] {434 nm};
\draw[->, very thick, poppurple] (5.5,0.38*\scale) -- (5.5,3.40*\scale) node[midway, right] {410 nm};
% Labels for series
\node[align=center, popdark] at (4,14.5) {Balmer series:\\transitions to $n=2$};
\end{tikzpicture}
\legende{Energy level diagram of hydrogen (vertical scale: 1 eV = 1 cm). Downward arrows indicate photon emission for the Balmer series.}
\end{popfigure}

\textbf{5. Conclusion.}
The four measured wavelengths correspond exactly (to within experimental and rounding errors) to photon energies equal to the differences $E_n - E_2$ between the quantised levels of hydrogen, for $n = 3, 4, 5, 6$. No other element possesses this particular set of level spacings, because each atom has its own unique energy level scheme determined by its nuclear charge and electron configuration. Since the discharge tube emits precisely the Balmer lines of hydrogen and no other prominent lines, we conclude that the unknown gas \textbf{is hydrogen}. The discrete line spectrum is direct experimental evidence that the energy of the electron in the atom is quantised: only photons of well-defined energies $hf = E_{\text{final}} - E_{\text{initial}}$ can be emitted.

\textbf{6. Why not sodium?}
If the tube contained sodium vapour, the students would have observed a spectrum dominated by an intense \cle{yellow doublet}: two very close lines at about $589.0\ \mathrm{nm}$ and $589.6\ \mathrm{nm}$, corresponding to the $3p \to 3s$ transitions in sodium (spin-orbit splitting). The measured wavelengths 656, 486, 434 and 410 nm do not match the sodium doublet, and no strong yellow line is seen. The spectrum is the ``identity card'' of the gas: observing the sodium doublet would rule out hydrogen, just as observing the Balmer series rules out sodium.

\textbf{7. Light from the stars.}
A star's hot, dense interior emits a continuous (black-body) spectrum. As this light crosses the cooler, low-density outer atmosphere (photosphere/chromosphere) of the star, atoms there absorb photons whose energies match exactly the differences between their own quantised levels, $hf = E_{\text{upper}} - E_{\text{lower}}$. The absorbed wavelengths are missing from the transmitted light, appearing as dark \cle{absorption lines} (Fraunhofer lines) superimposed on the continuous spectrum. By measuring the positions of these dark lines and comparing them with the emission/absorption spectra of elements measured in the laboratory (exactly as we compared our four lines with the hydrogen levels), astronomers identify the elements present in the star's atmosphere. In this way hydrogen, helium, sodium, iron and many other elements have been detected in the Sun and in distant stars, without any sample ever leaving the star. This is the foundation of \cle{astrophysical spectroscopy}.
\end{exoresolu}

\begin{savaistu}[The Schr\"odinger Model and the Hydrogen Atom]
The formula $E_n = -13.6/n^2\ \mathrm{eV}$ was first derived by Bohr (1913) using a semi-classical model. The full quantum-mechanical treatment by Schr\"odinger (1926) solves the wave equation for the electron in the Coulomb potential of the proton. It yields the same energy levels (depending only on $n$) but also predicts the orbital shapes (s, p, d, f) and the fine structure (small splittings due to relativistic effects and spin-orbit coupling) which are not captured by the simple Bohr formula. For the GCE Advanced Level, the Bohr formula is used for hydrogen-like atoms.
\end{savaistu}

\begin{savaistu}[Spectroscopy in Cameroon]
The University of Yaound\'e I and the Institute of Geological and Mining Research (IRGM) use spectroscopic techniques (X-ray fluorescence, atomic absorption) to analyse mineral ores and environmental samples. The principle is exactly the same: each element emits or absorbs light at characteristic wavelengths, allowing quantitative analysis of the composition of rocks, soils, and water.
\end{savaistu}

\newpage

\coursec{Methods and worked examples}

\courssub{Method 1: Interpreting Rutherford's alpha-scattering experiment}

\begin{methode}[Analysing alpha-particle scattering]
When a question presents scattering data or asks you to deduce properties of the nucleus:
\begin{enumerate}
\item Recall the set-up: a thin gold foil (a few hundred atoms thick) is bombarded by $\alpha$ particles (\ensuremath{\mathrm{^{4}_{2}He}} nuclei, charge $+2e$) from a radioactive source; scattered particles are detected by scintillations on a zinc sulphide screen mounted on a rotatable microscope.
\item State the three observations and their conclusions:
\begin{itemize}
\item \textbf{Most $\alpha$ particles pass straight through} (deviation $< 1^{\circ}$) $\Rightarrow$ the atom is mostly \cle{empty space}.
\item \textbf{A few are deflected through large angles} ($> 90^{\circ}$) $\Rightarrow$ the positive charge and nearly all the mass are concentrated in a tiny \cle{nucleus}.
\item \textbf{A very small fraction (about 1 in 8000) is turned back} ($\approx 180^{\circ}$) $\Rightarrow$ the nucleus is very small and dense compared with the atom ($R_{\text{nucleus}} \approx 10^{-15}\ \mathrm{m}$, $R_{\text{atom}} \approx 10^{-10}\ \mathrm{m}$).
\end{itemize}
\item For a \emph{head-on} collision, use conservation of energy: the initial kinetic energy of the $\alpha$ particle is entirely converted into electric potential energy at the distance of closest approach $d$:
\[ E_k = \frac{1}{4\pi\varepsilon_0}\,\frac{(2e)(Ze)}{d} \quad\Longrightarrow\quad d = \frac{2Ze^2}{4\pi\varepsilon_0\, E_k}. \]
\item Always quote distances in SI units (metres) and energies in joules (convert from MeV if needed: $1\ \mathrm{MeV} = 1.60\times10^{-13}\ \mathrm{J}$).
\end{enumerate}
\end{methode}

\begin{popfigure}
\begin{tikzpicture}[scale=1.0]
% Nucleus
\fill[poporange] (0,0) circle (0.18);
\node[below] at (0,-0.25) {\small nucleus $(+Ze)$};
% Straight-through paths
\draw[->, thick, popblue] (-4,1.5) -- (3.5,1.5);
\draw[->, thick, popblue] (-4,-1.5) -- (3.5,-1.5);
\node[popblue, right] at (3.5,1.5) {\small most pass through};
% Slightly deflected
\draw[->, thick, popgreen] (-4,0.7) -- (-0.3,0.62) -- (3.5,1.4);
% Large angle
\draw[->, thick, poppurple] (-4,0.35) -- (-0.25,0.28) -- (2.8,-1.8);
\node[poppurple, right] at (2.8,-1.8) {\small large angle};
% Back-scattered
\draw[->, thick, poppink] (-4,0.08) -- (-0.35,0.05) -- (-2.6,-1.4);
\node[poppink, left] at (-2.7,-1.5) {\small rare: bounced back};
% Incoming label
\node[popdark, left] at (-4,0.8) {\small $\alpha$ particles};
\end{tikzpicture}
\legende{Rutherford's experiment: the fate of alpha particles fired at a thin gold foil. Only those passing very close to the tiny, dense, positive nucleus are strongly deflected.}
\end{popfigure}

\begin{exoresolu}[Closest approach of an alpha particle to a gold nucleus]
An alpha particle of kinetic energy $5.0\ \mathrm{MeV}$ is fired directly at the nucleus of a gold atom ($Z = 79$). Assuming the nucleus remains stationary, calculate the distance of closest approach of the alpha particle to the centre of the nucleus. Take $\dfrac{1}{4\pi\varepsilon_0} = 9.0\times10^{9}\ \mathrm{N\,m^2\,C^{-2}}$, $e = 1.60\times10^{-19}\ \mathrm{C}$. Comment on the result, given that the radius of a gold nucleus is about $7\times10^{-15}\ \mathrm{m}$.
\tcblower
\textbf{Step 1 -- Convert the kinetic energy to joules.}
\[ E_k = 5.0\ \mathrm{MeV} = 5.0\times10^{6}\times 1.60\times10^{-19}\ \mathrm{J} = 8.0\times10^{-13}\ \mathrm{J}. \]

\textbf{Step 2 -- Apply conservation of energy.} Far away, the alpha particle has only kinetic energy. At the distance of closest approach $d$, it is momentarily at rest, so all its kinetic energy has become electric potential energy:
\[ E_k = \frac{1}{4\pi\varepsilon_0}\,\frac{q_\alpha\, Q_{\text{Au}}}{d} = \frac{1}{4\pi\varepsilon_0}\,\frac{(2e)(79e)}{d}. \]

\textbf{Step 3 -- Solve for $d$.}
\[ d = \frac{9.0\times10^{9}\times 2\times79\times(1.60\times10^{-19})^{2}}{8.0\times10^{-13}}. \]
Numerator: $9.0\times10^{9}\times158\times2.56\times10^{-38} = 3.64\times10^{-26}\ \mathrm{J\,m}$.
\[ d = \frac{3.64\times10^{-26}}{8.0\times10^{-13}} = 4.6\times10^{-14}\ \mathrm{m}. \]

\textbf{Step 4 -- Comment.} The distance of closest approach $d \approx 4.6\times10^{-14}\ \mathrm{m}$ is about six times larger than the nuclear radius ($7\times10^{-15}\ \mathrm{m}$): the alpha particle does not touch the nucleus, so treating the nucleus as a point charge is justified. It is also about $10^{4}$ times smaller than the atomic radius ($\sim 10^{-10}\ \mathrm{m}$), confirming that the nucleus occupies a minute fraction of the atom's volume.
\end{exoresolu}

\courssub{Method 2: Converting between joules and electron volts}

\begin{methode}[Using the electron volt]
\begin{enumerate}
\item Memorise the definition: \cle{one electron volt} is the kinetic energy gained by an electron accelerated from rest through a potential difference of one volt:
\[ 1\ \mathrm{eV} = e\times 1\ \mathrm{V} = 1.60\times10^{-19}\ \mathrm{J}. \]
\item To convert eV $\to$ J: \textbf{multiply} by $1.60\times10^{-19}$. To convert J $\to$ eV: \textbf{divide} by $1.60\times10^{-19}$.
\item For a particle of charge $q$ accelerated from rest through a p.d. $V$: $E_k = qV$ (in joules if $q$ is in coulombs; directly in eV if $q$ is counted in units of $e$).
\item Use eV for atomic-scale energies (a few eV); use MeV for nuclear-scale energies.
\end{enumerate}
\end{methode}

\begin{exoresolu}[Electron accelerated through a potential difference]
In an X-ray tube at a hospital in Yaound\'e, electrons are accelerated from rest through a potential difference of $12\ \mathrm{kV}$.
\begin{enumerate}
\item Calculate the kinetic energy of each electron on arrival at the anode, in electron volts and in joules.
\item Deduce the speed of the electron, assuming non-relativistic mechanics ($m_e = 9.11\times10^{-31}\ \mathrm{kg}$).
\end{enumerate}
\tcblower
\textbf{1. Kinetic energy.} The work done by the electric field equals the gain in kinetic energy: $E_k = qV = eV$.
\[ E_k = 1 \times 12\ \mathrm{kV} = 1.2\times10^{4}\ \mathrm{eV} = 12\ \mathrm{keV}. \]
In joules:
\[ E_k = 1.2\times10^{4}\times1.60\times10^{-19} = 1.92\times10^{-15}\ \mathrm{J}. \]

\textbf{2. Speed.} Using $E_k = \tfrac{1}{2}m_e v^2$:
\[ v = \sqrt{\frac{2E_k}{m_e}} = \sqrt{\frac{2\times1.92\times10^{-15}}{9.11\times10^{-31}}} = \sqrt{4.21\times10^{15}} = 6.5\times10^{7}\ \mathrm{m\,s^{-1}}. \]
This is about $0.22\,c$, so the non-relativistic approximation is acceptable (though only just). \textbf{Check:} the speed is below $c = 3.00\times10^{8}\ \mathrm{m\,s^{-1}}$, as it must be.
\end{exoresolu}

\courssub{Method 3: Linking photon energy, frequency and wavelength}

\begin{methode}[Using $E = hf = \dfrac{hc}{\lambda}$]
\begin{enumerate}
\item Write the two key relations: $E = hf$ (Planck--Einstein) and $c = f\lambda$, hence
\[ E = \frac{hc}{\lambda} \qquad\text{with } h = 6.63\times10^{-34}\ \mathrm{J\,s},\ c = 3.00\times10^{8}\ \mathrm{m\,s^{-1}}. \]
\item Make sure $\lambda$ is in \textbf{metres} (convert nm: $1\ \mathrm{nm} = 10^{-9}\ \mathrm{m}$) and $E$ in \textbf{joules}.
\item If the energy is given in eV, convert to joules \emph{before} using $E = hf$.
\item Useful shortcut: $\lambda(\mathrm{nm}) \approx \dfrac{1240}{E(\mathrm{eV})}$.
\item Remember the ordering: short wavelength $\Rightarrow$ high frequency $\Rightarrow$ \emph{high} photon energy.
\end{enumerate}
\end{methode}

\begin{exoresolu}[Photon energy of visible light]
The red line of the hydrogen spectrum has wavelength $\lambda = 656\ \mathrm{nm}$.
\begin{enumerate}
\item Calculate the frequency of this radiation.
\item Calculate the energy of one photon, in joules and in electron volts.
\item A lamp emits $2.0\ \mathrm{W}$ of this red light. How many photons does it emit per second?
\end{enumerate}
\tcblower
\textbf{1. Frequency.}
\[ f = \frac{c}{\lambda} = \frac{3.00\times10^{8}}{656\times10^{-9}} = 4.57\times10^{14}\ \mathrm{Hz}. \]

\textbf{2. Photon energy.}
\[ E = hf = 6.63\times10^{-34}\times4.57\times10^{14} = 3.03\times10^{-19}\ \mathrm{J}. \]
In electron volts:
\[ E = \frac{3.03\times10^{-19}}{1.60\times10^{-19}} = 1.89\ \mathrm{eV}. \]
(Shortcut check: $1240/656 = 1.89\ \mathrm{eV}$. \checkmark)

\textbf{3. Number of photons per second.} Power = energy emitted per second, so
\[ N = \frac{P}{E} = \frac{2.0}{3.03\times10^{-19}} = 6.6\times10^{18}\ \text{photons per second}. \]
The enormous number explains why light usually appears continuous even though it is emitted as discrete quanta.
\end{exoresolu}

\courssub{Method 4: Energy levels and transitions in the hydrogen atom}

\begin{methode}[Applying $E = hf = E_2 - E_1$]
\begin{enumerate}
\item Write the hydrogen level formula: $E_n = -\dfrac{13.6}{n^2}\ \mathrm{eV}$ ($n = 1, 2, 3, \dots$). $E_1 = -13.6\ \mathrm{eV}$ is the \cle{ground state}; $E_\infty = 0$ is the \cle{ionisation} limit.
\item Identify the initial level $n_i$ (higher energy) and final level $n_f$ (lower energy) for \textbf{emission}; the photon carries away the difference:
\[ hf = E_{n_i} - E_{n_f} = 13.6\left(\frac{1}{n_f^{2}} - \frac{1}{n_i^{2}}\right)\ \mathrm{eV}. \]
\item For \textbf{absorption}, the photon energy must match \emph{exactly} a difference $E_{n_f} - E_{n_i}$ with $n_f > n_i$.
\item Compute $\Delta E$ in eV, convert to joules ($\times 1.60\times10^{-19}$), then get $\lambda = hc/\Delta E$.
\item Check the spectral region: Lyman ($n_f=1$, ultraviolet), Balmer ($n_f=2$, visible/near-UV), Paschen ($n_f=3$, infrared).
\end{enumerate}
\end{methode}

\begin{popfigure}
\begin{tikzpicture}[scale=1.0]
% Energy axis
\draw[->, thick] (0,-6.4) -- (0,0.6) node[above] {\small $E$ (eV)};
% Levels
\draw[very thick, popblue] (1,-6) -- (5,-6) node[right] {\small $E_1 = -13.6$};
\draw[very thick, popblue] (1,-4.5) -- (5,-4.5) node[right] {\small $E_2 = -3.40$};
\draw[very thick, popblue] (1,-3.33) -- (5,-3.33) node[right] {\small $E_3 = -1.51$};
\draw[very thick, popblue] (1,-2.75) -- (5,-2.75) node[right] {\small $E_4 = -0.85$};
\draw[very thick, popdark, dashed] (1,-1.2) -- (5,-1.2) node[right] {\small $E_\infty = 0$};
% n labels
\node[left] at (1,-6) {\small $n=1$};
\node[left] at (1,-4.5) {\small $n=2$};
\node[left] at (1,-3.33) {\small $n=3$};
\node[left] at (1,-2.75) {\small $n=4$};
% Transitions
\draw[->, very thick, poppink] (2.2,-3.33) -- (2.2,-4.5);
\node[poppink, right] at (2.3,-3.9) {\small $3\to2$: red, $656\ \mathrm{nm}$};
\draw[->, very thick, poppurple] (3.2,-4.5) -- (3.2,-6);
\node[poppurple, right] at (3.3,-5.3) {\small $2\to1$: UV, $122\ \mathrm{nm}$};
\draw[->, very thick, popgreen] (4.3,-1.2) -- (4.3,-6);
\node[popgreen, right] at (4.4,-3.5) {\small ionisation: $13.6\ \mathrm{eV}$};
\end{tikzpicture}
\legende{Energy-level diagram of the hydrogen atom (not to scale vertically). Downward arrows: photon emission. Ionisation from the ground state requires $13.6\ \mathrm{eV}$.}
\end{popfigure}

\begin{exoresolu}[Balmer transition in hydrogen]
The energy levels of the hydrogen atom are given by $E_n = -\dfrac{13.6}{n^2}\ \mathrm{eV}$.
\begin{enumerate}
\item Calculate the energies of the levels $n = 2$ and $n = 3$.
\item An electron falls from level $n = 3$ to level $n = 2$. Show that the emitted photon has energy $1.89\ \mathrm{eV}$ and calculate its wavelength.
\item In which series and which region of the spectrum does this line appear?
\item What is the minimum energy needed to ionise a hydrogen atom initially in the level $n = 2$?
\end{enumerate}
\tcblower
\textbf{1. Level energies.}
\[ E_2 = -\frac{13.6}{2^2} = -3.40\ \mathrm{eV}, \qquad E_3 = -\frac{13.6}{3^2} = -1.51\ \mathrm{eV}. \]

\textbf{2. Emitted photon.} Applying $E = hf = E_2 - E_1$ in the generalised form $\Delta E = E_{n_i} - E_{n_f}$:
\[ \Delta E = E_3 - E_2 = (-1.51) - (-3.40) = 1.89\ \mathrm{eV}. \]
Convert to joules: $\Delta E = 1.89\times1.60\times10^{-19} = 3.02\times10^{-19}\ \mathrm{J}$. Then
\[ \lambda = \frac{hc}{\Delta E} = \frac{6.63\times10^{-34}\times3.00\times10^{8}}{3.02\times10^{-19}} = 6.59\times10^{-7}\ \mathrm{m} = 659\ \mathrm{nm}. \]
(The accepted value is $656\ \mathrm{nm}$; the small difference comes from rounding $13.6\ \mathrm{eV}$.)

\textbf{3. Series and region.} The transition ends on $n = 2$: it belongs to the \textbf{Balmer series}. With $\lambda \approx 656\ \mathrm{nm}$, it lies in the \textbf{visible} region -- it is the famous red H-$\alpha$ line of hydrogen.

\textbf{4. Ionisation from $n = 2$.} Ionisation means taking the electron from $E_2$ to $E_\infty = 0$:
\[ E_{\text{ion}} = 0 - (-3.40) = 3.40\ \mathrm{eV} = 5.44\times10^{-19}\ \mathrm{J}. \]
Any photon with $hf \geq 3.40\ \mathrm{eV}$ can ionise the atom from this level; the excess energy appears as kinetic energy of the freed electron.
\end{exoresolu}

\courssub{Method 5: Explaining emission and absorption line spectra}

\begin{methode}[Interpreting a line spectrum]
\begin{enumerate}
\item Identify the type of spectrum: a \textbf{continuous spectrum} (hot solid, liquid or dense gas), an \textbf{emission line spectrum} (bright coloured lines on a dark background, from a low-pressure gas discharge), or an \textbf{absorption line spectrum} (dark lines crossing a continuous spectrum, produced when white light passes through a cool gas).
\item Connect lines to transitions: each line corresponds to a photon of a \emph{single} energy, hence to a single transition between two energy levels. The existence of discrete lines is direct evidence for \cle{quantised energy levels}.
\item For absorption: the cool gas absorbs photons of exactly the same energies it would emit, so absorption lines occur at the \emph{same wavelengths} as the emission lines of the same element.
\item Use the fingerprint idea: each element has a unique set of levels, hence a unique line spectrum -- this is how the composition of stars and of gas samples is identified.
\end{enumerate}
\end{methode}

\begin{exoresolu}[Absorption by hydrogen gas]
White light from an incandescent lamp passes through a tube of cool, low-pressure hydrogen gas, then through a diffraction grating. The hydrogen atoms are nearly all in the ground state ($E_1 = -13.6\ \mathrm{eV}$, $E_2 = -3.40\ \mathrm{eV}$, $E_3 = -1.51\ \mathrm{eV}$, $E_4 = -0.85\ \mathrm{eV}$).
\begin{enumerate}
\item Explain why dark lines appear in the continuous spectrum.
\item Calculate the wavelength of the dark line corresponding to the transition $n = 1 \to n = 2$.
\item Explain why no dark line is observed at the wavelength of the $n = 2 \to n = 3$ transition.
\end{enumerate}
\tcblower
\textbf{1. Origin of the dark lines.} The lamp emits a continuous spectrum. Photons whose energy matches exactly a level difference of hydrogen (e.g. $E_2 - E_1$) are absorbed by the gas atoms, which are excited to higher levels. These directions are therefore depleted at those precise wavelengths: dark lines appear on the continuous background. The excited atoms re-emit photons of the same energy, but in \emph{random directions}, so the forward beam remains deficient at those wavelengths.

\textbf{2. Wavelength of the $1 \to 2$ line.}
\[ \Delta E = E_2 - E_1 = (-3.40) - (-13.6) = 10.2\ \mathrm{eV} = 10.2\times1.60\times10^{-19} = 1.63\times10^{-18}\ \mathrm{J}. \]
\[ \lambda = \frac{hc}{\Delta E} = \frac{6.63\times10^{-34}\times3.00\times10^{8}}{1.63\times10^{-18}} = 1.22\times10^{-7}\ \mathrm{m} = 122\ \mathrm{nm}. \]
This is the Lyman-$\alpha$ line, in the ultraviolet.

\textbf{3. Why no $2 \to 3$ absorption line.} Absorption from the level $n = 2$ requires atoms already in that level. In cool gas, virtually all atoms are in the ground state $n = 1$ (the level $n = 2$ lies $10.2\ \mathrm{eV}$ above it and is not thermally populated). Since almost no atom is in $n = 2$, photons of energy $E_3 - E_2 = 1.89\ \mathrm{eV}$ find no atoms able to absorb them, and no corresponding dark line appears. Absorption lines from excited states are only seen in \emph{hot} gases, where collisions populate the excited levels.
\end{exoresolu}

\courssub{Method 6: Ionisation, excitation and counting spectral lines}

\begin{methode}[Excitation, ionisation and number of lines]
\begin{enumerate}
\item \textbf{Excitation} = raising the atom from a lower to a higher level: the absorbed energy must equal a level difference exactly (photon absorption) or be at least equal to it (collision with an electron, which keeps the surplus as kinetic energy).
\item \textbf{Ionisation energy} from level $n$: $E_{\text{ion}} = -E_n$ (for hydrogen, $13.6/n^2\ \mathrm{eV}$).
\item \textbf{Number of distinct lines} emitted when atoms excited to level $n$ de-excite by all possible cascades:
\[ N = \frac{n(n-1)}{2}. \]
\item For a photon to be absorbed it must match a transition \emph{exactly}; a photon with more energy than the ionisation threshold can also ionise the atom.
\end{enumerate}
\end{methode}

\begin{exoresolu}[Bombarding hydrogen with electrons]
A beam of electrons of kinetic energy $12.75\ \mathrm{eV}$ passes through hydrogen gas whose atoms are all in the ground state ($E_n = -13.6/n^2\ \mathrm{eV}$).
\begin{enumerate}
\item Determine the highest level to which the hydrogen atoms can be excited.
\item How many different wavelengths can the gas subsequently emit? List the transitions.
\item Can a \emph{photon} of energy $12.75\ \mathrm{eV}$ excite the same atoms? Justify.
\end{enumerate}
\tcblower
\textbf{1. Highest level reached.} Compute the excitation energies from the ground state:
\begin{align*}
E_2 - E_1 &= -3.40 - (-13.6) = 10.2\ \mathrm{eV},\\
E_3 - E_1 &= -1.51 - (-13.6) = 12.09\ \mathrm{eV},\\
E_4 - E_1 &= -0.85 - (-13.6) = 12.75\ \mathrm{eV}.
\end{align*}
An electron of $12.75\ \mathrm{eV}$ can give up exactly $12.75\ \mathrm{eV}$ in a collision (it need not keep any surplus), so atoms can be raised to $n = 4$. The highest level is $n = 4$.

\textbf{2. Number of emission lines.} From $n = 4$ the atoms de-excite by all possible routes:
\[ N = \frac{n(n-1)}{2} = \frac{4\times3}{2} = 6 \text{ distinct wavelengths}. \]
The transitions are: $4\to3$, $4\to2$, $4\to1$, $3\to2$, $3\to1$, $2\to1$. For example, $2\to1$ gives Lyman-$\alpha$ ($122\ \mathrm{nm}$, UV) and $3\to2$ gives the red Balmer line ($656\ \mathrm{nm}$).

\textbf{3. Photon of $12.75\ \mathrm{eV}$?} A photon can only be absorbed if its energy matches \emph{exactly} a level difference (photons cannot be partially absorbed). Here $12.75\ \mathrm{eV}$ equals exactly $E_4 - E_1$, so -- remarkably -- this photon \textbf{can} be absorbed, exciting the atom to $n = 4$. However, a photon of, say, $12.5\ \mathrm{eV}$ would pass through without effect, whereas an electron of $12.5\ \mathrm{eV}$ could still excite the atom to $n = 3$ and keep $0.41\ \mathrm{eV}$ of kinetic energy. This contrast between \emph{all-or-nothing} photon absorption and \emph{flexible} collisional excitation is a favourite GCE question.
\end{exoresolu}

\courssub{Method 7: Using the Schr\"odinger (quantum) model of the hydrogen atom}

\begin{methode}[Working with quantum numbers]
\begin{enumerate}
\item In the Schr\"odinger model the electron is described by a \cle{wave function} $\psi$; $|\psi|^2$ gives the \emph{probability density} of finding the electron at a point -- we speak of an \cle{orbital} (region of high probability), not a definite orbit.
\item Know the four quantum numbers and their allowed values:
\begin{itemize}
\item $n = 1, 2, 3, \dots$ (principal): fixes the energy, $E_n = -13.6/n^2\ \mathrm{eV}$, and roughly the size of the orbital;
\item $\ell = 0, 1, \dots, n-1$ (orbital): fixes the shape ($\ell = 0$: $s$, spherical; $\ell = 1$: $p$; $\ell = 2$: $d$);
\item $m_\ell = -\ell, \dots, 0, \dots, +\ell$ (magnetic): orientation of the orbital;
\item $m_s = \pm\tfrac{1}{2}$ (spin).
\end{itemize}
\item Apply the \cle{Pauli exclusion principle}: no two electrons in an atom share the same set of four quantum numbers $\Rightarrow$ each orbital holds at most 2 electrons.
\item To check whether a proposed set $(n, \ell, m_\ell, m_s)$ is allowed, verify each range in turn.
\end{enumerate}
\end{methode}

\begin{exoresolu}[Quantum numbers in the hydrogen atom]
\begin{enumerate}
\item For the level $n = 3$ of hydrogen, list all the allowed values of $\ell$ and, for each, the allowed values of $m_\ell$. How many distinct orbitals does the shell $n = 3$ contain, and what is the maximum number of electrons it can hold?
\item State, with reasons, whether each of the following sets of quantum numbers is allowed: (i) $(2, 1, -1, +\tfrac12)$; (ii) $(2, 2, 0, -\tfrac12)$; (iii) $(3, 1, 2, +\tfrac12)$.
\item Contrast the description of the electron in the ground state of hydrogen given by Bohr's model and by Schr\"odinger's model.
\end{enumerate}
\tcblower
\textbf{1. Shell $n = 3$.} The rule $\ell = 0, 1, \dots, n-1$ gives $\ell = 0, 1, 2$:
\begin{center}
\begin{tabular}{|c|c|c|c|}
\hline
\rowcolor{popblueL} $\ell$ & Sub-shell & Allowed $m_\ell$ & Number of orbitals \\
\hline
$0$ & $3s$ & $0$ & $1$ \\
\hline
$1$ & $3p$ & $-1, 0, +1$ & $3$ \\
\hline
$2$ & $3d$ & $-2, -1, 0, +1, +2$ & $5$ \\
\hline
\end{tabular}
\end{center}
Total: $1 + 3 + 5 = 9$ orbitals (in general $n^2$). By the Pauli principle each orbital holds at most 2 electrons (with opposite spins), so the shell $n = 3$ can hold at most $2n^2 = 18$ electrons. Hydrogen of course has only one electron, but the counting applies to any atom.

\textbf{2. Validity of the sets.}
\begin{itemize}
\item (i) $(2, 1, -1, +\tfrac12)$: $n = 2$ allows $\ell = 0, 1$; $\ell = 1$ allows $m_\ell = -1, 0, 1$; $m_s = +\tfrac12$ is fine. \textbf{Allowed} (a $2p$ electron).
\item (ii) $(2, 2, 0, -\tfrac12)$: for $n = 2$, the maximum value of $\ell$ is $n - 1 = 1$, so $\ell = 2$ is impossible. \textbf{Forbidden.}
\item (iii) $(3, 1, 2, +\tfrac12)$: $\ell = 1$ requires $m_\ell \in \{-1, 0, +1\}$, so $m_\ell = 2$ is impossible. \textbf{Forbidden.}
\end{itemize}

\textbf{3. Bohr versus Schr\"odinger.} In \textbf{Bohr's model} the ground-state electron follows a definite circular orbit of radius $r_1 = 5.3\times10^{-11}\ \mathrm{m}$ with a definite speed -- an orbit like a planet around the Sun. In \textbf{Schr\"odinger's model} the electron has no definite trajectory: it is described by a wave function, and $|\psi|^2$ tells us only the \emph{probability} of finding it at each point. For the ground state ($n = 1, \ell = 0$, the $1s$ orbital) this probability cloud is spherically symmetric, and its most probable radius happens to equal the Bohr radius $5.3\times10^{-11}\ \mathrm{m}$. Both models give the same energy $E_1 = -13.6\ \mathrm{eV}$, but only the Schr\"odinger model correctly accounts for the shapes of orbitals, the spectra of many-electron atoms, and the fine details of the hydrogen spectrum. The quantum model replaces certainty of position with probability -- a profound change in our picture of the atom.
\end{exoresolu}

\begin{savaistu}[Helium: discovered in the Sun before the Earth]
The line spectrum is a true chemical fingerprint. In 1868, astronomers studying the spectrum of the Sun's chromosphere during an eclipse found a yellow line ($\lambda \approx 588\ \mathrm{nm}$) that matched no known element on Earth. They concluded it belonged to a new element, named \emph{helium} after \emph{helios}, the Greek word for the Sun. Helium was only identified on Earth 27 years later, in 1895. The same reasoning is used today to determine the composition of stars, of interstellar gas clouds and of the atmospheres of exoplanets -- all from the positions of their spectral lines.
\end{savaistu}

\begin{attention}[Frequent errors at the GCE]
\begin{itemize}
\item Mixing units: always convert eV to joules before using $E = hf$, and nm to metres before using $c = f\lambda$.
\item Sign errors: level energies are \emph{negative}; the emitted photon energy is $E_{n_i} - E_{n_f}$ with $n_i > n_f$ -- a positive quantity.
\item Believing an absorbed photon can have \emph{any} energy greater than a level gap: photon absorption is all-or-nothing and requires an exact match (unlike electron collisions).
\item Confusing the Rutherford conclusions: large-angle scattering is rare \emph{because} the nucleus is tiny, not because the alpha particles are slow.
\item Forgetting that $E = hf = E_2 - E_1$ applies to \emph{one} photon emitted in \emph{one} transition between two levels.
\end{itemize}
\end{attention}

\newpage

\coursec{Exercises}

\courssub{I check my knowledge}

\begin{exercice}[MCQ: atomic models and spectra]
For each question, choose the single correct answer.
\begin{enumerate}
\item In Rutherford's alpha scattering experiment, the fact that most alpha particles passed through the gold foil with little or no deflection shows that:
\begin{enumerate}
\item the atom is mostly empty space;
\item the nucleus is positively charged;
\item electrons orbit the nucleus;
\item the atom is indivisible.
\end{enumerate}
\item The very small number of alpha particles deflected through angles greater than $90^{\circ}$ is explained by:
\begin{enumerate}
\item collisions with electrons;
\item the repulsion by a small, dense, positively charged nucleus;
\item diffraction by the atoms of the foil;
\item the plum-pudding distribution of charge.
\end{enumerate}
\item The electron volt is a unit of:
\begin{enumerate}
\item potential difference;
\item charge;
\item energy;
\item power.
\end{enumerate}
\item A photon is emitted by a hydrogen atom when the electron:
\begin{enumerate}
\item jumps from a lower to a higher energy level;
\item falls from a higher to a lower energy level;
\item remains in its ground state;
\item leaves the nucleus.
\end{enumerate}
\end{enumerate}
\end{exercice}

\begin{exercice}[True or false — justify]
State whether each statement is true or false, and justify your answer in one or two sentences.
\begin{enumerate}
\item In Thomson's plum-pudding model, the positive charge is concentrated in a tiny central nucleus.
\item The emission spectrum of a gas consists of a continuous band of colours.
\item The absorption lines of a gas occur at exactly the same wavelengths as its emission lines.
\item In the Schrodinger model, the electron moves round the nucleus on a well-defined circular orbit.
\end{enumerate}
\end{exercice}

\begin{exercice}[Key definitions]
Define each of the following terms clearly, as you would in the GCE examination:
\begin{enumerate}
\item the \cle{electron volt} (eV);
\item an \cle{emission line spectrum};
\item the \cle{ground state} of an atom;
\item \cle{ionisation energy};
\item an \cle{orbital} in the Schrodinger model.
\end{enumerate}
\end{exercice}

\begin{exercice}[Direct calculations]
\begin{enumerate}
\item Convert $\SI{13.6}{eV}$ into joules. (Take $e = \SI{1.60e-19}{C}$.)
\item Calculate the energy of a photon of frequency $\SI{6.0e14}{Hz}$. Give your answer in joules and in electron volts. (Take $h = \SI{6.63e-34}{J.s}$.)
\item An electron in a hydrogen atom falls from the level $E_2 = \SI{-3.4}{eV}$ to the level $E_1 = \SI{-13.6}{eV}$. Calculate the energy of the emitted photon in joules.
\end{enumerate}
\end{exercice}

\courssub{I practise}

\begin{exercice}[Rutherford scattering: distance of closest approach]
An alpha particle of charge $+2e$ and kinetic energy $\SI{5.0}{MeV}$ is fired head-on towards a gold nucleus ($Z = 79$), assumed stationary.
\begin{enumerate}
\item Write down the expression for the electric potential energy of the alpha particle at a distance $r$ from the nucleus.
\item By applying conservation of energy, show that the distance of closest approach is $r_{\min} = \dfrac{2Ze^{2}}{4\pi\varepsilon_{0}E_{k}}$.
\item Calculate $r_{\min}$. (Take $\dfrac{1}{4\pi\varepsilon_{0}} = \SI{9.0e9}{N.m^{2}.C^{-2}}$.)
\item Compare your answer with the radius of a gold atom (about $\SI{1.3e-10}{m}$) and comment on what this tells us about the size of the nucleus.
\end{enumerate}
\end{exercice}

\begin{exercice}[Photon emitted by hydrogen]
The ground state of the hydrogen atom has energy $E_1 = \SI{-13.6}{eV}$ and the first excited state has energy $E_2 = \SI{-3.4}{eV}$.
\begin{enumerate}
\item Calculate the energy of the photon emitted when the electron falls from $E_2$ to $E_1$, in eV and in joules.
\item Deduce the frequency of the emitted photon.
\item Calculate its wavelength in air.
\item State the region of the electromagnetic spectrum to which this radiation belongs.
\end{enumerate}
(Take $h = \SI{6.63e-34}{J.s}$, $c = \SI{3.00e8}{m.s^{-1}}$, $e = \SI{1.60e-19}{C}$.)
\end{exercice}

\begin{exercice}[Identifying a transition]
A photon of wavelength $\SI{434}{nm}$ is emitted by a hydrogen atom. The energy levels of hydrogen are given by $E_n = -\dfrac{13.6}{n^{2}}\ \mathrm{eV}$.
\begin{enumerate}
\item Calculate the energy of this photon in electron volts.
\item By computing the energies of the first few levels, identify the two levels $n$ between which the transition takes place.
\item To which series (Lyman, Balmer or Paschen) does this line belong? Justify your answer.
\end{enumerate}
\end{exercice}

\begin{exercice}[Ionisation of hydrogen]
\begin{enumerate}
\item Explain what is meant by the \cle{ionisation energy} of the hydrogen atom.
\item Show that the minimum frequency of radiation that can ionise a hydrogen atom from its ground state is about $\SI{3.3e15}{Hz}$.
\item Calculate the corresponding wavelength and state the region of the spectrum concerned.
\item Explain why radiation of lower frequency cannot ionise the atom, however intense the beam may be.
\end{enumerate}
\end{exercice}

\courssub{I prepare for the GCE}

\begin{exercice}[GCE-style structured question: the nuclear atom and hydrogen spectra]
\begin{enumerate}
\item Describe Rutherford's alpha scattering experiment: apparatus, observations and conclusions. Explain, with reference to the observations, why Thomson's plum-pudding model had to be rejected. \textbf{(6 marks)}
\item The energy levels of the hydrogen atom are given by $E_n = -\dfrac{13.6}{n^{2}}\ \mathrm{eV}$. Sketch the energy level diagram for $n = 1$ to $n = 4$ and $n = \infty$, labelling the energy of each level. \textbf{(3 marks)}
\item On your diagram, mark with arrows the transitions giving rise to the Lyman series and the Balmer series, and state the spectral region of each series. \textbf{(3 marks)}
\item Calculate the longest wavelength of the Balmer series. \textbf{(4 marks)}
\item State one success and one limitation of the Bohr model of the hydrogen atom compared with the Schrodinger model. \textbf{(2 marks)}
\end{enumerate}
\textbf{(Total: 18 marks)}
\end{exercice}

\begin{exercice}[GCE-style structured question: emission and absorption spectra]
\begin{enumerate}
\item Describe, with the aid of a labelled diagram, how the emission spectrum of hydrogen gas at low pressure can be observed in the school laboratory. \textbf{(4 marks)}
\item Explain, in terms of energy levels and the equation $E = hf = E_2 - E_1$, why the spectrum consists of discrete lines rather than a continuous band. \textbf{(4 marks)}
\item White light is passed through a tube of cool hydrogen gas and the transmitted light is analysed. Explain why only certain wavelengths are absorbed, and why the absorption lines coincide exactly with the emission lines of hydrogen. \textbf{(4 marks)}
\item A line in the hydrogen spectrum has wavelength $\SI{656}{nm}$. Calculate the energy difference between the two levels involved, in eV, and identify the transition using $E_n = -\dfrac{13.6}{n^{2}}\ \mathrm{eV}$. \textbf{(4 marks)}
\end{enumerate}
\textbf{(Total: 16 marks)}
\end{exercice}

\begin{exercice}[Essay: from Bohr to Schrodinger]
Write a structured essay (about one page) describing how the Schrodinger model changed the picture of the electron in the hydrogen atom. Your essay should:
\begin{enumerate}
\item recall briefly the Bohr picture of quantised circular orbits and its successes; \textbf{(3 marks)}
\item explain the meaning of the wave function, the probability density and the concept of an orbital; \textbf{(5 marks)}
\item introduce the quantum numbers ($n$, $\ell$, $m_{\ell}$) and state what each one describes; \textbf{(4 marks)}
\item explain why the idea of a definite electron trajectory must be abandoned, and give one advantage of the Schrodinger model over the Bohr model. \textbf{(3 marks)}
\end{enumerate}
\textbf{(Total: 15 marks)}
\end{exercice}

\newpage

\coursec{Solutions to the exercises}

\begin{corrige}[Exercise 1 — MCQ: atomic models and spectra]
\begin{enumerate}
\item \textbf{Answer (a): The atom is mostly empty space.} If the atom were a uniform ball of matter (as in Thomson's model), every alpha particle would interact strongly and be deflected. The fact that the overwhelming majority cross the foil undeflected means they pass far from any massive charged region: the atom is essentially empty.
\item \textbf{Answer (b): The repulsion by a small, dense, positively charged nucleus.} Electrons are far too light to turn back a heavy alpha particle; a large-angle deflection requires a strong Coulomb repulsion at very short range, hence a concentrated positive charge. The plum-pudding model, with charge spread out, cannot produce such strong fields.
\item \textbf{Answer (c): Energy.} $1\ \mathrm{eV} = e \times 1\ \mathrm{V} = \SI{1.60e-19}{J}$ is the kinetic energy gained by an electron accelerated through $\SI{1}{V}$.
\item \textbf{Answer (b): Falls from a higher to a lower energy level.} The energy lost by the electron is carried away by the emitted photon: $hf = E_2 - E_1$ with $E_2 > E_1$. A jump from lower to higher level requires \emph{absorption} of a photon.
\end{enumerate}
\end{corrige}

\begin{corrige}[Exercise 2 — True or false — justify]
\begin{enumerate}
\item \textbf{False.} In Thomson's plum-pudding model the positive charge is spread uniformly throughout the whole volume of the atom (the ``pudding''), with the electrons embedded in it. The concentration of positive charge in a tiny nucleus is precisely Rutherford's conclusion, which destroyed Thomson's model.
\item \textbf{False.} The emission spectrum of a low-pressure gas consists of discrete bright lines at well-defined wavelengths, separated by dark regions. A continuous band is produced by hot solids, liquids or dense gases (e.g.\ a filament lamp).
\item \textbf{True.} Absorption corresponds to the same upward transitions ($E_1 \to E_2$) as emission does downward ($E_2 \to E_1$), so the photon energies $hf = |E_2 - E_1|$ — and hence the wavelengths — are identical. This is why the dark Fraunhofer-like lines of a gas coincide exactly with its bright emission lines.
\item \textbf{False.} In the Schrodinger model the electron has no definite trajectory. It is described by a wave function whose squared modulus gives the probability density of finding the electron at a point; one speaks of an \emph{orbital} (a region of high probability), not of a well-defined circular orbit.
\end{enumerate}
\end{corrige}

\begin{corrige}[Exercise 3 — Key definitions]
\begin{enumerate}
\item \textbf{Electron volt:} The electron volt is the kinetic energy gained by a particle of charge $e$ (one electron) accelerated from rest through a potential difference of one volt. $1\ \mathrm{eV} = \SI{1.60e-19}{J}$.
\item \textbf{Emission line spectrum:} The set of discrete bright lines, each of a definite wavelength, emitted by atoms of a gas at low pressure when they are excited (by heating or an electric discharge) and their electrons fall back to lower energy levels.
\item \textbf{Ground state:} The state of lowest possible energy of an atom ($n = 1$ for hydrogen, $E_1 = \SI{-13.6}{eV}$); it is the most stable state, in which the atom is normally found unless it has absorbed energy.
\item \textbf{Ionisation energy:} The minimum energy that must be supplied to an atom in its ground state to remove an electron completely (to $n = \infty$, $E = 0$). For hydrogen it is $\SI{13.6}{eV}$.
\item \textbf{Orbital:} In the Schrodinger model, a region of space around the nucleus, described by a wave function $\psi$ and a set of quantum numbers, within which the probability $|\psi|^{2}$ of finding the electron is high. It replaces the classical idea of a definite orbit.
\end{enumerate}
\end{corrige}

\begin{corrige}[Exercise 4 — Direct calculations]
\begin{enumerate}
\item $E = 13.6 \times \SI{1.60e-19}{J} = \boxed{\SI{2.18e-18}{J}}$.
\item In joules: $E = hf = 6.63\times 10^{-34} \times 6.0\times 10^{14} = \boxed{\SI{3.98e-19}{J}} \approx \SI{4.0e-19}{J}$.

In electron volts: $E = \dfrac{3.98\times 10^{-19}}{1.60\times 10^{-19}} = \boxed{\SI{2.5}{eV}}$.
\item The photon carries the energy difference:
\[ hf = E_2 - E_1 = -3.4 - (-13.6) = 10.2\ \mathrm{eV}. \]
Converting: $E = 10.2 \times 1.60\times 10^{-19} = \boxed{\SI{1.63e-18}{J}}$.
\end{enumerate}
\end{corrige}

\begin{corrige}[Exercise 5 — Rutherford scattering: distance of closest approach]
\begin{enumerate}
\item The gold nucleus carries charge $Q = +Ze$ and the alpha particle $q = +2e$. The electric potential energy of the pair at separation $r$ is
\[ E_p = \frac{1}{4\pi\varepsilon_0}\,\frac{qQ}{r} = \frac{1}{4\pi\varepsilon_0}\,\frac{2Ze^{2}}{r}. \]
\item \textbf{Proof (examinable).} Far from the nucleus the alpha particle has kinetic energy $E_k$ and negligible potential energy ($E_p \to 0$ as $r \to \infty$). At the distance of closest approach $r_{\min}$, the particle is instantaneously at rest (head-on collision), so $E_k' = 0$ and all the energy is potential. Conservation of mechanical energy gives
\[ E_k = \frac{1}{4\pi\varepsilon_0}\,\frac{2Ze^{2}}{r_{\min}}
\quad\Longrightarrow\quad
r_{\min} = \frac{2Ze^{2}}{4\pi\varepsilon_0\,E_k}. \]
\item Numerical application. Convert the kinetic energy:
\[ E_k = 5.0\times 10^{6}\ \mathrm{eV} \times 1.60\times 10^{-19}\ \mathrm{J/eV} = \SI{8.0e-13}{J}. \]
Then
\[ r_{\min} = \frac{9.0\times 10^{9} \times 2 \times 79 \times (1.60\times 10^{-19})^{2}}{8.0\times 10^{-13}}
= \frac{9.0\times 10^{9} \times 158 \times 2.56\times 10^{-38}}{8.0\times 10^{-13}}. \]
Numerator: $9.0\times 10^{9} \times 158 \times 2.56\times 10^{-38} = \SI{3.64e-26}{J\,m}$. Hence
\[ r_{\min} = \frac{3.64\times 10^{-26}}{8.0\times 10^{-13}} = \boxed{\SI{4.6e-14}{m}}. \]
\item Comparison: $\dfrac{r_{\min}}{r_{\text{atom}}} = \dfrac{4.6\times 10^{-14}}{1.3\times 10^{-10}} \approx 3.5\times 10^{-4}$, i.e.\ the closest approach is about $3000$ times smaller than the atomic radius. Since even a head-on $\SI{5}{MeV}$ alpha particle cannot get closer than about $\SI{4.6e-14}{m}$, the positive charge and nearly all the mass must be concentrated in a nucleus of radius smaller than this — typically of order $10^{-14}$ to $10^{-15}\ \mathrm{m}$ — while the atom itself is of order $10^{-10}\ \mathrm{m}$. The atom is indeed mostly empty space.
\end{enumerate}
\end{corrige}

\begin{corrige}[Exercise 6 — Photon emitted by hydrogen]
\begin{enumerate}
\item Energy of the photon:
\[ E = E_2 - E_1 = -3.4 - (-13.6) = \boxed{10.2\ \mathrm{eV}}. \]
In joules: $E = 10.2 \times 1.60\times 10^{-19} = \boxed{\SI{1.63e-18}{J}}$.
\item Frequency: $f = \dfrac{E}{h} = \dfrac{1.63\times 10^{-18}}{6.63\times 10^{-34}} = \boxed{\SI{2.46e15}{Hz}}$.
\item Wavelength: $\lambda = \dfrac{c}{f} = \dfrac{3.00\times 10^{8}}{2.46\times 10^{15}} = \boxed{\SI{1.22e-7}{m}} = \SI{122}{nm}$.
\item $\lambda \approx \SI{122}{nm} < \SI{400}{nm}$: this radiation lies in the \textbf{ultraviolet} region. (It is in fact the first line of the Lyman series, Lyman-$\alpha$.)
\end{enumerate}
\end{corrige}

\begin{corrige}[Exercise 7 — Identifying a transition]
\begin{enumerate}
\item Photon energy in joules:
\[ E = \frac{hc}{\lambda} = \frac{6.63\times 10^{-34} \times 3.00\times 10^{8}}{434\times 10^{-9}} = \SI{4.58e-19}{J}. \]
In electron volts:
\[ E = \frac{4.58\times 10^{-19}}{1.60\times 10^{-19}} = \boxed{\SI{2.86}{eV}}. \]
\item Compute the first levels with $E_n = -\dfrac{13.6}{n^{2}}\ \mathrm{eV}$:
\begin{center}
\begin{tabular}{|c|ccccc|}
\hline
\rowcolor{popblueL} $n$ & 1 & 2 & 3 & 4 & 5 \\
\hline
$E_n$ (eV) & $-13.6$ & $-3.40$ & $-1.51$ & $-0.85$ & $-0.54$ \\
\hline
\end{tabular}
\end{center}
Look for a difference equal to $\SI{2.86}{eV}$:
\[ E_5 - E_2 = -0.54 - (-3.40) = 2.86\ \mathrm{eV}. \]
The transition is therefore $\boxed{n = 5 \rightarrow n = 2}$.
\item The lower level of the transition is $n = 2$: the line belongs to the \textbf{Balmer series} (all transitions down to $n = 2$, lying in the visible/near-UV). Lyman lines end on $n = 1$ (ultraviolet) and Paschen lines on $n = 3$ (infrared).
\end{enumerate}
\end{corrige}

\begin{corrige}[Exercise 8 — Ionisation of hydrogen]
\begin{enumerate}
\item The \cle{ionisation energy} of hydrogen is the minimum energy required to remove the electron completely from the atom when it is initially in its ground state, i.e.\ to take it from $E_1 = \SI{-13.6}{eV}$ to $E_{\infty} = 0$. It equals $\SI{13.6}{eV} = \SI{2.18e-18}{J}$.
\item The photon must supply at least this energy: $hf_{\min} = |E_1|$, so
\[ f_{\min} = \frac{13.6 \times 1.60\times 10^{-19}}{6.63\times 10^{-34}}
= \frac{2.18\times 10^{-18}}{6.63\times 10^{-34}}
= \boxed{\SI{3.3e15}{Hz}}. \]
\item Corresponding wavelength:
\[ \lambda = \frac{c}{f_{\min}} = \frac{3.00\times 10^{8}}{3.3\times 10^{15}} = \boxed{\SI{9.1e-8}{m}} = \SI{91}{nm}, \]
which lies in the \textbf{ultraviolet} region.
\item A photon of frequency $f < f_{\min}$ carries an energy $hf < \SI{13.6}{eV}$, which is insufficient to take the electron from the ground state to $E = 0$. Since the electron can only absorb one photon at a time and the intermediate levels are quantised (no level exists at exactly the absorbed energy for most frequencies), the photon simply is not absorbed. Increasing the intensity only increases the \emph{number} of photons, not the energy of each one; no single photon has enough energy, so no ionisation occurs. This is exactly the same one-photon-one-electron logic as in the photoelectric effect.
\end{enumerate}
\end{corrige}

\begin{corrige}[Exercise 9 — GCE-style structured question: the nuclear atom and hydrogen spectra]
\textbf{(a) Experiment, observations, conclusions (6 marks).}
\emph{Apparatus (1 mark):} A radioactive source emitting a narrow collimated beam of alpha particles is directed at a very thin gold foil inside an evacuated chamber; a detector (zinc sulphide screen viewed with a microscope, or a counter) can be rotated around the foil to count scattered particles at each angle.

\emph{Observations (2 marks):} (i) The great majority of alpha particles pass through the foil with little or no deflection; (ii) A small fraction are deflected through appreciable angles; (iii) A very small fraction (about 1 in several thousand) are deflected through more than $90^{\circ}$, some almost straight back.

\emph{Conclusions (2 marks):} (i) Most particles are undeflected $\Rightarrow$ the atom is mostly empty space; (ii) Large deflections of heavy, fast alpha particles require a very strong repulsive force $\Rightarrow$ all the positive charge and nearly all the mass are concentrated in a tiny central \cle{nucleus}; (iii) The rarity of large deflections $\Rightarrow$ the nucleus is extremely small compared with the atom ($\sim 10^{-14}\ \mathrm{m}$ versus $\sim 10^{-10}\ \mathrm{m}$).

\emph{Rejection of Thomson's model (1 mark):} In the plum-pudding model the positive charge is spread over the whole atomic volume, so the electric field anywhere is far too weak to reverse the motion of an alpha particle; large-angle scattering would be impossible. Since such scattering is observed, the model must be rejected.

\textbf{(b) Energy level diagram (3 marks).} $E_n = -\dfrac{13.6}{n^{2}}\ \mathrm{eV}$ gives
\[ E_1 = -13.6,\quad E_2 = -3.40,\quad E_3 = -1.51,\quad E_4 = -0.85\ \mathrm{eV},\quad E_{\infty} = 0. \]
The diagram must show horizontal lines getting closer together as $n$ increases, each labelled with $n$ and its energy in eV, with $E = 0$ at the top (ionisation limit). \emph{(1 mark for correct energies, 1 mark for converging spacing, 1 mark for clear labels/scale.)}

\textbf{(c) Series (3 marks).} Lyman series: arrows from $n = 2, 3, 4, \dots$ \emph{down to} $n = 1$ — ultraviolet. Balmer series: arrows from $n = 3, 4, 5, \dots$ \emph{down to} $n = 2$ — visible (and near UV). \emph{(1 mark for each correctly drawn series, 1 mark for both regions.)}

\textbf{(d) Longest Balmer wavelength (4 marks).} The longest wavelength corresponds to the \emph{smallest} energy difference, i.e.\ the transition $n = 3 \to n = 2$:
\[ \Delta E = E_3 - E_2 = -1.51 - (-3.40) = 1.89\ \mathrm{eV} = 1.89 \times 1.60\times 10^{-19} = \SI{3.02e-19}{J}. \]
Then
\[ \lambda_{\max} = \frac{hc}{\Delta E} = \frac{6.63\times 10^{-34} \times 3.00\times 10^{8}}{3.02\times 10^{-19}} = \boxed{\SI{6.6e-7}{m}} \approx \SI{656}{nm}, \]
the red Balmer-$\alpha$ line. \emph{(1 mark for identifying $3 \to 2$, 1 mark for $\Delta E$, 1 mark for the formula, 1 mark for the numerical value with unit.)}

\textbf{(e) Success and limitation of Bohr (2 marks).} \emph{Success (1 mark):} The Bohr model predicts quantised energy levels $E_n = -13.6/n^{2}\ \mathrm{eV}$ which reproduce the hydrogen line spectrum (Lyman, Balmer, Paschen series) quantitatively. \emph{Limitation (1 mark):} It treats the electron as a classical particle on a definite circular orbit, which contradicts quantum mechanics; it fails for atoms with more than one electron and cannot explain fine structure or relative line intensities, whereas the Schrodinger model describes the electron by a wave function and orbitals and applies in principle to all atoms.
\end{corrige}

\begin{corrige}[Exercise 10 — GCE-style structured question: emission and absorption spectra]
\textbf{(a) Observing the emission spectrum (4 marks).} Hydrogen gas at low pressure is sealed in a discharge tube connected to a high-voltage supply; the gas glows. The light is passed through a narrow slit and then through a prism (or a diffraction grating), and the spectrum is observed on a screen or with a telescope. The labelled diagram should show: discharge tube + high voltage $\to$ slit $\to$ prism/grating $\to$ screen showing a few bright coloured lines on a dark background. \emph{(1 mark apparatus, 1 mark excitation by discharge, 1 mark labelled diagram, 1 mark description of discrete lines.)}

\textbf{(b) Why discrete lines (4 marks).} In the atom the electron can only occupy certain allowed energy levels $E_1, E_2, \dots$ — energy is \cle{quantised}. When the electron falls from a level $E_2$ to a lower level $E_1$, a single photon is emitted with
\[ hf = E_2 - E_1. \]
Since only definite values of $E_2 - E_1$ exist, only definite frequencies $f$ (and wavelengths $\lambda = c/f$) can be emitted. Each allowed transition gives one bright line; hence the spectrum is a set of discrete lines and not a continuous band. \emph{(1 mark quantised levels, 1 mark equation, 1 mark link between definite energy gaps and definite frequencies, 1 mark conclusion.)}

\textbf{(c) Absorption spectrum (4 marks).} White light contains photons of all wavelengths. Passing through cool hydrogen, a photon is absorbed only if its energy $hf$ is \emph{exactly} equal to the difference between two levels, $hf = E_2 - E_1$, so that it can raise an electron from a lower to a higher level; photons of other energies pass through unaffected. Hence only certain wavelengths are removed, giving dark lines. Because absorption ($E_1 \to E_2$) and emission ($E_2 \to E_1$) involve the \emph{same} pairs of levels and therefore the same energy differences, the absorption lines occur at exactly the same wavelengths as the emission lines. \emph{(1 mark photon energies in white light, 1 mark resonance condition, 1 mark dark lines, 1 mark coincidence justified by identical level differences.)}

\textbf{(d) The $\SI{656}{nm}$ line (4 marks).}
\[ \Delta E = \frac{hc}{\lambda} = \frac{6.63\times 10^{-34} \times 3.00\times 10^{8}}{656\times 10^{-9}} = \SI{3.03e-19}{J} = \frac{3.03\times 10^{-19}}{1.60\times 10^{-19}} = \boxed{\SI{1.89}{eV}}. \]
Using $E_n = -13.6/n^{2}\ \mathrm{eV}$: $E_2 = -3.40\ \mathrm{eV}$, $E_3 = -1.51\ \mathrm{eV}$, and
\[ E_3 - E_2 = 1.89\ \mathrm{eV}, \]
so the line corresponds to the transition $\boxed{n = 3 \rightarrow n = 2}$ (Balmer-$\alpha$, red). \emph{(1 mark formula, 1 mark value in eV, 1 mark computing levels, 1 mark identifying $3 \to 2$.)}
\end{corrige}

\begin{corrige}[Exercise 11 — Essay: from Bohr to Schrodinger]
\emph{Indicative content and mark scheme.}

\textbf{(a) The Bohr picture (3 marks).} In Bohr's model (1913) the electron moves around the proton on circular orbits without radiating, but only on certain \emph{allowed} orbits for which the angular momentum is quantised, $m_e v r = n\hbar$. Each orbit has a definite energy $E_n = -13.6/n^{2}\ \mathrm{eV}$, and radiation is emitted or absorbed only in jumps between orbits, with $hf = E_2 - E_1$. \emph{Successes:} It explained the stability of the atom and reproduced quantitatively the wavelengths of the hydrogen spectrum (Lyman, Balmer, Paschen series). \emph{(1 mark quantised orbits, 1 mark energy formula/jumps, 1 mark success with spectra.)}

\textbf{(b) Wave function, probability density, orbital (5 marks).} In the Schrodinger model the electron is described by a \cle{wave function} $\psi(x,y,z)$, obtained by solving the Schrodinger equation with the Coulomb potential of the nucleus. $\psi$ itself has no direct physical meaning, but $|\psi|^{2}\,dV$ gives the \emph{probability} of finding the electron in a small volume $dV$: $|\psi|^{2}$ is the \cle{probability density}. Where $|\psi|^{2}$ is large the electron is likely to be found; where it is zero the electron is never found. A stationary solution characterised by a set of quantum numbers defines an \cle{orbital}: the region of space in which the electron has a high probability (say $90\%$) of being found — for example the spherical 1s orbital of the ground state, or the dumb-bell shaped p orbitals. \emph{(1 mark wave function, 2 marks probability density, 2 marks orbital with example.)}

\textbf{(c) Quantum numbers (4 marks).} Each orbital is labelled by three quantum numbers:
\begin{itemize}
\item $n = 1, 2, 3, \dots$: The \emph{principal} quantum number — it fixes essentially the energy of the level ($E_n = -13.6/n^{2}\ \mathrm{eV}$ in hydrogen) and the average size of the orbital;
\item $\ell = 0, 1, \dots, n-1$: The \emph{orbital} (azimuthal) quantum number — it gives the magnitude of the orbital angular momentum and the \emph{shape} of the orbital (s, p, d, \dots);
\item $m_{\ell} = -\ell, \dots, 0, \dots, +\ell$: The \emph{magnetic} quantum number — it describes the \emph{orientation} of the orbital in space relative to an applied magnetic field.
\end{itemize}
\emph{(1 mark each correct quantum number with meaning, 1 mark for allowed ranges.)}

\textbf{(d) Abandoning the trajectory; advantage (3 marks).} According to Heisenberg's uncertainty principle, the position and momentum of an electron cannot both be known precisely at the same instant; the very notion of a definite path followed by the electron therefore loses its meaning. Experiments (electron diffraction) confirm that the electron behaves as a wave. The Schrodinger model replaces the orbit by a probability cloud: one can only state \emph{where the electron is likely to be}, never exactly where it is. \emph{Advantage over Bohr:} The model applies in principle to all atoms and molecules (not only hydrogen), predicts the correct quantisation of angular momentum (including states of zero angular momentum, impossible for Bohr orbits), and explains chemical bonding and line intensities. \emph{(1 mark uncertainty/no trajectory, 1 mark probability-cloud picture, 1 mark valid advantage.)}
\end{corrige}

\newpage

\coursec{Summary sheet}

\begin{popfigure}
\begin{tikzpicture}[font=\small,
  root/.style={rectangle, rounded corners, draw=popblue, fill=popblue, text=white, font=\bfseries\small, align=center, inner sep=4pt},
  br/.style={rectangle, rounded corners, draw=popdark, fill=popblueL, align=center, font=\scriptsize, inner sep=3pt},
  leaf/.style={align=center, font=\tiny, text=popdark, inner sep=1pt}
]
\node[root] (atom) at (0,0) {The Atom};
\node[br] (models) at (-4.6,2.2) {Atomic models};
\node[br] (ruth) at (0,3.0) {Rutherford scattering};
\node[br] (spectra) at (4.6,2.2) {Spectra};
\node[br] (levels) at (-4.6,-2.2) {Energy levels};
\node[br] (photon) at (0,-3.0) {Photon exchange};
\node[br] (schrod) at (4.6,-2.2) {Schr\"odinger model};
\draw[popline, thick] (atom) -- (models);
\draw[popline, thick] (atom) -- (ruth);
\draw[popline, thick] (atom) -- (spectra);
\draw[popline, thick] (atom) -- (levels);
\draw[popline, thick] (atom) -- (photon);
\draw[popline, thick] (atom) -- (schrod);
\node[leaf, below=1pt of models, align=center]{Dalton, Thomson\\ Rutherford: nucleus\\ Bohr: quantised orbits};
\node[leaf, above=1pt of ruth, align=center]{$\alpha$ on gold foil\\ Most pass through, few bounce back\\ Nucleus: tiny, dense, $+$};
\node[leaf, below=1pt of spectra, align=center]{Emission: bright lines\\ Absorption: dark lines\\ Each element unique};
\node[leaf, above=1pt of levels, align=center]{Quantised $E_n$\\ Ground, excited, ionisation\\ H: $E_n=-13.6/n^2$ eV};
\node[leaf, below=1pt of photon, align=center]{$E=hf=E_2-E_1$\\ $c=f\lambda$\\ $1\ \mathrm{eV}=1.60\times10^{-19}$ J};
\node[leaf, above=1pt of schrod, align=center]{Electron cloud\\ Orbitals, probability\\ Quantum numbers};
\end{tikzpicture}
\legende{Mind map of the chapter ``The Atom''.}
\end{popfigure}

\begin{aretenir}[The Atom — Essentials]

\textbf{Key definitions}
\begin{itemize}
\item \cle{Atom}: smallest neutral particle of an element; a tiny dense \cle{nucleus} (protons $+$ neutrons) surrounded by electrons.
\item \cle{Emission spectrum}: bright coloured lines on a dark background, emitted by a hot low-pressure gas; \cle{absorption spectrum}: dark lines on a continuous background, when white light crosses a cool gas. The lines occur at the \emph{same} wavelengths for a given element — its optical ``fingerprint''.
\item \cle{Energy level}: one of the discrete allowed energies of an atom. The lowest is the \cle{ground state}; higher ones are \cle{excited states}; $E=0$ corresponds to \cle{ionisation}.
\item \cle{Electron volt}: $\SI{1}{eV} = 1.60\times 10^{-19}\ \mathrm{J}$, the energy gained by an electron accelerated through $\SI{1}{V}$.
\item \cle{Photon}: quantum of light of energy $E=hf$.
\end{itemize}

\textbf{Laws and formulas to know}
\begin{itemize}
\item Rutherford's conclusions: the atom is mostly \emph{empty space}; the nucleus is very small ($\sim 10^{-15}$ m against $\sim 10^{-10}$ m for the atom), positively charged, and contains nearly all the mass.
\item Photon energy: $E = hf = \dfrac{hc}{\lambda}$, with $h = 6.63\times10^{-34}\ \mathrm{J\,s}$, $c = 3.00\times10^{8}\ \mathrm{m\,s^{-1}}$.
\item Transition equation: $\boxed{E = hf = E_2 - E_1}$. \textbf{Emission}: electron falls $E_2 \to E_1$ ($E_2>E_1$), photon emitted. \textbf{Absorption}: photon absorbed only if $hf$ matches \emph{exactly} a gap $E_2-E_1$.
\item Hydrogen levels (Bohr): $E_n = -\dfrac{13.6}{n^2}\ \mathrm{eV}$, $n = 1, 2, 3, \dots$ Ionisation energy from the ground state: $\SI{13.6}{eV}$.
\item Useful shortcut: $hc \approx 1240\ \mathrm{eV\,nm}$, so $E(\mathrm{eV}) \approx \dfrac{1240}{\lambda(\mathrm{nm})}$.
\item Schr\"odinger model: the electron is described by a \emph{wave function}; only the \cle{probability} of presence (orbital / electron cloud) is defined — no definite orbit. States are labelled by quantum numbers ($n$, $\ell$, $m_\ell$, $m_s$); quantisation of energy follows naturally from the wave equation.
\end{itemize}

\textbf{Methods}
\begin{itemize}
\item \emph{Photon from a transition}: compute $\Delta E = E_{\text{high}} - E_{\text{low}}$ (in eV), convert to joules ($\times 1.60\times10^{-19}$), then $f = \Delta E/h$ and $\lambda = c/f$.
\item \emph{Can a photon be absorbed?} Compare $hf$ with all gaps from the initial level: absorption occurs only for an \emph{exact} match; if $hf \geq |E_{\text{initial}}|$, the atom is ionised and the electron carries away the surplus as kinetic energy.
\item \emph{Counting emission lines} from level $n$: the electron cascades down; the number of possible transitions is $\dfrac{n(n-1)}{2}$.
\item \emph{Interpreting a spectrum}: bright lines $\Rightarrow$ emission (hot gas); dark lines at the same $\lambda$ $\Rightarrow$ absorption (cool gas in front of a continuous source).
\end{itemize}

\textbf{Common mistakes}
\begin{itemize}
\item Mixing up $E_2 - E_1$: the photon energy is always \emph{positive}; take higher level minus lower level.
\item Forgetting to convert eV $\to$ J before using $f = E/h$.
\item Using $\lambda$ in nm or \AA\ without converting to metres in $c = f\lambda$.
\item Believing any photon can be absorbed: only exact energy gaps work (absorption is selective, unlike ionisation).
\item Confusing the \emph{sign} of levels: bound states have $E_n < 0$; a transition \emph{down} releases energy.
\item Saying the electron ``orbits like a planet'' in the Schr\"odinger model: it is a probability cloud, not a trajectory.
\item Thinking Rutherford's large deflections were frequent: they were \emph{rare} ($\sim 1$ in $8000$), which is precisely why the nucleus must be tiny.
\end{itemize}

\textbf{GCE tips}
\begin{itemize}
\item Always state the \emph{three} observations of Rutherford's experiment (most $\alpha$ pass straight; some deflect through small angles; a very few rebound) and link each to a conclusion.
\item In $E = hf = E_2 - E_1$ questions, draw the level diagram, mark the arrow (up = absorption, down = emission), and give $\lambda$ with its unit and 3 significant figures.
\item Quote constants exactly: $h = 6.63\times10^{-34}\ \mathrm{J\,s}$, $e = 1.60\times10^{-19}\ \mathrm{C}$, $c = 3.00\times10^{8}\ \mathrm{m\,s^{-1}}$.
\item For ``explain why the spectrum is made of discrete lines'', answer: because energy levels are \emph{quantised}, so only photons of definite energies $E_2-E_1$ are emitted or absorbed.
\end{itemize}
\end{aretenir}

\findecours

\end{document}
